\section{Integral and $\lambda$-ring controls on the bridges}\label{sec:root-findings}
The results in this section are deductions checked in this run. They are not claims of historical priority. We distinguish an identity of characters, a decomposition in a Grothendieck group, an actual direct-sum decomposition, and a dimension identity.

\begin{proposition}[G-01: The three meanings of the second operation]\label{G-01}
Let $E$ be a finite $G$-set. Write $b_2(E)$ for the $G$-set of two-element subsets and $s_2(E)$ for the $G$-set of two-element multisets. In the Burnside ring one has
\[
 [s_2(E)]=[b_2(E)]+[E],\qquad [E]^2=[E]+[\operatorname{Inj}(\{1,2\},E)].
\]
The identity $[E]^2=2[b_2(E)]+[E]$ does not hold in general. If $V$ is a representation over a field of characteristic zero, the correct special $\lambda$-ring identities are
\[
 [\operatorname{Sym}^2V]=[V]^2-[\wedge^2V],\quad
 [\operatorname{Sym}^2V]-[\wedge^2V]=\psi^2[V],\quad
 [V]^2=2[\wedge^2V]+\psi^2[V].
\]
Consequently the expression $2\lambda^2(E)+E=E^2$ in cote 133, scan 019, line 53 is valid after taking cardinalities, but is not a universal equivariant or special-$\lambda$-ring identity. The supplied study's assertion that the two equivariant descriptions agree needs this correction.
\end{proposition}
\begin{proof}
Partition a multiset of cardinality two into the diagonal multisets $\{e,e\}$ and the subsets $\{e,e'\}$ with $e\ne e'$. This partition is $G$-equivariant. Similarly partition ordered pairs into diagonal and off-diagonal pairs. For a counterexample take $G=C_2$ and its regular two-element set $E$. Then $E\times E$ is a disjoint union of two free $C_2$-orbits, whereas $b_2(E)$ is one fixed point. Thus the asserted equality would read $2[E]=[E]+2[\mathrm{pt}]$. Applying the fixed-point homomorphism for the whole group gives $0=2$, a contradiction.
For representations, first prove the universal identity on $\mathrm{GL}_r$ and then pull it back along the representation of $G$. Write the character of its defining representation on the diagonal torus as $x_1+\cdots+x_r$. The symmetric and exterior squares have characters $\sum_{i\leq j}x_ix_j$ and $\sum_{i<j}x_ix_j$. Their sum is $(\sum_i x_i)^2$, and their difference is $\sum_i x_i^2=\psi^2[V]$. Character injectivity in characteristic zero proves the formulas. In the same $C_2$ example the permutation representation is $1\oplus\epsilon$, and $\wedge^2V=\epsilon$, whereas linearization of $b_2(E)$ is $1$. This explicitly shows why the subset operation cannot be silently identified with exterior power. Taking dimensions sends $\psi^2[V]$ to $\dim V$ and recovers the numerical identity. None of this determines the ink on a scan which has not been supplied.
\end{proof}

\begin{proposition}[G-02: The terminal exterior-power variable is necessary]\label{G-02}
For $i\geq2$, a universal formula for $\lambda^i(xy)$ in special $\lambda$-rings cannot in general omit $\lambda^i(x)$ or $\lambda^i(y)$. Thus the terminal index in cote 133, scan 006, lines 25--26 must be $i$ if those lines state the universal product formula.
\end{proposition}
\begin{proof}
Set $y=1$, for which $\lambda^1(1)=1$ and $\lambda^j(1)=0$ for $j>1$. A formula using only $x,\lambda^2(x),\ldots,\lambda^{i-1}(x)$ would express $\lambda^i(x)$ in the subring generated by these earlier operations. Evaluate at the defining representation of $\mathrm{GL}_i$ and restrict to its diagonal torus. This would express the elementary symmetric polynomial $e_i$ in $\mathbb Z[e_1,\ldots,e_{i-1}]$, contrary to the algebraic independence of $e_1,\ldots,e_i$ (equivalently the fundamental theorem of symmetric polynomials). Interchange $x$ and $y$ for the other endpoint. The argument establishes the required mathematical reading, not a new visual transcription.
\end{proof}

\begin{proposition}[G-03: Finite $\gamma$-data do not determine multiplicities]\label{G-03}
For every integer $D\geq0$ there are two honest representations $A_D,B_D$ of the split group scheme $\mathbb G_{m,\mathbb Z}$, each of rank $2^D$, which have the same class modulo $F^{D+1}_\gamma$ but different weight multiplicities. Over $\mathbb Q$ their Chern characters agree in degrees $0,\ldots,D$.
\end{proposition}
\begin{proof}
Let $L$ denote the standard character and put $r=D+1$. Define
\[
 A_D=\bigoplus_{\substack{0\leq j\leq r\\j\text{ even}}}L^{\otimes j\,\oplus\binom rj},\qquad
 B_D=\bigoplus_{\substack{0\leq j\leq r\\j\text{ odd}}}L^{\otimes j\,\oplus\binom rj}.
\]
The two binomial parity sums are $2^{r-1}$. In the representation ring $\mathbb Z[L,L^{-1}]$ the difference of the classes is $(1-L)^r$, a product of $r$ rank-zero elements and hence an element of $F^r_\gamma$. The constant weight has multiplicity one in $A_D$ and zero in $B_D$, so the characters differ. With $c=c_1(L)$, the difference of the Chern characters is $(1-e^c)^r$, whose lowest possible degree is $r$. This proves the assertion. For a line bundle on a scheme, applying the associated-bundle map gives the identical statement in its $K_0$. This does not assert that the two resulting bundles are nonisomorphic on every scheme.
\end{proof}

\begin{proposition}[G-04: An integral coefficient model and its exact specialization gate]\label{G-04}
Let $W_{\mathbb Z}$ be a finite free module, let $\rho:\mathrm{GL}_{N,\mathbb Z}\to\mathrm{GL}(W_{\mathbb Z})$ be a homogeneous polynomial representation of positive degree $e$, and let $f\in W_{\mathbb Z}$. Define $X_{\mathbb Z}$ as the affine schematic closure of its orbit. Its coordinate algebra is torsion-free, hence flat over $\mathbb Z$. In degree $d$ its orbit-coefficient map is an explicit finite integer matrix
\[
 A_d:\operatorname{Sym}^d(W_{\mathbb Z}^*)\longrightarrow
 \mathbb Z[g_{11},\ldots,g_{NN}]_{\leq ed},\qquad P\longmapsto P(\rho(g)f).
\]
Its rank over $\mathbb Q$ is the degree-$d$ coordinate-ring dimension. If $\delta_d$ is any nonzero maximal-rank minor of $A_d$, then this dimension and the degree-$d$ kernel commute with base change over $\mathbb Z[1/\delta_d]$. At a prime dividing $\delta_d$, rank must be tested anew; the fiber of the integral closure need not be the schematic orbit closure of the reduced point.
\end{proposition}
\begin{proof}
The coordinate homomorphism of the orbit map has target $\mathbb Z[\mathrm{GL}_N]=\mathbb Z[g_{ij},\det(g)^{-1}]$, a torsion-free $\mathbb Z$-module. Its kernel $I$ defines the schematic closure and its image $\mathbb Z[W]/I$ is torsion-free. A module over the principal ideal domain $\mathbb Z$ is flat if it is torsion-free, proving the first assertion. Since the representation is polynomial, each degree-$d$ homogeneous coordinate polynomial maps to a polynomial of degree at most $ed$ in the entries of $g$. Equality on invertible matrices is equality in the polynomial ring because localization by the nonzero polynomial $\det(g)$ is injective. This proves the finite coefficient-matrix description.
Let $r=\operatorname{rank}_{\mathbb Q}A_d$. All $(r+1)$-minors vanish as integers, while an $r$-minor $\delta_d$ is nonzero. Over $\mathbb Z[1/\delta_d]$, use that invertible block for row and column elimination. The vanishing of larger minors makes the remaining Schur complement zero, so the matrix is equivalent to $\operatorname{diag}(I_r,0)$. Its image and kernel are direct summands there and commute with every base change. Without inversion, an inclusion of lattices need not remain injective modulo a prime: already multiplication by $p$ on $\mathbb Z$ has rank one over $\mathbb Q$ and rank zero modulo $p$. Thus no unsupported all-primes orbit-fiber conclusion follows. Homogeneity sends degree $d$ to degree $ed$, so the kernel is homogeneous and the stated grading is well defined. If $r=0$, take $\delta_d=1$.
\end{proof}

\begin{proposition}[G-05: Ternary saturation and the three-row exclusion]\label{G-05}
Work over $\mathbb C$, and let $n\geq2$, $N=n^2$, and $\mathrm{Det}_n\subseteq\operatorname{Sym}^n(\mathbb C^N)^*$ be the determinant orbit closure. Then $\mathrm{Det}_n$ contains every degree-$n$ form supported on a fixed three-dimensional subspace. Consequently, for every partition $\lambda\vdash nd$ with $\ell(\lambda)\leq3$,
\[
 \operatorname{mult}_{S_\lambda\mathbb C^N}\mathbb C[\mathrm{Det}_n]_d
 =\operatorname{mult}_{S_\lambda\mathbb C^N}\operatorname{Sym}^d(\operatorname{Sym}^n\mathbb C^N).
\]
No multiplicity obstruction with at most three rows can separate any other orbit closure of a degree-$n$ form in the same ambient space, including a padded permanent, from $\mathrm{Det}_n$.
\end{proposition}
\begin{proof}
We use the established linear determinantal representation theorem for smooth plane curves: a smooth plane degree-$n$ curve together with a line bundle $L$ of degree $g-1$ with $H^0(L)=0$ admits an $n\times n$ matrix of linear forms whose determinant is its defining form up to a nonzero scalar (Beauville, 2000, Proposition 3.1(a)). Such an $L$ exists over $\mathbb C$: for $g\geq1$, the effective locus in $\operatorname{Pic}^{g-1}$ is the image of the symmetric power $C^{(g-1)}$, of dimension at most $g-1$, whereas $\operatorname{Pic}^{g-1}$ has dimension $g$; for $g=0$, take the degree-$-1$ line bundle. Scaling one row absorbs the nonzero scalar. Hence every smooth ternary degree-$n$ form is obtained by a linear substitution in $\det_n$. Any endomorphism substitution is a limit of invertible substitutions because $\mathrm{GL}_N$ is Zariski dense in $\operatorname{End}(\mathbb C^N)$; continuity of the polynomial action therefore puts these forms in $\mathrm{Det}_n$. Smooth ternary forms form a dense open set, so its closedness puts every ternary form there.
Let $W=\operatorname{Sym}^d(\operatorname{Sym}^n\mathbb C^N)$, and choose the Borel convention whose highest-weight vector has nonnegative weight $\lambda=(\lambda_1,\lambda_2,\lambda_3,0,\ldots,0)$. In a monomial basis of $W$, zero weight in positions $4,\ldots,N$ forces every coefficient variable to involve only the first three basis vectors. Thus the entire highest-weight space $W_\lambda^{\mathrm{hw}}$ is contained in $\operatorname{Sym}^d(\operatorname{Sym}^n\mathbb C^3)$. A nonzero polynomial in these coefficient variables cannot vanish on all ternary forms. By the first paragraph, the degree-$d$ determinant ideal therefore meets $W_\lambda^{\mathrm{hw}}$ trivially. Exactness and complete reducibility in characteristic zero give the asserted multiplicity equality. Every other closed subvariety has a degree-$d$ coordinate ring which is a quotient of $W$, so its multiplicity cannot exceed this common ambient value. This proves the exclusion. The assertion is about complex characteristic-zero multiplicities, not an unqualified statement in positive characteristic.
\end{proof}

\begin{proposition}[G-06: The four-variable gate is nonvacuous]\label{G-06}
For $n\geq4$, the closure $\mathcal D_{n,4}$ of degree-$n$ quaternary forms which are determinants of $n\times n$ matrices of linear forms has dimension at most $2n^2+2$. Since
\[
 \binom{n+3}{3}-(2n^2+2)=\binom{n-1}{3}>0,
\]
this is a proper subvariety of the space of quaternary forms. The determinant ideal in $\operatorname{Sym}(\operatorname{Sym}^n\mathbb C^{n^2})$ therefore has a nonzero isotypic component with a partition of exactly four rows in some finite degree.
\end{proposition}
\begin{proof}
Consider the polynomial map
\[
 (A_1,A_2,A_3,A_4)\longmapsto\det(x_1A_1+x_2A_2+x_3A_3+x_4A_4)
\]
from an affine space of dimension $4n^2$. The group
$H=\{(P,Q)\in\mathrm{GL}_n^2:\det(P)=\det(Q)\}$ acts by $A_i\mapsto PA_iQ^{-1}$, preserving this map. It has dimension $2n^2-1$. Its generic stabilizer has dimension one. Indeed, for a tuple with $A_1=I$, a diagonal $A_2$ with pairwise distinct eigenvalues, and an $A_3$ with every off-diagonal entry nonzero, a stabilizer has $Q=P$; commuting with $A_2$ makes $P$ diagonal, and commuting with $A_3$ makes it scalar. The scalar pairs $(cI,cI)$ always stabilize every tuple. Upper semicontinuity of stabilizer dimension, or the rank calculation for the commutation equations on this exhibited open set, gives generic dimension one. Thus generic fibers contain orbits of dimension $2n^2-2$, and the dimension-of-fibers theorem yields the stated image bound. The displayed identity follows by expanding the binomial coefficient.
Let $\pi$ set all but the first four variables to zero. For every form in $\mathrm{Det}_n$, its image under $\pi$ lies in $\mathcal D_{n,4}$: approximate the form by invertible substitutions of $\det_n$, apply $\pi$, and use the closedness of $\mathcal D_{n,4}$. Choose a nonzero homogeneous equation $P$ of this proper cone. Then $P\circ\pi$ is a nonzero homogeneous equation of $\mathrm{Det}_n$, using only first-four coefficient variables. The ideal of $\mathcal D_{n,4}$ is $\mathrm{GL}_4$-stable; choose a highest-weight vector in a nonzero irreducible submodule of its degree-$d$ part. Embedded in the full coordinate ring, it is also a $\mathrm{GL}_{n^2}$ highest-weight vector: raising operators involving a variable beyond the fourth annihilate it. Its highest weight has at most four rows, and G-05 rules out at most three. Hence its length is exactly four. Existence of such an equation does not establish that it is nonzero on the padded permanent.
\end{proof}

\begin{proposition}[G-07: An exact four-row determinant-closure formula]\label{G-07}
Set $D=320112$. For every $d\geq0$ and partition $\lambda\vdash4d$ with $\ell(\lambda)\leq4$, padded by zeros to length four, the degree-$d$ determinant-closure multiplicity satisfies
\[
 m_{\mathrm{Det}_4}(\lambda,d)
 =a_\lambda(d[4])-a_{\lambda-(D,D,D,D)}((d-D)[4]).
\]
The second term is defined to be zero if $d<D$ or $\lambda_4<D$. In particular the four-row sector is ambient for $d<D$, and at $d=D$ the multiplicity of $\lambda=(D^4)$ is exactly one less than its ambient plethysm multiplicity.
\end{proposition}
\begin{proof}
Leal, Lozano Huerta, and Vite (2024), Theorem 2 and Table 2, identify the closure of linear determinantal quartic surfaces as an irreducible reduced hypersurface of degree $D$ in $\mathbb P(\operatorname{Sym}^4\mathbb C^4)^*$. Table 2 identifies this component $F_1$ by $a=(5,5,5,5)$, $b=(6,6,6,6)$, so every matrix entry has degree one. Its affine cone has a prime principal ideal $(F)$ in the polynomial ring $A=\operatorname{Sym}(\operatorname{Sym}^4\mathbb C^4)$, with $F$ homogeneous of degree $D$.
The line $\mathbb CF$ is stable under $\mathrm{GL}_4$. Every algebraic one-dimensional character of $\mathrm{GL}_4$ is a power of the determinant. Since scalar $zI$ acts on $A_D$ by $z^{4D}$, this character is $\det^D$. Thus multiplication by $F$ gives the exact sequence of $\mathrm{GL}_4$-modules
\[
 0\longrightarrow \det^D\otimes A_{d-D}\xrightarrow{\ \cdot F\ }A_d
 \longrightarrow(A/(F))_d\longrightarrow0.
\]
For $d<D$ its first term is zero.
Let $\pi$ restrict forms in sixteen variables to their first four variables. The image $\pi(\mathrm{Det}_4)$ is precisely the affine cone just described. One inclusion follows by continuity from the orbit of $\det_4$. For the converse, every four-variable linear determinant is an endomorphism substitution in $\det_4$, hence belongs to its orbit closure; the same is true of every limit of such forms, and such a form is fixed by $\pi$. It follows that the intersection of the full determinant ideal with the first-four coefficient polynomial ring is exactly $(F)$. Highest-weight spaces for partitions of at most four rows use only these coefficient variables, by the weight argument of G-05. Applying complete reducibility to the exact sequence and shifting a highest weight by the determinant character proves the formula. At $d=D$ the source is $\det^D\otimes\mathbb C$, so the deficit is one.
\end{proof}
\paragraph{Source correction for G-07.}
The attached Landsberg source, section \texttt{zhanglee}, theorem \texttt{zlthm}, records $640224$. The later result used here is the degree of the reduced prime divisor. Leal--Lozano Huerta--Vite, section 4.1 immediately before Theorem 2, explicitly gives $P_{20,2}=\tfrac12D_{23,14}$. This is the precise cycle-multiplicity distinction; using the unreduced Noether--Lefschetz coefficient as the polynomial degree doubles the answer. This correction concerns the cited literature, not a purported visual correction of cote 133.

\begin{proposition}[G-08: A uniform low-variable determinantal construction]\label{G-08}
Let $m\geq1$ and $k\geq2$, put $a=\lfloor k/2\rfloor$, $b=\lceil k/2\rceil$, and
\[
 B(m,k)=\binom{m+a}{a}+\binom{m+b}{b}-1.
\]
Over any commutative ring, every degree-$m$ homogeneous polynomial $q$ in $k$ variables admits a matrix of linear forms of size $B(m,k)$ with determinant
\[
 \ell^{B(m,k)-m}q,
\]
where $\ell$ may first be an independent variable and then be specialized to any linear form. The construction is polynomial in the coefficients of $q$ and defined over $\mathbb Z$.
Consequently, over $\mathbb C$, whenever $n\geq B(m,k)$, no multiplicity obstruction of length at most $k$ can separate the orbit closure of an arbitrarily padded degree-$m$ form $\ell^{n-m}f$ from $\mathrm{Det}_n$.
\end{proposition}
\begin{proof}
Split the variables into disjoint lists $x=(x_1,\ldots,x_a)$ and $y=(y_1,\ldots,y_b)$, and write
\[
 q=\sum_{|\alpha|+|\beta|=m}c_{\alpha\beta}x^\alpha y^\beta.
\]
Make a rooted directed tree with one vertex for every $x$-monomial of degree at most $m$. The parent of a nonconstant monomial is obtained by deleting one occurrence of its largest-index variable; orient the edge away from the constant root and label it by that variable. There is exactly one root-to-vertex path, with label product $x^\alpha$. Make the corresponding $y$-tree with every edge reversed, so the unique path from vertex $\beta$ to its constant root has product $y^\beta$. Connect vertex $\alpha$ in the first tree to vertex $\beta$ in the second by an edge of constant label $c_{\alpha\beta}$ exactly when $|\alpha|+|\beta|=m$. Every source-to-sink path uses exactly one connecting edge; its label product is $c_{\alpha\beta}x^\alpha y^\beta$. Therefore its path polynomial is exactly $q$. The graph is acyclic and has
\[
 s=\binom{m+a}{a}+\binom{m+b}{b}=B(m,k)+1
\]
vertices.
Choose a topological ordering with source first and sink last, and let $A$ be its strictly upper triangular weighted adjacency matrix. Then $\det(I-A)=1$ and $(I-A)^{-1}=I+A+\cdots+A^{s-1}$. The $(1,s)$ entry is the source-to-sink path polynomial $q$. The adjugate identity therefore expresses $q$, up to the fixed sign $(-1)^{s+1}$, as the determinant of the $(s-1)\times(s-1)$ minor of $I-A$ obtained by deleting row $s$ and column $1$. Change the sign of one row if needed. Its entries are affine linear in $x,y$.
Replace each constant entry $c$ of this minor by $c\ell$. Its determinant is the degree-$B(m,k)$ homogenization of $q$, hence is $\ell^{B(m,k)-m}q$. This identity is universal over $\mathbb Z[c_{\alpha\beta},x,y,\ell]$. Padding the matrix with $n-B(m,k)$ diagonal entries $\ell$ proves the corresponding size-$n$ statement.
For the multiplicity assertion, restrict any invertible substitution of $\ell^{n-m}f$ to $k$ coordinate variables. The result has the form $L^{n-m}q$, with $q$ a homogeneous degree-$m$ form in these variables. The construction puts this result in $\mathrm{Det}_n$. Continuity gives the same inclusion for the entire padded orbit closure. If $P$ is any highest-weight polynomial of length at most $k$ vanishing on $\mathrm{Det}_n$, it uses only coefficient variables supported on the first $k$ underlying variables (in the sense made precise in G-05), and therefore vanishes on the other orbit closure as well. Thus the highest-weight space of the determinant ideal embeds in that of the padded-form ideal. Subtracting their dimensions from the common ambient highest-weight multiplicity gives $m_{\mathrm{padded}}(\lambda,d)\leq m_{\mathrm{Det}_n}(\lambda,d)$ in every degree.
\end{proof}
\paragraph{Quantified implication.}
For $k=4$, the construction gives $B(m,4)=(m+1)(m+2)-1$. More generally $B(m,k)=O_k(m^{\lceil k/2\rceil})$. Thus for a fixed number of rows this sector is excluded once the padding exceeds the displayed bound. This is an explicit restricted lower-bound obstruction theorem, not a general separation result and not an efficient solution of arbitrary plethysm.

\begin{proposition}[G-09: A sharper low-row exclusion by shared exponent chains]\label{G-09}
Let $m\geq1$ and $k\geq3$. Define the following integer bounds. If $k=2r+1$ is odd, set
\[
 C(m,k)=2\binom{m+r}{r}-1.
\]
If $k=2r+2\geq4$ is even, for $1\leq h\leq m+1$ put $Q=\lfloor m/h\rfloor$, $s_0=\min(h-1,m)$ and
\[
\begin{aligned}
 C_h(m,k)&=\sum_{q=0}^{Q}\binom{m-qh+r}{r}+Q(h-1)
          +\sum_{s=0}^{s_0}\binom{m-s+r}{r}-1,\\
 C(m,k)&=\min_{1\leq h\leq m+1}C_h(m,k).
\end{aligned}
\]
Over a field, for every nonzero linear form $L$ and homogeneous degree-$m$ form $q$ in $k$ variables, $L^{C(m,k)-m}q$ has a linear determinantal representation of size $C(m,k)$. In coordinates with $L=x_0$ the construction is defined over $\mathbb Z$ and polynomial in the coefficients of $q$. For fixed $k\geq3$,
\[
 C(m,k)=O_k\bigl(m^{(k-1)/2}\bigr).
\]
Consequently, over $\mathbb C$, if $n\geq C(m,k)$, there is no multiplicity obstruction with at most $k$ rows separating any padded degree-$m$ form from $\mathrm{Det}_n$. This strengthens G-08 in the specified range.
\end{proposition}
\paragraph{Prior-result boundary.}
The asymptotic determinantal-representation order used here is established: Boralevi, van Doornmalen, Draisma, Hochstenbach, and Plestenjak (2017), Theorem 1.3, prove order $m^{v/2}$ for a fixed number $v\geq2$ of affine variables, with a matching generic lower bound. Their proof in section 4 uses connected monomial spaces and an even/odd binary-digit split. We do not claim that order as new. The proof below gives the displayed finite integer bound by a different shared-chain construction, records its coefficient-ring scope, and proves the resulting GCT multiplicity exclusion. These are the forms needed here that the cited theorem does not state.
\begin{proof}
Extend $L$ to a basis of linear forms and write it as $x_0$. Dehomogenize $q$ by setting $x_0=1$, obtaining an arbitrary polynomial $p$ of degree at most $m$ in $k-1$ variables. It suffices to construct an affine linear determinant of size $C(m,k)$ equal to $p$. Homogenizing each matrix entry by $x_0$ then gives $x_0^{C(m,k)-m}q$.

For odd $k=2r+1$, split the dehomogenized variables into lists $x$ and $y$ of $r$ variables each. Use the two monomial trees from G-08, each with $\binom{m+r}{r}$ vertices, and connect every pair $(\alpha,\beta)$ with $|\alpha|+|\beta|\leq m$ by its coefficient in $p$. Every source-to-sink path is unique for its two monomial indices, so its path polynomial is $p$. The cofactor construction in G-08 has size one less than the total number of vertices, giving the displayed bound.

For even $k=2r+2$, write the variables as $x_1,\ldots,x_r,y_1,\ldots,y_r,z$ and fix $h$. Decompose the exponent of $z$ uniquely as $e=qh+s$, $0\leq s<h$. On the left, make a directed chain with $Qh$ edges, all labelled $z$, beginning at the source. At its vertex of index $qh$ attach the outward $x$-monomial tree of degrees at most $m-qh$, identifying its constant root with that chain vertex. Every such tree vertex indexed by $(q,\alpha)$ has a unique source path, with product $z^{qh}x^\alpha$. Chain vertices whose indices are not multiples of $h$ are transit vertices and receive no connecting edges. The left graph has
\[
 (Qh+1)+\sum_{q=0}^Q\left[\binom{m-qh+r}{r}-1\right]
 =\sum_{q=0}^Q\binom{m-qh+r}{r}+Q(h-1)
\]
vertices.

On the right, make a chain with vertices $0,\ldots,s_0$, directed toward vertex $0$, the sink, and again label its edges by $z$. At vertex $s$ attach the inward $y$-monomial tree of degrees at most $m-s$, identifying the root with that vertex. Every tree vertex indexed by $(s,\beta)$ has a unique sink path with product $y^\beta z^s$. The right graph has $\sum_{s=0}^{s_0}\binom{m-s+r}{r}$ vertices. Connect left vertex $(q,\alpha)$ to right vertex $(s,\beta)$ with the constant coefficient of $x^\alpha y^\beta z^{qh+s}$ in $p$ when $|\alpha|+|\beta|+qh+s\leq m$.

All connecting edges go from left to right. Thus every source-to-sink path uses exactly one such edge, and the uniqueness of the two tree paths and the division $e=qh+s$ gives exactly the required monomial with its coefficient, once. The path polynomial is $p$. Taking the same cofactor of $I-A$ as in G-08 yields an affine linear determinant of size $C_h(m,k)$. This graph has a path capable of realizing $z^m$ and a connecting edge, so its number of vertices is at least $m+2$; hence $C_h(m,k)\geq m+1$. The odd bound is also greater than $m$. Homogenization therefore uses a nonnegative exponent. Minimizing over the finite choices of $h$ proves the stated representation bound.

For the asymptotic statement in the even case take $h=\lceil\sqrt m\rceil$. The number of binomial summands is at most $m/h+1+h\leq2\sqrt m+2$, each bounded by $\binom{m+r}{r}=O_r(m^r)$, while $Q(h-1)\leq m$. Thus $C(m,k)=O_r(m^{r+1/2})$. The odd case is immediate from its binomial formula. These exponents are $(k-1)/2$ in both cases.

For the GCT consequence restrict a substituted padded form to any fixed $k$ underlying variables. If its restricted padding $L$ is nonzero, apply the determinant construction and pad with $n-C(m,k)$ diagonal copies of $L$. If $L=0$, the form is zero because $n\geq C(m,k)>m$. Thus every restriction lies in $\mathrm{Det}_n$. Apply the highest-weight ideal argument from G-08, which uses exactly this containment, to obtain the multiplicity inequality in every degree.
\end{proof}
\paragraph{Integral scope and growth of the remaining gate.}
With $L=x_0$ the graph and cofactor identity are universal over the coefficient ring. Extending an arbitrary nonzero $L$ to coordinates is a field operation; over a general base one instead works on a chart where a coefficient of $L$ is invertible and chooses the resulting local basis. No global polynomial choice of that basis is asserted. If $n=m^c$ with fixed $c>0$, every fixed $k\geq3$ satisfying $(k-1)/2<c$ is excluded for sufficiently large $m$. This is a quantitative restriction on possible multiplicity obstructions, not a claim that the remaining higher-row obstruction exists.
\begin{proposition}[G-10: An explicit degree bound for the four-row dimension gaps]\label{G-10}
For $n\geq4$ put
\[
 r_n=2n^2+2,\qquad B_n=n^{r_n-1},\qquad
 d_n=r_n(B_n-1)+1.
\]
For every integer $d\geq d_n$, some partition $\lambda\vdash nd$
with exactly four rows satisfies
\[
 m_\lambda(\operatorname{per}_n)>m_\lambda(\det_n).
\]
More generally, the same explicit threshold $d_n$ applies to every
padded pair $(m,n)$ for which the dimension criterion
$\dim\mathcal P_{m,4}+3>r_n$ of R-14 has been established.
The degree bound is very large and does not specify which partition
has the gap. It makes the preceding ``sufficiently large degree''
assertions effective; it gives no larger range of padded matrix sizes.
\end{proposition}
\begin{proof}
Write $Y=\mathbb P(\mathcal D_{n,4})$, $s=\dim Y$, and
$\Delta=\deg Y$. The pencil coefficient map projectivizes to a
rational map
\[
 \phi:\mathbb P^{4n^2-1}\dashrightarrow Y
\]
given by homogeneous polynomials of degree $n$. Its image is dense.
We first check the degree estimate
\begin{equation}\label{G10:image-degree}
 \Delta\leq n^s\leq n^{r_n-1}=B_n.
\end{equation}
Choose $s$ general hyperplanes in the target. Their intersection
with $Y$ consists of $\Delta$ distinct points in a dense open subset
contained in the actual image of $\phi$. Such an open subset exists
because the dominant image is constructible. Each point therefore
has a nonempty fiber in the domain of definition $U$ of $\phi$.
The common zero locus $Z$ of the $s$ pulled-back degree-$n$
equations in $\mathbb P^{4n^2-1}$ meets $U$ in the disjoint union
of these fibers. Every irreducible component of $Z$ meeting $U$
maps to just one of the target points. At least $\Delta$ distinct
components of $Z$ are consequently required. Each has degree at
least one, while repeated projective B\'ezout gives that the sum
of degrees of all components of $Z$ is at most $n^s$.
This proves the first inequality. Components contained in the
base locus only contribute to that upper bound and cannot cover
any of the nonempty fibers in $U$.
The second inequality follows from G-06, which gives $s\leq r_n-1$.
This is the usual image-degree argument with its base locus
accounted for explicitly; compare Sombra (1997), section~1,
printed preprint p.~5, and Proposition~2.10, pp.~11--12,
whose displayed statement treats an everywhere regular map.

The all-degree Hilbert bound of Chardin (1989), the unnumbered
central theorem on printed p.~306, applies to the homogeneous
prime ideal of the reduced irreducible variety $Y$ over
$\mathbb C$. Thus for every $d\geq1$,
\begin{equation}\label{G10:det-Hilbert}
 \dim\mathbb C[\mathcal D_{n,4}]_d
 \leq\Delta\binom{d+s}{s}
 \leq B_n\binom{d+r_n-1}{r_n-1}.
\end{equation}
We use this cited bound rather than assuming that the determinant
coordinate ring is Cohen--Macaulay.

Let $X$ be either $\mathcal P_{n,4}$ or a padded image
$\mathcal Z_{m,n}$ satisfying the stated criterion. It is a
reduced irreducible standard graded affine cone and has dimension
at least $r_n+1$, by R-15 or R-14 respectively. Linear Noether
normalization over the infinite field $\mathbb C$ supplies
$\dim X$ algebraically independent linear coordinate functions.
Any $r_n+1$ of these generate a polynomial subalgebra of its
coordinate ring. Counting its degree-$d$ monomials gives
\begin{equation}\label{G10:per-Hilbert}
 \dim\mathbb C[X]_d\geq\binom{d+r_n}{r_n}.
\end{equation}
The ratio of this binomial to the one in
\eqref{G10:det-Hilbert} is $(d+r_n)/r_n$.
When $d\geq r_n(B_n-1)+1$, that ratio is strictly greater
than $B_n$. Hence the Hilbert function of $X$ is strictly larger
in every such degree.

Decomposing the two degree-$d$ coordinate modules into
$\mathrm{GL}_4$ irreducibles, their dimension difference is
\[
 \sum_{\lambda\vdash nd,\,\ell(\lambda)\leq4}
 \bigl(m_\lambda(X)-m_\lambda(\mathcal D_{n,4})\bigr)
 \dim S_\lambda\mathbb C^4>0.
\]
At least one multiplicity difference is positive. The
highest-weight restriction argument of R-14 identifies these
multiplicities with those of the original orbit closures.
G-05 excludes positive differences with at most three rows.
The remaining partition has exactly four rows, as asserted.
\end{proof}
\paragraph{Effectivity and its limit.}
At the displayed degree, the finite coefficient and isotypic-rank
algorithms of R-05--R-06 can enumerate the partitions and find one
with a positive difference. The proof therefore bounds a finite
search, but its size is not polynomial. In the quartic case the
explicit family at degree $320112+s$ in R-12 is substantially
stronger than this uniform bound.


\begin{proposition}[G-11: Explicit rectangular multiplicity gaps]\label{G-11}
For $n\geq4$ define the following positive integers:
\[
\begin{gathered}
 b_n=\binom{n+3}{3},\qquad \beta_n=6b_n n^{30},\qquad
 T_n=\binom{b_n+\beta_n}{b_n}-1,\\
 E_n=4T_n(\beta_n!),\qquad u_n=2n^2-13,\\
 Q_n=(nE_n)^{u_n-1},\qquad K_n=u_n(Q_n-1)+1.
\end{gathered}
\]
For every integer $k\geq K_n$, put $d=E_n k$ and
$\lambda=(nd/4,nd/4,nd/4,nd/4)$. Then
\begin{equation}\label{G11:rectangle}
 m_\lambda(\operatorname{per}_n)>m_\lambda(\det_n).
\end{equation}
More generally, let $n>m\geq4$. The same conclusion, with
$\operatorname{per}_n$ replaced by
$\ell^{n-m}\operatorname{per}_m$, holds if
\begin{equation}\label{G11:padded-hypotheses}
 \dim\mathcal P_{m,4}+3>2n^2+2
 \quad\hbox{and}\quad n<4m/3.
\end{equation}
In particular it holds for $n=m+1$ and every $m\geq8$, and
for $n=\lfloor\rho m\rfloor$, all sufficiently large $m$,
and any fixed $1<\rho<\sqrt5/2$.
The degree, its progression, and the rectangular partition are
explicit. No useful computational size or polynomial algorithm
is asserted for these bounds.
\end{proposition}
\paragraph{Background used.}
We use the Kempf--Ness closed-orbit criterion in Losev (2006),
Theorem 4, p.~1; the stable-locus and quotient statements of
Mukai (2003), Definition 5.12, Proposition 5.15, Theorem 5.16,
and Corollary 5.17, pp.~165--167; and Derksen (2001),
Theorem 1.1, Proposition 1.2, and Example 2.1, printed p.~956.
The latter gives the displayed bound $\beta_n$ for generators
of the invariant ring of quaternary degree-$n$ forms.
The argument below proves the particular stability witnesses,
Veronese bound, and GCT deduction that those sources do not state.
\begin{proof}
Write $G=\mathrm{SL}_4(\mathbb C)$ and let $V_n$ be the space
of degree-$n$ quaternary forms. First we exhibit stable points
in the cones whose dimensions we compare. The form
\[
 f_n=x_0^n+x_1^n+x_2^n+x_3^n
\]
belongs to both $\mathcal P_{n,4}$ and $\mathcal D_{n,4}$.
Indeed factor $x_0^n+x_1^n$ and $x_2^n+x_3^n$ into binary
linear factors, place them on the diagonal and directed cycle
of an $n\times n$ pencil, and use its two perfect matchings.
For the determinant, multiply one cycle entry by the sign
$(-1)^{n-1}$.

Equip $V_n$ with an $\mathrm{SU}_4$-invariant Hermitian inner
product and use unit monomial vectors. Distinct monomials are
orthogonal. In the normalized sum of the four pure powers,
the mean exponent of each variable is $n/4$. Its moment map
therefore vanishes on the diagonal trace-zero Lie algebra.
It also vanishes on the off-diagonal directions: an operator
$x_j\partial/\partial x_i$, $i\ne j$, cannot carry one of
these pure powers to another when $n\geq3$. Thus the moment
map is zero. The Kempf--Ness criterion makes the $G$-orbit
closed. Its infinitesimal stabilizer is zero, since the
monomials $x_i^{n-1}x_j$ and $x_i^n$ in the derivative of
$f_n$ are distinct for $n\geq3$ and force all matrix entries
to vanish. Its stabilizer is consequently finite. This is a
stable point in each cone.

For the padded case put $a=n-m>0$ and suppose $3a<m$.
Let $e_0$ be the unit monomial $x_0^n$ and $e_i$ the unit
monomial proportional to $x_0^a x_i^m$, $1\leq i\leq3$.
Set
\[
 \theta_0=\frac{m-3a}{4m},\qquad
 \theta_1=\theta_2=\theta_3=\frac{n}{4m},\qquad
 F=\sum_{i=0}^3\sqrt{\theta_i}\,e_i.
\]
All four numbers are positive and sum to one. The mean
exponent of $x_0$ is $n\theta_0+3a\theta_1=n/4$, and
that of $x_i$ is $m\theta_i=n/4$. No two supported monomials
differ by a root $e_i-e_j$, because their differences are
multiples of $m\geq4$. The same moment-map calculation
therefore gives zero. We have $F=x_0^a q$, where
$q=\sum_{i=0}^3 c_i x_i^m$ and all $c_i\ne0$.
The preceding binary-factor construction puts $q$ in
$\mathcal P_{m,4}$, so $F\in\mathcal Z_{m,n}$.

Here is a direct finite-stabilizer check. For a linear vector
field $D_A=\sum_i L_i\partial/\partial x_i$, the equation
$D_A(F)=0$, after division by $x_0^{a-1}$, is
\[
 aL_0q+x_0D_Aq=0.
\]
Reduction modulo $x_0$ shows $L_0=d_0x_0$, because $a\ne0$
and $q(0,x_1,x_2,x_3)\ne0$. It follows that
$D_Aq=-ad_0q$. The distinct mixed monomials
$x_i^{m-1}x_j$ force every off-diagonal entry of $A$ to
vanish. The $x_0^m$ coefficient gives $(m+a)d_0=0$,
and the remaining pure powers give every other diagonal
entry zero. Thus the Lie stabilizer is zero and $F$ is
stable by Kempf--Ness.

For every reduced irreducible cone $X$ just considered,
the stable locus is a nonempty open subset and its quotient
fibers are the $15$-dimensional $G$-orbits (Mukai, loc. cit.).
Hence the dimension-of-fibers theorem yields
\begin{equation}\label{G11:quotient-dimension}
 \dim\mathbb C[X]^G=\dim X-15.
\end{equation}
In particular, writing $A_D=\mathbb C[\mathcal D_{n,4}]^G$
and $A_X=\mathbb C[X]^G$, G-06, R-15, and the first
hypothesis in \eqref{G11:padded-hypotheses} imply
\[
 \dim A_D\leq u_n,\qquad \dim A_X\geq u_n+1.
\]
The invariant functor is exact in characteristic zero.
Thus both invariant rings are graded quotients of
$\mathbb C[V_n]^G$, including when $X$ is padded; see also
Mukai, Proposition 5.10 and its proof, pp.~164--165.
Derksen's bound, specialized to $q=4$ and symmetric degree
$n$, says that this ambient invariant ring is generated in
degrees at most $\beta_n=6b_n n^{30}$. The dual convention
for forms makes no difference: the contragredient
transforms the action by an automorphism of $\mathrm{SL}_4$.
We may take at most $T_n$ homogeneous generators, because
the vector space of all coefficient polynomials in $b_n$
variables of positive degree at most $\beta_n$ has dimension
$T_n$.

We check explicitly that the $E_n$th Veronese of each ring
is generated in its first degree. More generally, suppose
an algebra has $r\leq T_n$ homogeneous generators with
positive degrees $d_i\leq\beta_n$, and put
$B_0=\beta_n!$. In a monomial in these generators of total
degree $kE_n$, remove as many pure blocks of
$B_0/d_i$ copies of generator $i$ as possible. Every removed
block has degree $B_0$. The remaining exponents are less
than $B_0/d_i$, so the residual degree is less than
$rB_0\leq T_n B_0$. It is an integer multiple $hB_0$,
because the original and removed degrees are such multiples;
thus $0\leq h<T_n$. There are $4kT_n-h$ pure blocks.
Combine the residual monomial with $4T_n-h$ of these blocks
to form a monomial of degree $E_n=4T_nB_0$, then divide
the other $4(k-1)T_n$ blocks into groups of $4T_n$.
This factors the original monomial into $k$ monomials of
degree $E_n$. Applying this to spanning monomials proves
the claim, also after passing to a quotient.

Consequently $A_D^{(E_n)}$ and $A_X^{(E_n)}$ are standard
graded domains, with dimensions unchanged. Set
$s=\dim\operatorname{Proj}A_D^{(E_n)}\leq u_n-1$.
This projective dimension is positive: block padding embeds
$\mathcal D_{4,4}$, whose dimension is $34$ by G-07,
into $\mathcal D_{n,4}$. By \eqref{G11:quotient-dimension}
the latter invariant ring has dimension at least $19$,
so $s\geq18$.
Choose a basis of its degree-one piece represented by
homogeneous invariant polynomials of degree $E_n$ on $V_n$.
Their composition with the determinant-pencil map is a
collection of homogeneous polynomials of degree $nE_n$.
They give a dominant rational map from
$\mathbb P^{4n^2-1}$ to $\operatorname{Proj}A_D^{(E_n)}$.
Dominance follows from dominance of the pencil map and
surjectivity of the affine quotient (Mukai, Theorem 5.9);
the Veronese morphism has this projective image. G-10's
base-locus-aware image-degree argument gives
\[
 \deg\operatorname{Proj}A_D^{(E_n)}
 \leq(nE_n)^s\leq Q_n.
\]
Chardin's bound in G-10 therefore gives, for every $k\geq1$,
\[
 \dim(A_D)_{E_nk}
 \leq Q_n\binom{k+u_n-1}{u_n-1}.
\]
Linear Noether normalization of the standard graded ring
$A_X^{(E_n)}$, whose dimension is at least $u_n+1$, gives
\[
 \dim(A_X)_{E_nk}\geq\binom{k+u_n}{u_n}.
\]
The ratio of the binomials is $(k+u_n)/u_n>Q_n$ when
$k\geq K_n$. Thus the first invariant dimension is strictly
smaller than the second.

Finally, in coefficient degree $d$, every polynomial
$\mathrm{GL}_4$ weight has size $nd$. The only irreducible
polynomial representation with trivial restriction to
$\mathrm{SL}_4$ is a determinant power; it has shape
$(nd/4)^4$, when $4\mid nd$, and is one-dimensional.
Here $E_n$ is divisible by four. The invariant dimension
difference is therefore exactly the difference of the two
multiplicities for the displayed rectangular partition.
The highest-weight restriction argument of G-05 and R-14
identifies these with the full orbit-closure multiplicities.
This proves \eqref{G11:rectangle} and its padded version.
R-14 and R-16 supply the stated particular ranges; all
satisfy $n<4m/3$.
\end{proof}
\paragraph{Relation to the rectangular bridge.}
The Haar projector onto $\mathrm{SU}_4$ invariants isolates
precisely these rectangular highest weights. Equivalently,
the corresponding character integral is the dimension of
this invariant graded piece. The proof makes that positive
rank exceed the determinant rank; it does not replace an
orbit-closure multiplicity by an ambient Kronecker
coefficient. The quartic formula G-07/R-12 is much sharper
where it applies. These deductions make no literature-priority
claim and are independent of RH and either proposed GRH input.

\begin{proposition}[G-12: A padding exclusion for the rectangular sector]\label{G-12}
Let $m\geq1$, $n>m$, and $n\geq4m/3$. For every degree-$m$
form $f$, every coordinate-ring degree $d\geq1$ with
$4\mid nd$, and $\lambda=(nd/4)^4$, one has
\begin{equation}\label{G12:exclusion}
 m_\lambda(\ell^{n-m}f)\leq m_\lambda(\det_n).
\end{equation}
If $n>4m/3$, the multiplicity on the left is zero.
The equality case $n=4m/3$ is included in the inequality.
Thus the rectangular gaps of G-11 cannot extend to or beyond
this padding boundary, even though nonrectangular four-row
questions remain.
\end{proposition}
\begin{proof}
As in G-05, restrict substituted forms to four variables.
They have shape $L^a q$, where $a=n-m$ and $q$ has degree
$m$. If $L=0$, the form is zero. Otherwise a change of
coordinates takes $L$ to $x_0$; it suffices to use a
$\mathrm{GL}_4$ change, since the determinant cone and each
homogeneous invariant ideal are stable under this group.
Act on the variables by the one-parameter subgroup
\[
 \eta(t)=\operatorname{diag}(t^3,t^{-1},t^{-1},t^{-1})
 \quad\hbox{of }\mathrm{SL}_4.
\]
A monomial in $x_0^a q$ with $x_0$ exponent $a+j$ has
weight $3a-m+4j$. If $3a>m$, all these weights are positive,
so $\eta(t)x_0^a q\longrightarrow0$ as $t\longrightarrow0$.
Every homogeneous positive-degree invariant therefore
vanishes on every such padded form and on their orbit
closure. Exactness of invariants identifies its invariant
coordinate ring with $\mathbb C$, whose positive-degree pieces vanish,
which proves the zero assertion.

If $3a=m$, the same limit instead equals
\[
 F_0=x_0^a q(0,x_1,x_2,x_3).
\]
Every ternary degree-$m$ form belongs to $\mathcal D_{m,3}$
by G-05 (including the elementary degrees one and two).
Embedding it and adding $a$ diagonal copies of $x_0$
shows $F_0\in\mathcal D_{n,4}$; the argument also holds for
limits because the block-padding map is polynomial.
Let $P$ be any homogeneous $\mathrm{SL}_4$ invariant
vanishing on $\mathcal D_{n,4}$. Since $P$ is constant
along $\eta$, continuity gives
$P(x_0^a q)=P(F_0)=0$. Thus the determinant ideal intersected
with the ambient invariant ring is contained in the ideal
of the restricted padded orbit closure intersected with
that ring. Both invariant coordinate rings are quotients
of the same ambient invariant ring, by exactness of
invariants. This containment gives a surjection from the
determinant invariant coordinate ring to the padded one,
and hence the degreewise dimension inequality.
Invariant dimensions are precisely the rectangular
multiplicities, as in G-11; highest-weight restriction
transfers the inequality to the full orbit closures.
\end{proof}
\paragraph{Source and scope.}
The strict inequality is also the rectangular specialization
of the usual padding weight constraint, but the proof above
gives its one-parameter-subgroup form and handles the
boundary by the ternary determinant theorem. This is an
exclusion for this specified rectangular sector. It does
not compare all four-row multiplicities or arbitrary
Kronecker coefficients.

\begin{proposition}[G-13: An explicit finite bound for the first rectangular gap]\label{G-13}
For $n\geq4$ put
\[
 \beta=6\binom{n+3}{3}n^{30},\qquad u=2n^2-13,\qquad
 C=(n\beta^2)^u,
\]
and define
\begin{equation}\label{G13:first-bound}
 s_*=(u+1)(C-1)+1,\qquad d_* = \beta s_*.
\end{equation}
There is an integer $d$ with $1\leq d\leq d_*$ and $4\mid nd$
such that, for $\lambda=(nd/4)^4$,
\[
 m_\lambda(\operatorname{per}_n)>m_\lambda(\det_n).
\]
The same conclusion holds with $\operatorname{per}_n$ replaced by
$\ell^{n-m}\operatorname{per}_m$ whenever $n>m\geq4$ and
\[
 \dim\mathcal P_{m,4}+3>2n^2+2,
 \qquad n<4m/3.
\]
In particular it applies to $n=m+1$, $m\geq8$, and, by R-18,
to $n=\lfloor\rho m\rfloor$ for all sufficiently large $m$
and each fixed $1<\rho<\sqrt{3/2}$.

The bound satisfies $\log d_*=O(n^2\log n)$. This is an
existence bound within a finite interval. It does not specify
the individual successful degree or replace G-11's assertion
at every degree of its displayed progression.
\end{proposition}

\paragraph{Background and scope.}
We use the stable points, exact invariant quotients, and quotient
dimensions established in G-11, together with Derksen (2001),
Theorem~1.1 and Example~2.1, printed p.~956. We also use the
projective image-degree argument of G-10 and Chardin's
unnumbered main theorem (1989), printed p.~306. The argument
below compares a weighted degree filtration with the ordinary
degree filtration in an affine generator embedding; it needs
no factorial Veronese.

\begin{proof}
Let $G=\mathrm{SL}_4(\mathbb C)$ and let $V_n$ be the space
of quaternary forms of degree $n$. Write
\[
 A_D=\mathbb C[\mathcal D_{n,4}]^G,
 \qquad A_X=\mathbb C[X]^G,
\]
where $X$ is $\mathcal P_{n,4}$ or the padded cone in the
statement. The stable-point and dimension arguments of G-11
give
\begin{equation}\label{G13:dims}
 t:=\dim A_D\leq u,\qquad \dim A_X\geq u+1.
\end{equation}
These are positively graded domains, with degree-zero part
$\mathbb C$, and both are graded quotients of
$\mathbb C[V_n]^G$. Choose homogeneous generators
$g_1,\ldots,g_r$ of the latter ring with positive degrees
$e_i\leq\beta$. Their images generate both quotient rings.

First bound the cumulative graded dimension of $A_D$. The
images of the generators give a closed affine embedding
\[
 Y_D=\operatorname{Spec} A_D\lhook\joinrel\longrightarrow
 \mathbb A^r.
\]
Compose this embedding with the determinant-pencil map and
the affine quotient. The resulting map from
$\mathbb A^{4n^2}$ has coordinates
$g_i(\det L)$, homogeneous of degrees $ne_i\leq n\beta$
in the pencil parameters. Its image is dense in $Y_D$.
To homogenize it, adjoin a source coordinate $z$ and use
the degree-$n\beta$ coordinate polynomials
\[
 z^{n\beta},\qquad
 z^{n\beta-ne_i}g_i(\det L),\quad 1\leq i\leq r.
\]
They define a rational map from $\mathbb P^{4n^2}$ whose
image closure is the ordinary projective closure
$\overline{Y_D}\subseteq\mathbb P^r$. It has dimension
$t$. G-10's image-degree bound, allowing a base locus, gives
\begin{equation}\label{G13:affine-degree}
 \deg\overline{Y_D}\leq(n\beta)^t\leq(n\beta)^u.
\end{equation}

For $d\geq0$, let
\[
 H_D^{\leq}(d)=\sum_{j=0}^d\dim(A_D)_j.
\]
Every element of $(A_D)_j$ has a preimage homogeneous of
weighted degree $j$ in $\mathbb C[T_1,\ldots,T_r]$ with
$\deg T_i=e_i$. Since every $e_i\geq1$, each monomial in
that preimage has ordinary degree at most $j$. It follows
that the direct sum of the graded pieces through $d$ is
contained in the space of functions on $Y_D$ represented
by ordinary polynomials of degree at most $d$.
This is an inclusion of filtered vector spaces, not an
identification of the two gradings.

The latter ordinary filtered dimension equals the
degree-$d$ Hilbert function of $\overline{Y_D}$:
homogenization with its first coordinate, followed by
dehomogenization on the dense affine chart, identifies the
two quotients. The projective closure is integral over
$\mathbb C$. Also $t\geq19$ by the fixed block-padding and
stable-quotient observation in G-11, so the positive-dimension
hypothesis in Chardin's stated theorem holds. That theorem
and \eqref{G13:affine-degree} yield
\begin{equation}\label{G13:upper}
 H_D^{\leq}(d)
 \leq(n\beta)^t\binom{d+t}{t}
 \leq(n\beta)^u\binom{d+u}{u}.
\end{equation}

For the other ring, a maximal algebraically independent
subset of the images of $g_1,\ldots,g_r$ has size
$\dim A_X$. Indeed those images generate its fraction
field, and maximality makes every remaining generator
algebraic over the chosen subset. By \eqref{G13:dims},
choose $u+1$ algebraically independent images. Each has
degree at most $\beta$. Their monomials of ordinary total
degree at most $s$ are linearly independent and have
graded degree at most $\beta s$. Consequently
\begin{equation}\label{G13:lower}
 H_X^{\leq}(\beta s)
 :=\sum_{j=0}^{\beta s}\dim(A_X)_j
 \geq\binom{s+u+1}{u+1}.
\end{equation}

For every nonnegative integer $s$, the exact binomial ratio
satisfies
\begin{align*}
 \frac{\binom{s+u+1}{u+1}}{\binom{\beta s+u}{u}}
 &=\frac{s+u+1}{u+1}
     \prod_{i=1}^u\frac{s+i}{\beta s+i}\\
 &\geq\frac{s+u+1}{(u+1)\beta^u}.
\end{align*}
The last inequality follows term by term from
$\beta(s+i)\geq\beta s+i$. At $s=s_*$ one has
\[
 s_*+u+1=(u+1)C+1,
\]
so the final lower bound is strictly greater than
$C/\beta^u=(n\beta)^u$. Equations
\eqref{G13:upper}--\eqref{G13:lower} therefore give
\[
 H_X^{\leq}(d_*)>H_D^{\leq}(d_*).
\]
The degree-zero dimensions are both one. Hence for some
$1\leq d\leq d_*$ the individual graded pieces satisfy
$\dim(A_X)_d>\dim(A_D)_d$.

As in G-11, the only polynomial $\mathrm{GL}_4$
representation of weight size $nd$ that is trivial on
$\mathrm{SL}_4$ is the one-dimensional determinant power
of shape $(nd/4)^4$. If $4\nmid nd$, both invariant
pieces are zero. Thus the successful degree necessarily
satisfies $4\mid nd$, and the invariant dimension gap is
the desired rectangular multiplicity gap. The first-four-
variable highest-weight argument in G-05/R-14 transfers
it to the full orbit closures.

Finally $\beta=O(n^{33})$ and $u=O(n^2)$, so the explicit
formula \eqref{G13:first-bound} gives
$\log d_*=O(n^2\log n)$. The particular padded ranges
follow from R-14/R-16 and R-18, respectively; their
dimension gaps and $n<4m/3$ are the same hypotheses already
used to construct the stable padded point in G-11.
\end{proof}

\paragraph{Complexity lower-bound context.}
The multiplicity statements G-10--G-13 locate particular representation
sectors and degrees. They do not improve the established quadratic
lower bound $\overline{\operatorname{dc}}(\operatorname{per}_m)
\geq m^2/2$ of Landsberg, Manivel, and Ressayre (2013);
the inspected April 2010 author version states it as Theorem~1.0.1,
p.~1 (G-B16). No claim about the strongest currently proposed lower
bound is needed here. A noncontainment result by itself does not
specify the four-row or rectangular multiplicity differences proved
in the present restricted parameter ranges.

\begin{proposition}[G-14: First-jet dominance and a stronger padding boundary]\label{G-14}
Let $m\geq1$, and write a quaternary form as
\[
 q(x_0,x_1,x_2,x_3)=\sum_{j=0}^m x_0^j q_j(x_1,x_2,x_3).
\]
The first-jet map
\[
 j_1:\mathcal D_{m,4}\longrightarrow
   \operatorname{Sym}^m\mathbb C^3\oplus
   \operatorname{Sym}^{m-1}\mathbb C^3,
 \qquad q\longmapsto(q_0,q_1),
\]
has Zariski-dense image.

Consequently let $n>m\geq1$, put $a=n-m$, and let $f$ be
any degree-$m$ form in the common ambient GCT setting.
For every $d\geq1$ and every partition $\lambda\vdash nd$
with at most four rows and
\begin{equation}\label{G14:boundary}
 \lambda_1\leq ad+1,
\end{equation}
one has
\[
 m_\lambda(\ell^a f)\leq m_\lambda(\det_n).
\]
If $\lambda_1<ad$, the multiplicity on the left is zero.
Thus every positive multiplicity gap in this sector must
satisfy $\lambda_1\geq ad+2$.
\end{proposition}

\paragraph{Scope and background.}
The proof uses the coefficient grading and first-four-variable
highest-weight restriction of G-05, together with elementary
cyclic cofactors and the finite geometric-product identity
already proved in R-16. It requires characteristic zero.
The jet-dominance calculation is given in full; no higher-jet
dominance or Lefschetz assertion is assumed.

\begin{proof}
The case $m=1$ of the first assertion is immediate. Suppose
$m\geq2$, choose a primitive $m$th root of unity $\zeta$,
and let $A$ be the ternary cyclic pencil with diagonal
entries
\[
 \ell_i=x_1-\zeta^i x_2,\qquad 0\leq i<m,
\]
and entries $x_3$ in positions $(i,i+1)$, with indices
modulo $m$. Its determinant has the two perfect matchings
and equals
\[
 x_1^m-x_2^m+(-1)^{m-1}x_3^m.
\]
We claim that the cofactors of $A$ span every ternary form
of degree $m-1$.

For each $0\leq k\leq m-1$, the cofactors using exactly
$k$ cycle entries are, up to nonzero signs, $x_3^k$ times
the $m$ cyclic translates of a consecutive product of
$h=m-1-k$ of the binary forms $\ell_i$. This follows from
the unique matching after deleting the relevant row and
column. A consecutive product, after a cyclic shift,
has the form
\[
 F_h(x_1,x_2)=\prod_{i=0}^{h-1}(x_1-\zeta^i x_2).
\]
Every coefficient in it is nonzero. Indeed its coefficient
of $x_1^{h-r}x_2^r$ is
\[
 (-1)^r\zeta^{r(r-1)/2}
   \prod_{j=1}^r\frac{1-\zeta^{h-r+j}}{1-\zeta^j},
 \qquad 0\leq r\leq h,
\]
with the empty product interpreted as one. Every exponent
in a nonempty numerator or denominator lies between $1$
and $h<m$. Thus no factor is zero. The cyclic substitution
$x_2\mapsto\zeta x_2$ has distinct characters on the
monomials with $0\leq r\leq h$. Fourier projection of the
$m$ translates therefore spans every binary form of degree
$h$. Varying $k$ proves the claim, since
\[
 \operatorname{Sym}^{m-1}\mathbb C^3
 =\bigoplus_{k=0}^{m-1}x_3^k
        \operatorname{Sym}^{m-1-k}\langle x_1,x_2\rangle.
\]

Now consider the parameter map
\[
 (A',B)\longmapsto j_1\det(A'+x_0B),
\]
where $A'$ is an arbitrary ternary linear pencil and $B$
is a constant $m\times m$ matrix. At the point $(A,0)$,
the derivative in the $A'$ directions changes $q_0$ by
arbitrary linear-form multiples of the cofactors of $A$.
By the claim these span all of
$\operatorname{Sym}^m\mathbb C^3$. It does not change
$q_1$, since the coefficient of $x_0$ is linear in $B$.
The derivative in the $B$ directions changes $q_1$ by
arbitrary constant linear combinations of the cofactors,
which span $\operatorname{Sym}^{m-1}\mathbb C^3$, and
does not change $q_0$. The differential is consequently
surjective onto the direct sum. The characteristic-zero
Jacobian criterion for dominance proves the first assertion.

For the multiplicity consequence, work first in four
variables and use the positive coefficient weight convention
of G-05/G-11 (equivalently, dualize the coordinate-ring
representations consistently). In a polynomial irreducible
of shape $\lambda$, every weight has each coordinate at
most $\lambda_1$. For example, the Schur construction uses
$\lambda_1$ exterior-power columns, and a given basis vector
can occur at most once in each column. This proves the
required bound on the $x_0$ coordinate of every weight.

Let $R=\mathbb C[\operatorname{Sym}^n\mathbb C^4]$ be
the coefficient polynomial ring, and let $R_{d,\lambda}$
be its degree-$d$ $\lambda$-isotypic subspace. Take
\[
 P\in I(\mathcal D_{n,4})\cap R_{d,\lambda}.
\]
The weight bound says that every coefficient monomial
occurring in $P$ has total $x_0$ weight at most
$\lambda_1$. Evaluate $P$ at a form $x_0^a q$. A
coefficient coming from $q_j$ has $x_0$ weight $a+j$;
a product of $d$ coefficients therefore has weight
\[
 ad+j_1+\cdots+j_d.
\]
All coefficient variables of smaller individual $x_0$
weight vanish on $x_0^a q$. If $\lambda_1<ad$, nothing
survives and every polynomial in $R_{d,\lambda}$ vanishes
on every such padded form.

Under \eqref{G14:boundary}, the restriction
$P(x_0^a q)$ depends only on $q_0,q_1$, because every
surviving coefficient monomial has
$j_1+\cdots+j_d\leq1$. It is a polynomial function of
these first-jet data. For every $q\in\mathcal D_{m,4}$,
block padding gives $x_0^a q\in\mathcal D_{n,4}$,
including the closure points. Hence this jet polynomial
vanishes on the dense image of $j_1$ just proved, and
therefore vanishes for every pair $(q_0,q_1)$. It follows
that $P(x_0^a q)=0$ for every quaternary form $q$ of
degree $m$.

To allow any padding form $L$, use a $\mathrm{GL}_4$
change taking $L$ to $x_0$. The determinant ideal and
the isotypic subspace are both stable under that group,
so the preceding conclusion applies to the corresponding
translate of $P$. Thus $P$ vanishes on every $L^a q$,
and in particular on the restricted padded orbit closure
associated with $\ell^a f$.
We have proved the containment of the determinant-ideal
$\lambda$-isotypic part in the corresponding padded ideal
part. Passing to the characteristic-zero semisimple
degree-$d$ quotients gives the stated multiplicity
inequality. The first-four-variable highest-weight
restriction identifies these multiplicities with those
in the full ambient orbit closures, as in G-05/R-14.
The same argument gives the zero assertion when
$\lambda_1<ad$.
\end{proof}

\paragraph{A precisely stated higher-jet continuation.}
For $0\leq h\leq m$, suppose that the projection
\[
 j_h:\mathcal D_{m,4}\longrightarrow
   \bigoplus_{j=0}^h\operatorname{Sym}^{m-j}\mathbb C^3
\]
has dense image. The identical coefficient-weight argument
then proves the same multiplicity inequality whenever
$\lambda_1\leq ad+h$: every surviving coefficient
monomial depends only on those jets. This is a conditional
continuation with an explicit geometric hypothesis.
The first-jet case above is proved without that hypothesis.
No dominance assertion for $h=2$ or $h=3$ is made here.

\begin{proposition}[G-15: Uniform third-jet dominance]\label{G-15}
For every $m\geq1$, the map
\[
 j_{\leq3}:\mathcal D_{m,4}\longrightarrow
   \bigoplus_{j=0}^{\min\{3,m\}}
          \operatorname{Sym}^{m-j}\mathbb C^3,
 \qquad \sum_{j=0}^m t^j q_j\longmapsto(q_0,q_1,q_2,q_3),
\]
with nonexistent components omitted, has Zariski-dense image.
For $m\geq4$, order three is the largest order for which the
full jet map can be dominant.

Consequently, if $n>m\geq1$, $a=n-m$, $d\geq1$, and
$\lambda\vdash nd$ has at most four rows, then
\[
 \lambda_1\leq ad+3
 \quad\Longrightarrow\quad
 m_\lambda(\ell^a f)\leq m_\lambda(\det_n)
\]
for every degree-$m$ form $f$ in the common ambient setting.
Every positive gap in this sector therefore satisfies
$\lambda_1\geq ad+4$. In particular a rectangular gap
with $\lambda=(nd/4)^4$ requires
\[
 d(4m-3n)>12.
\]
Since $4\mid nd$ for this rectangular shape, the last
quantity is divisible by four; equivalently,
$d(4m-3n)\geq16$ is necessary.
\end{proposition}

\paragraph{Background and proof strategy.}
We use the standard $\mathfrak{sl}_2$ decompositions
\[
 \operatorname{End}(\operatorname{Sym}^{m-1}\mathbb C^2)
       =\bigoplus_{l=0}^{m-1}\operatorname{Sym}^{2l}\mathbb C^2,
 \qquad
 \operatorname{Sym}^{r}(\operatorname{Sym}^2\mathbb C^2)
       =\bigoplus_{b=0}^{\lfloor r/2\rfloor}
             \operatorname{Sym}^{2r-4b}\mathbb C^2.
\]
These are the Clebsch--Gordan and harmonic decompositions;
see Fulton and Harris (1991), section~11.2, Exercises~11.11
and 11.14, printed pp.~151 and 153. We specify their
covariants, the finite orthogonal-polynomial calculation,
and the required rank argument below. No generic strong
Lefschetz property or Gulliksen--Neg\aa rd theorem is
assumed in this proof.

\begin{proof}
The case $m=1$ is immediate. Fix $m\geq2$ and let
$\rho$ be the irreducible representation of
$\mathfrak{sl}_2$ on $V=\operatorname{Sym}^{m-1}\mathbb C^2$.
Use standard generators $H,E,F$, with
$[H,E]=2E$, $[H,F]=-2F$, and $[E,F]=H$.
Choose the invariant Hermitian metric for which
$\rho(H)$ is Hermitian and $\rho(E)^*=\rho(F)$.
Its $m$ distinct $H$-eigenvalues are
\[
 \lambda_j=m-1-2j,\qquad 0\leq j<m.
\]
For $z=(z_1,z_2,z_3)$ put
\[
 A(z)=z_1H+z_2(E+F)+\mathrm i z_3(E-F),\qquad
 q=z_1^2+z_2^2+z_3^2,
\]
and consider the explicit linear pencil
\begin{equation}\label{G15:pencil}
 M(t,z)=tI_m-\rho(A(z)).
\end{equation}
The invariant quadratic form is normalized so that
$q(H)=1$. In the defining two-dimensional representation,
$A(z)^2=q(z)I_2$.

Let $\mathcal H_l$ be the space of harmonic homogeneous
polynomials of degree $l$ in $z_1,z_2,z_3$. It is the
irreducible representation of highest weight $2l$ and
dimension $2l+1$. The harmonic decomposition reads
\begin{equation}\label{G15:harmonic}
 \mathbb C[z_1,z_2,z_3]
     =\bigoplus_{l\geq0}\mathcal H_l\otimes\mathbb C[q].
\end{equation}
One can obtain this directly by applying the Laplacian to
$q^b h_l$: for harmonic $h_l$,
\[
 \Delta(q^b h_l)=2b(2l+2b+1)q^{b-1}h_l.
\]
This gives the direct decomposition and its dimensions in
characteristic zero. The indicated irreducibility is the
standard highest-weight calculation for the adjoint
three-dimensional representation.

\paragraph{1. Canonical harmonic matrix covariants.}
Write the multiplicity-free decomposition of
$\operatorname{End}(V)$ as
$\bigoplus_{l=0}^{m-1}\mathcal E_l$, where
$\mathcal E_l\simeq\mathcal H_l$, and let $\pi_l$ be its
equivariant projection. The trace Hermitian form makes
these distinct summands orthogonal. The polynomial map
\[
 A\longmapsto\pi_l\bigl(\rho(A)^l\bigr)
\]
is nonzero and has harmonic coefficient functions of
degree $l$. Indeed the symmetrized multiplication map
from the $l$th symmetric power of the adjoint
representation to $\operatorname{End}(V)$ is equivariant.
Its highest-weight vector $E^l$ maps to
$\rho(E)^l\ne0$ for $l<m$, and this vector has highest
weight $2l$, so lies in $\mathcal E_l$. The harmonic
decomposition says that this is the unique copy of
highest weight $2l$ in symmetric degree $l$.

The value of this covariant at $H$ is nonzero. Otherwise
equivariance and homogeneity would make it vanish on
every scalar multiple of the semisimple orbit of $H$,
which is the dense open set $q\ne0$, contradicting its
nonvanishing. This value is diagonal, since it has
Cartan weight zero.

These diagonal vectors have an elementary description.
Put the uniform positive inner product on functions on
$\{\lambda_0,\ldots,\lambda_{m-1}\}$:
\[
 \langle f,g\rangle=\frac1m\sum_{j=0}^{m-1}
                          f(\lambda_j)\overline{g(\lambda_j)}.
\]
Let $p_0,\ldots,p_{m-1}$ be the real orthonormal
polynomials obtained from $1,x,\ldots,x^{m-1}$, with
positive leading coefficients; thus $p_0=1$.
The zero-weight vector in $\mathcal E_l$ is proportional
to $p_l(\rho(H))$. To check the assertion, note that
$\rho(H)^r$ lies in the sum of the zero-weight spaces of
$\mathcal E_0,\ldots,\mathcal E_r$ by the symmetric-power
decomposition. The top projection is nonzero by the
preceding argument. Induction therefore identifies their
successive spans with those of $1,x,\ldots,x^r$ on the
spectrum. Orthogonality of the endomorphism summands
is precisely orthogonality for the displayed uniform
inner product. This proves the assertion, including
its degree and nonzero leading coefficient.

Normalize the covariant $T_l(A)$ so that
$T_l(H)=p_l(\rho(H))$. Its coefficient functions span
all of $\mathcal H_l$, since it is a nonzero map between
the corresponding irreducible representations.
Equivariance and \eqref{G15:harmonic} now give uniquely
\begin{equation}\label{G15:adjexpansion}
 \pi_l\operatorname{adj} M(t,z)=T_l(A(z))F_l(t,q),
 \qquad 0\leq l<m,
\end{equation}
where $F_l$ is a polynomial homogeneous of weighted
degree $m-1-l$ for $\deg t=1$, $\deg q=2$.

\paragraph{2. The scalar factors are trailing Jacobi determinants.}
Multiplication by $x$ on the functions on the $m$ nodes,
expressed in the orthonormal basis $p_0,\ldots,p_{m-1}$,
is a real symmetric tridiagonal matrix $J$.
Tridiagonality follows because $xp_l$ is orthogonal
to every $p_r$ with $r<l-1$. Its off-diagonal entries
$b_1,\ldots,b_{m-1}$ are strictly positive: they are
the ratios of the positive leading coefficients of
successive orthonormal polynomials. The nodes and
measure are symmetric under $x\mapsto-x$, so $p_l$
has parity $l$. Thus every diagonal entry
$\langle xp_l,p_l\rangle$ is zero.

Let $J^{[l+1]}$ denote the trailing principal submatrix
on indices $l+1,\ldots,m-1$, with the empty determinant
equal to one, and define
\[
 Q_l(t)=\det(tI-J^{[l+1]}),\qquad 0\leq l<m.
\]
The spectrum of $J$ is the node set. If
$D(t)=\prod_j(t-\lambda_j)$, evaluating
\eqref{G15:adjexpansion} at $H$ and taking its
orthonormal trace coefficient gives
\[
 F_l(t,1)=\frac1m\sum_jp_l(\lambda_j)
                         \frac{D(t)}{t-\lambda_j}.
\]
On the other hand the spectral expression for the
$(l,0)$ entry of $(tI-J)^{-1}$ is
\[
 \frac1m\sum_j\frac{p_l(\lambda_j)}{t-\lambda_j}.
\]
Expansion of the corresponding tridiagonal cofactor
shows that its numerator is
$b_1\cdots b_l Q_l(t)$: the unique path from index
zero to index $l$ contributes the indicated product,
and all remaining rows and columns give the trailing
principal determinant. Therefore
\begin{equation}\label{G15:jacobi}
 F_l(t,1)=b_1\cdots b_l Q_l(t).
\end{equation}
In particular every factor is nonzero. Weighted
homogeneity then determines $F_l(t,q)$ from this
univariate identity. All $Q_l$ have the parity of
their degree $m-1-l$.

We record the elementary spectral facts needed below.
Every nonempty trailing block is real symmetric,
tridiagonal, and has positive off-diagonal entries.
An eigenvector is determined by its first coordinate
using the three-term recurrence; that coordinate cannot
be zero. Hence all eigenvalues are simple. Conjugation
by the alternating-sign diagonal matrix negates the
block, so its eigenvalues occur in opposite pairs.
An even-sized block has no zero eigenvalue, as its
determinant at zero is a nonzero signed product of
alternate squared off-diagonal entries. An odd-sized
block has a simple zero eigenvalue. These assertions
also hold for blocks of sizes zero or one with the
obvious interpretations.

\paragraph{3. Cofactor multiples span every third jet.}
Let $S=\mathbb C[t,z_1,z_2,z_3]$, and let $\mathcal C_m$
be the degree-$m$ space spanned by all linear-form
multiples of the entries of $\operatorname{adj} M$.
It is the differential image of the determinant-pencil
map at \eqref{G15:pencil}. Formula
\eqref{G15:adjexpansion} implies that it contains
\[
 tF_l(t,q)\mathcal H_l\qquad(0\leq l<m).
\]
It also contains
\begin{equation}\label{G15:lowerproduct}
 qF_{l+1}(t,q)\mathcal H_l\qquad(0\leq l\leq m-2).
\end{equation}
For clarity, multiplication of a harmonic polynomial of
degree $l+1$ by a linear form has a nonzero projection
onto $q\mathcal H_l$. Indeed
$\Delta(z_i h)=2\partial_i h$; these derivatives span
$\mathcal H_l$ by equivariance and nonvanishing.
The space $\mathcal C_m$ is invariant, so projecting
its products onto this irreducible summand stays in
$\mathcal C_m$. Multiplication by the invariant
$F_{l+1}$ gives \eqref{G15:lowerproduct}.
Similarly the top harmonic component of
$\mathcal H_1\mathcal H_{m-1}$ supplies
$\mathcal H_m\subseteq\mathcal C_m$, since
$F_{m-1}$ is a nonzero constant.

The quotient $S_m/t^4S_{m-4}$ is exactly the third-jet
space, with $S_{m-4}=0$ if $m<4$. Its harmonic
decomposition has, for each $0\leq l\leq m$, the
multiplicity space
\begin{equation}\label{G15:targetspace}
 \operatorname{span}\{t^j q^{(D-j)/2}:
         0\leq j\leq\min\{3,D\},\ j\equiv D\pmod2\},
 \qquad D=m-l,
\end{equation}
tensored with $\mathcal H_l$. It has dimension one
for $D=0,1$, and dimension two for $D\geq2$.
The cases $D=0$ and $D=1$ are already supplied by
$\mathcal H_m$ and $tF_{m-1}\mathcal H_{m-1}$.

For $D\geq2$, specialize $q=1$. On the fixed
multiplicity space \eqref{G15:targetspace}, this is
an isomorphism onto the corresponding powers of $t$
modulo $t^4$. By \eqref{G15:jacobi}, it suffices to
show that $tQ_l(t)$ and $Q_{l+1}(t)$ are independent
modulo $t^4$.

If $D$ is even, $Q_l$ has odd degree $D-1$ and a
simple root at zero, while $Q_{l+1}$ has even degree
$D-2$ and nonzero value at zero. Thus $tQ_l$ starts
with a nonzero $t^2$ term, while $Q_{l+1}$ has a
nonzero constant term. They are independent in the
even two-dimensional space modulo $t^4$.

If $D$ is odd, then $D\geq3$. The block
$K=J^{[l+1]}$ is real symmetric and invertible,
of even size $D-1$. Write $e$ for its first
coordinate vector. The diagonal resolvent identity gives
\[
 \frac{Q_{l+1}(t)}{Q_l(t)}
    =e^{\mathsf T}(tI-K)^{-1}e.
\]
The alternating-sign symmetry fixes $e$ and negates
$K$, so all diagonal entries of odd powers of
$K^{-1}$ vanish. Expanding the inverse at zero gives
\[
 \frac{Q_{l+1}(t)}{tQ_l(t)}
   =-e^{\mathsf T}K^{-2}e
     -t^2 e^{\mathsf T}K^{-4}e
     +O(t^4).
\]
The coefficient of $t^2$ is strictly negative, since
$e^{\mathsf T}K^{-4}e=\|K^{-2}e\|^2>0$. Thus
$Q_{l+1}$ cannot be a constant multiple of $tQ_l$
modulo $t^4$. This proves independence in the odd
two-dimensional space, completing all cases of
\eqref{G15:targetspace}.

Consequently the differential image $\mathcal C_m$
surjects onto $S_m/t^4S_{m-4}$. The determinant-pencil
parameter map followed by third-jet projection has a
surjective differential at the explicit point
\eqref{G15:pencil}. In characteristic zero its image
is dense. This proves third-jet dominance for every
$m\geq2$, and hence for all $m\geq1$.

For $m\geq4$ the third-jet target dimension is
\[
 \sum_{j=0}^3\binom{m-j+2}{2}=2m^2+2.
\]
The fourth-jet target adds
$\binom{m-2}{2}>0$, whereas G-06 bounds
$\dim\mathcal D_{m,4}\leq2m^2+2$. The fourth-jet
map, and therefore every higher full jet map, cannot
be dominant. This establishes the maximal-order
assertion.

Finally apply the coefficient-weight argument of
G-14 with the now proved third-jet dominance. A
degree-$d$ coefficient monomial surviving on
$x_0^a q$ has $x_0$ weight
$ad+j_1+\cdots+j_d$. If $\lambda_1\leq ad+3$,
its evaluation depends only on the first four
coefficient forms $q_0,q_1,q_2,q_3$. Their dense
realization by determinants and block padding force
the same isotypic ideal containment as in G-14.
This proves the multiplicity inequality, including
the elementary cases $m<3$ where all available
jets are included. The rectangular inequality is
the equivalent condition
$nd/4>ad+3$, namely $d(4m-3n)>12$.
\end{proof}

\paragraph{What this resolves.}
The argument proves the needed rank for one explicit
pencil in every size and hence proves dominance. It
does not assume or assert a general strong Lefschetz
theorem for arbitrary Artinian determinantal algebras.
The positivity above is the squared norm for an
ordinary finite-dimensional real symmetric Jacobi
matrix; no arithmetic purity or RH input is used.
No literature-priority claim is made for this deduction.

\begin{proposition}[G-16: Explicit nonrectangular labels throughout the dimension-certified range]\label{G-16}
Let $n>m\geq3$, put $a=n-m$, and suppose
\begin{equation}\label{G16:dimensionhypothesis}
 \dim\mathcal P_{m,4}+3>2n^2+2.
\end{equation}
Define
\[
 b=\binom{n+3}{3}\binom{a+3}{3},\qquad L=n+3a,
 \qquad \beta=6bL^{30},\qquad u=2n^2-10,
\]
and set
\begin{equation}\label{G16:firstdegree}
 C=((n+a)\beta^2)^u,\qquad
 s_*=(u+1)(C-1)+1,\qquad d_* =\beta s_*.
\end{equation}
There is an integer $1\leq d\leq d_*$ with $4\mid md$ such that,
for the explicit shape
\begin{equation}\label{G16:label}
 \lambda_{m,n,d}=
 \left(ad+\frac{md}{4},\frac{md}{4},\frac{md}{4},\frac{md}{4}\right),
\end{equation}
one has
\[
 m_{\lambda_{m,n,d}}(\ell^a\operatorname{per}_m)
       >m_{\lambda_{m,n,d}}(\det_n).
\]
The bound has $\log d_*=O(n^2\log n)$ uniformly for $3\leq m<n$.
There is no additional condition $n<4m/3$ in this proposition.

An explicit progression of successful degrees is also available. Put
\[
 T=\binom{b+\beta}{b}-1,\quad E=4T(\beta!),\quad
 Q=((n+a)E)^{u-1},\quad K=u(Q-1)+1.
\]
For every integer $k\geq K$, the same inequality holds at
$d=Ek$ and the label \eqref{G16:label}.
Because $a>0$, every displayed label is nonrectangular.
\end{proposition}

\paragraph{Scope and sources.}
This applies to every padded pair certified by R-14, R-16, R-18,
R-19, or R-20, including those at or above $n=4m/3$, where G-12
excludes rectangular gaps. The construction uses a Segre--Veronese
coordinate ring and the same Kempf--Ness and reductive-quotient
facts used in G-11. For the quantitative bounds, we use Derksen
(2001), Theorem~1.1 and Proposition~1.2, printed p.~956, and the
image-degree/Chardin argument of G-10 and G-13. The new representation
requires a different coefficient-degree bound, which is checked below.

\begin{proof}
Let $V=\mathbb C^4$, $G=\mathrm{SL}(V)$, and put
\[
 Y_D=\mathcal D_{n,4},\qquad
 Y_P=\{L^a q:L\in V,\ q\in\mathcal P_{m,4}\}^{\rm cl}.
\]
These are the irreducible cones used in R-14, with
\[
 \dim Y_D=2n^2+2,\qquad
 \dim Y_P=\dim\mathcal P_{m,4}+3.
\]
For either cone $Y\subseteq\operatorname{Sym}^n V$, define
\begin{equation}\label{G16:enhancedcone}
 X_Y=\{F\otimes v^a:F\in Y,\ v\in V^*\}^{\rm cl}
 \subseteq W:=\operatorname{Sym}^n V\otimes\operatorname{Sym}^a V^*.
\end{equation}

\paragraph{1. The graded ring and the selected representation.}
The projectivization of \eqref{G16:enhancedcone} is the closed
Segre--Veronese image of $\mathbb P(Y)\times\mathbb P(V^*)$.
In particular $X_Y$ is irreducible and
\begin{equation}\label{G16:xdimension}
 \dim X_Y=\dim Y+3.
\end{equation}
Its coordinate ring has the exact grading
\begin{equation}\label{G16:segrering}
 \mathbb C[X_Y]_d=
       \mathbb C[Y]_d\otimes\operatorname{Sym}^{ad}V.
\end{equation}
Indeed each linear coordinate restricts to a product of a degree-one
coefficient function of $F$ and a degree-$a$ polynomial function of $v$.
Their degree-$d$ products span the right side: the first factor is
generated in degree one, and every degree-$ad$ monomial in $v$ can
be split into $d$ degree-$a$ monomials. The only additional relations
are exactly the restrictions defining this image. This proves
\eqref{G16:segrering} directly, without a projective-normality assumption
on $Y$.

To keep the duals explicit, use the actual coordinate-ring decomposition
\[
 \mathbb C[Y]_d=
       \bigoplus_{\lambda\vdash nd}m_\lambda(Y)S_\lambda V^*.
\]
It is the dual of the positive-weight orbit-span convention used
elsewhere in the paper. Schur's lemma for $G$ and
\eqref{G16:segrering} give
\begin{equation}\label{G16:extraction}
 \dim\mathbb C[X_Y]_d^G=
 \begin{cases}
 m_{\lambda_{m,n,d}}(Y),&4\mid md,\\
 0,&4\nmid md.
 \end{cases}
\end{equation}
For full detail,
\[
 (S_\lambda V^*\otimes\operatorname{Sym}^{ad}V)^G
   =\operatorname{Hom}_G(S_\lambda V,\operatorname{Sym}^{ad}V).
\]
This has dimension one precisely when
$\lambda=(ad+r,r,r,r)$ for an integer $r\geq0$, and otherwise is
zero. The fixed size $|\lambda|=nd$ forces $4r=md$, proving
the formula and ruling out any additional labels in this extraction.

\paragraph{2. Stable points with arbitrary padding.}
Let $x_0,x_1,x_2,x_3$ be an orthonormal basis of $V$ and
$e_0^*$ its first dual vector. Set
\[
 \gamma=\sqrt{\binom{n}{m}},\qquad
 F=x_0^a\bigl(x_0^m+\gamma x_1^m+
                         \gamma x_2^m+\gamma x_3^m\bigr),
 \qquad w=F\otimes(e_0^*)^a.
\]
The degree-$m$ factor of $F$ belongs to both
$\mathcal P_{m,4}$ and $\mathcal D_{m,4}$. In fact each of its two
binary summands factors into $m$ linear factors over $\mathbb C$;
put those factors on the diagonal and the cycle of an $m\times m$
pencil. Its permanent is the sum of the two products, and its
determinant is the same sum after absorbing $(-1)^{m-1}$ into one
cycle factor. Block padding then gives $F\in Y_P\cap Y_D$.
Thus the same tensor $w$ lies in both enhanced cones.

Use the Hermitian structures induced from tensor powers. The four
monomial tensors in $w$ are mutually orthogonal and have equal norm:
$\|x_0^n\|=1$, while
$\|x_0^a x_i^m\|^2=a!m!/n!$ for $i>0$, which is canceled by
$\gamma^2$. Their torus weights are exactly
\[
 ne_0-ae_0=me_0,\qquad
 ae_0+me_i-ae_0=me_i\quad(1\leq i\leq3).
\]
Their mean weight is $(m/4)(1,1,1,1)$, so the traceless diagonal
moment map is zero. The off-diagonal terms vanish as well. On the
first tensor factor, the four monomial weights differ by
$m(e_i-e_j)$, never by one root when $m\geq3$. On the second
factor, an off-diagonal derivative is either zero or orthogonal
to $(e_0^*)^a$. Therefore $w$ has zero $G$ moment map.

Its projective stabilizer has zero Lie algebra. To verify this
without a generic-stabilizer assertion, suppose
$A\in\mathfrak{sl}_4$ and $A\cdot F=\kappa F$.
Write $F=x_0^a q$. Dividing by $x_0^{a-1}$ and reducing modulo
$x_0$ shows that $A\cdot x_0$ is a multiple $\eta x_0$, since
$a>0$ and $q|_{x_0=0}\ne0$. It follows that
$A\cdot q=(\kappa-a\eta)q$.
For the Fermat sum $q$, the monomials $x_j^{m-1}x_i$, $i\ne j$,
are all distinct because $m\geq3$. Their coefficients force every
off-diagonal entry of $A$ to vanish. The four pure-power coefficients
then force every diagonal entry to be the same, and trace zero
forces $A=0$. If $A$ stabilizes the projective tensor line of $w$,
the elementary tangent property of a nonzero decomposable tensor
forces it to stabilize the line of $F$; hence again $A=0$.
Thus the projective stabilizer is finite.

Kempf--Ness now makes $w$ a stable point of each enhanced cone.
Their stable loci are nonempty open sets. Using
\eqref{G16:xdimension}, the two affine invariant quotients therefore
have dimensions
\begin{equation}\label{G16:quotdims}
 \begin{split}
 \dim A_D&=u=2n^2-10,\\
 \dim A_P&=\dim\mathcal P_{m,4}-9>u,
 \end{split}
 \qquad
 A_D=\mathbb C[X_{Y_D}]^G,\quad A_P=\mathbb C[X_{Y_P}]^G.
\end{equation}
The subtraction is $15=\dim G$. The padding weights were canceled
in the tensor construction, so this calculation needs no inequality
between $n$ and $4m/3$.

\paragraph{3. A quantitative invariant bound for this representation.}
The dimension of $W$ is $b$. In the matrix coordinates on
$\mathrm{SL}_4$, the action on $V^*$ is given by the transpose
adjugate, whose entries have degree three. The action on $W$
therefore has matrix coefficients of degree at most
$L=n+3a$. Its kernel is finite: its differential is a nonzero
representation of the simple Lie algebra $\mathfrak{sl}_4$, hence
injective, and a zero-dimensional algebraic kernel is finite.
For example the vector $x_0^n\otimes(e_1^*)^a$ has a nonzero
diagonal Lie-algebra weight, verifying nontriviality.

Derksen's Proposition~1.2 applies to the embedding
$\mathrm{SL}_4\subseteq\mathbb A^{16}$ defined by the degree-four
equation $\det-1=0$. Thus its parameters are $t=16$, group dimension
$15$, equation-degree bound $4$, and action-degree bound $L$, giving
$\sigma(W)\leq4L^{15}$.
Theorem~1.1, with invariant-quotient dimension at most $b$, then gives
generators of $\mathbb C[W]^G$ of degrees at most
\[
 \max\{2,(3/8)b(4L^{15})^2\}\leq6bL^{30}=\beta.
\]
Exactness of invariants for this reductive group makes $A_D$ and
$A_P$ graded quotients of that invariant ring. Choose common
homogeneous generators $g_i$ of degrees $1\leq e_i\leq\beta$.

The determinant enhanced cone has a dense parametrization
\[
 (Z,v)\longmapsto\det(Z_1x_0+Z_2x_1+Z_3x_2+Z_4x_3)\otimes v^a.
\]
Its coordinates are homogeneous of degree $n+a$ in the
$4n^2+4$ parameters. This parameter degree is distinct from the
group-action coefficient degree $n+3a$ above.
The generator coordinates $g_i$ composed with this map have degree
$(n+a)e_i\leq(n+a)\beta$.
Homogenizing them with one additional source coordinate gives a
dominant rational map to the ordinary projective closure of
$\operatorname{Spec}A_D$ in its generator embedding. The image-degree
argument of G-10, including its treatment of a base locus, yields
\[
 \deg\overline{\operatorname{Spec}A_D}
      \leq((n+a)\beta)^u.
\]

Let $H_D^{\leq}(d)=\sum_{j=0}^d\dim(A_D)_j$, and define
$H_P^{\leq}$ similarly. The graded pieces through degree $d$
are contained in the ordinary degree-at-most-$d$ generator filtration,
because all generator weights are at least one. Chardin's bound on
the ordinary projective closure therefore gives
\begin{equation}\label{G16:cumulativeupper}
 H_D^{\leq}(d)\leq((n+a)\beta)^u\binom{d+u}{u}.
\end{equation}
By \eqref{G16:quotdims}, at least $u+1$ images of the generators in
$A_P$ are algebraically independent. Their monomials of ordinary
total degree at most $s$ have graded degree at most $\beta s$;
hence
\begin{equation}\label{G16:cumulativelower}
 H_P^{\leq}(\beta s)\geq\binom{s+u+1}{u+1}.
\end{equation}
As in G-13, the exact binomial comparison is
\[
 \frac{\binom{s+u+1}{u+1}}{\binom{\beta s+u}{u}}
 =\frac{s+u+1}{u+1}\prod_{i=1}^u\frac{s+i}{\beta s+i}
 \geq\frac{s+u+1}{(u+1)\beta^u}.
\]
At $s=s_*$ this is strictly greater than
$C/\beta^u=((n+a)\beta)^u$. Equations
\eqref{G16:cumulativeupper}--\eqref{G16:cumulativelower} give a
strict cumulative gap through $d_*$. Degree zero contributes one
on both sides, so there is a positive degree $d\leq d_*$ with
$\dim(A_P)_d>\dim(A_D)_d$.
Formula \eqref{G16:extraction} identifies this gap with exactly the
label \eqref{G16:label} and forces $4\mid md$.
The four-variable highest-weight restriction of G-05/R-14 transfers
it to the full padded-permanent and determinant orbit closures.
Since $a<n$, $b=O(n^6)$, $L=O(n)$, and $\beta=O(n^{36})$;
\eqref{G16:firstdegree} consequently has
$\log d_*=O(n^2\log n)$.

\paragraph{4. An explicit progression, if individual degrees are desired.}
There are at most $T=\binom{b+\beta}{b}-1$ generators, since the
ambient polynomial space in positive degrees through $\beta$ has
that dimension. The elementary generator-grouping proof in G-11
therefore makes both $E$th Veronese rings standard graded for
$E=4T(\beta!)$. Their dimensions are still $u$ and at least $u+1$.
The degree-one invariant functions of the first Veronese, composed
with the enhanced determinant parametrization, have degree
$(n+a)E$. The projective image has dimension $u-1$ and degree at
most $Q=((n+a)E)^{u-1}$.
Chardin's theorem and linear Noether normalization give respectively
\[
 \dim(A_D)_{Ek}\leq Q\binom{k+u-1}{u-1},\qquad
 \dim(A_P)_{Ek}\geq\binom{k+u}{u}.
\]
Their binomial ratio is $(k+u)/u$, strictly larger than $Q$ for
$k\geq u(Q-1)+1$. This proves the claimed progression; $4\mid E$
ensures all labels in it are integral.
\end{proof}

\paragraph{Examples where ordinary rectangles are excluded.}
The pairs $(m,n)=(15,20)$ and $(16,22)$ from R-17 satisfy
\eqref{G16:dimensionhypothesis}. G-16 produces labels
$\lambda=(35r,15r,15r,15r)$ at degrees $d=4r$ for the first pair,
and labels $\lambda=(10d,4d,4d,4d)$ for the second, at the
successful degrees specified above. G-12 precludes ordinary
rectangular gaps for both pairs. The same explicit label family
applies to the unbounded nonrectangular strip supplied by R-20.


\begin{theorem}[G-17: Exact criterion for four-row stretched directions]\label{G-17}
Assume $n>m\geq3$, put $a=n-m$, and suppose
\[
 \dim\mathcal P_{m,4}+3>2n^2+2.
\]
Fix $s\geq1$ and a partition $\nu\vdash ns$ of length at most four,
appending zero parts when necessary. There is a positive integer $j$
such that
\begin{equation}\label{G17:gap}
 m_{j\nu}(\ell^a\operatorname{per}_m)>m_{j\nu}(\det_n),
 \qquad\text{in coordinate degree }sj,
\end{equation}
if and only if
\begin{equation}\label{G17:criterion}
 \nu_4>0\qquad\hbox{and}\qquad\nu_1>as.
\end{equation}
If either condition fails, no positive stretch $j$ gives a gap.
If both hold, an explicit progression of stretches gives gaps at
every sufficiently large term, as specified below. This does not
assert a gap at every sufficiently large integer stretch.

Here are effective bounds in the affirmative case. Put
\[
 \eta=(\nu_1-\nu_4,\nu_2-\nu_4,\nu_3-\nu_4,0),\qquad e=|\eta|,
 \qquad r_\eta=\#\{(i,j):1\leq i<j\leq4,\ \eta_i>\eta_j\},
\]
and define
\[
 \begin{aligned}
 b&=\binom{\binom{n+3}{3}+s-1}{s}\dim S_\eta\mathbb C^4,
 &u&=2n^2+r_\eta-13,\\
 L&=ns+3e,&p&=ns+e,\qquad \beta=6bL^{30}.
 \end{aligned}
\]
The Schur dimension is explicitly
$\prod_{i<j}(\eta_i-\eta_j+j-i)/(j-i)$.
One successful stretch satisfies
\begin{equation}\label{G17:firststretch}
 1\leq j\leq
 j_*:=\beta\bigl[(u+1)((p\beta^2)^u-1)+1\bigr].
\end{equation}
For a progression, let
\[
 T=\binom{b+\beta}{b}-1,\quad E=4T(\beta!),\quad
 Q=(pE)^{u-1},\quad K=u(Q-1)+1.
\]
Then \eqref{G17:gap} holds for every stretch $j=Ek$ with $k\geq K$.
\end{theorem}

\paragraph{Meaning of the criterion.}
For any one of the dimension-certified padded pairs, this determines
exactly which rational directions in the four-row partition simplex
contain a stretched multiplicity obstruction. Its two excluded faces
come from low-row determinant survival and the padding weight bound.
In particular this is a classification of all such directions,
not only a further special choice of the one-longer-row labels in G-16.
The coordinate degrees and stretch factors remain large; no efficient
multiplicity-positivity algorithm is asserted.

\begin{proof}
Necessity follows from earlier ideal containments and includes equality.
If $\nu_4=0$, every $j\nu$ has at most three rows, and G-05 gives
the full ambient multiplicity on the determinant side. If
$\nu_1\leq as$, then in coordinate degree $d=sj$ one has
$(j\nu)_1\leq ad$. The coefficient-weight argument of G-14 applies
and gives the opposite weak multiplicity inequality for every $j$.
Thus it remains to prove sufficiency under \eqref{G17:criterion}.

Use $V=\mathbb C^4$ and the two degree-$n$ cones $Y_D,Y_P$ of G-16.
We first give the auxiliary flag construction, including its coordinate
ring, so that no projective-normality or multiplicity-free assertion
is left implicit.

\paragraph{1. The auxiliary cone and its Cartan coordinate ring.}
Let $v_\eta$ be a nonzero extremal vector of weight $-\eta$ in
$S_\eta V^*$. Concretely, put
$c_i=\eta_i-\eta_{i+1}$ for $i=1,2,3$ and take the product of
the vectors
\[
 (e_0^*\wedge\cdots\wedge e_{i-1}^*)^{\otimes c_i}.
\]
It lies in the single Cartan component $S_\eta V^*$ of the indicated
tensor product of exterior powers. Its projective orbit records the
nested subspaces of dimensions $i$ for which $c_i>0$.
This is the corresponding partial flag variety, a closed incidence
subvariety of a product of Grassmannians. Its dimension is
$r_\eta$: its stabilizer is the block parabolic for the equal-part
blocks of $\eta$, and the missing off-diagonal blocks contribute
one parameter for each pair counted in $r_\eta$.
If $\eta=0$, the projective orbit is a point and $r_\eta=0$.

Let $\mathcal C_\eta$ be the affine cone over this projective orbit.
Its degree-$d$ coordinate ring is
\begin{equation}\label{G17:cartanring}
 \mathbb C[\mathcal C_\eta]_d=S_{d\eta}V.
\end{equation}
Here is a direct representation-theoretic proof. The dual of its
degree-$d$ coordinate space is the span of the group translates of
$v_\eta^{\otimes d}$. This is an extremal vector of weight $-d\eta$.
Its weight space at the extreme of the tensor-product weight order
is one dimensional, and it generates exactly the irreducible Cartan
component $S_{d\eta}V^*$. Taking the dual proves
\eqref{G17:cartanring}. Products give its usual Cartan multiplication.
When $\eta=0$, the cone is a line and each graded piece is
one dimensional, so the same formula holds.

For either cone $Y$, define the enhanced Veronese cone
\begin{equation}\label{G17:enhancement}
 X_Y=\{F^{\otimes s}\otimes v:F\in Y,\ v\in\mathcal C_\eta\}^{\rm cl}
 \subseteq W:=\operatorname{Sym}^s(\operatorname{Sym}^n V)
                      \otimes S_\eta V^*.
\end{equation}
Its projectivization is the Segre product of the $s$th Veronese
embedding of $\mathbb P(Y)$ with the partial flag variety. Hence
\begin{equation}\label{G17:dimension}
 \dim X_Y=\dim Y+r_\eta.
\end{equation}
Pullback of degree-$d$ coordinate functions gives exactly
\begin{equation}\label{G17:ring}
 \mathbb C[X_Y]_d=\mathbb C[Y]_{sd}\otimes S_{d\eta}V.
\end{equation}
The first factor is the Veronese subring of the standard graded
coordinate ring of $Y$, and the second is \eqref{G17:cartanring}.
Both are generated by their degree-one pieces in these gradings;
products of paired coordinates span their diagonal tensor product.
The pullback into the product coordinate ring is injective, proving
\eqref{G17:ring} without assuming that $Y$ is normal.

Write the actual degree-$sd$ coordinate representations as
$\bigoplus_{\lambda\vdash nsd}m_\lambda(Y)S_\lambda V^*$.
The invariant multiplicity in \eqref{G17:ring} is
\[
 \dim\operatorname{Hom}_{\mathrm{SL}_4}
                   (S_\lambda V,S_{d\eta}V).
\]
It is one if $\lambda=d\eta+t(1,1,1,1)$, and zero otherwise.
The size $nsd$ forces
$4t=(ns-e)d=4\nu_4d$, so $\lambda=d\nu$.
Consequently
\begin{equation}\label{G17:extraction}
 \dim\mathbb C[X_Y]_d^{\mathrm{SL}_4}=m_{d\nu}(Y).
\end{equation}
This is a single-label extraction, with original coordinate degree
$sd$, rather than a sum over different labels or degrees.

\paragraph{2. Stability for every strict direction in the criterion.}
Define four positive rational numbers
\[
 \theta_0=\frac{\nu_1/s-a}{m},\qquad
 \theta_i=\frac{\nu_{i+1}}{sm}\quad(1\leq i\leq3).
\]
Positivity is exactly \eqref{G17:criterion}, and their sum is one.
Let $u_0=x_0^n$ and
$u_i=\sqrt{\binom nm}\,x_0^a x_i^m$ for $i>0$.
These are orthonormal monomial vectors for the tensor-induced metric.
Put
\[
 F=\sum_{i=0}^3\sqrt{\theta_i}\,u_i,
 \qquad w=F^{\otimes s}\otimes v_\eta,
\]
normalizing $v_\eta$ to have norm one. The degree-$m$ factor of $F$
is a Fermat sum with four nonzero coefficients. The binary factorization
and cyclic-pencil construction in G-16 therefore put $F$ in both
$Y_D$ and $Y_P$, so $w$ belongs to both enhanced cones.

The diagonal moment of $F$, before removing the scalar part, is
\[
 a e_0+m\sum_i\theta_i e_i=\nu/s.
\]
The moment of a unit pure tensor power is $s$ times the moment of
the original vector; this follows by applying the Lie operator to
each factor and taking the tensor inner product. Thus the first
factor of $w$ contributes $\nu$. The auxiliary extremal vector
contributes its weight $-\eta$. Their sum is
$\nu-\eta=\nu_4(1,1,1,1)$, whose traceless part is zero.
Off-diagonal moments vanish on $F$ because its monomial weights
differ by $m$ times a root, with $m\geq3$, and on $v_\eta$ because
it is a weight vector. They therefore vanish on the tensor product.

The projective stabilizer of this $F$ has zero Lie algebra by the
mixed-monomial calculation in G-16; all four coefficients remain
nonzero, as required there. A Lie operator stabilizing the line of
$F^{\otimes s}$ must stabilize the line of $F$: in the symmetric
algebra on $\operatorname{Sym}^n V$ the equation
$sF^{s-1}(A\cdot F)=cF^s$ implies $A\cdot F=(c/s)F$.
A projective stabilizer of the decomposable tensor $w$ consequently
stabilizes that same line. Thus its projective stabilizer is finite.
Kempf--Ness gives a stable point of each enhanced cone.

The invariant quotients have dimensions
\begin{equation}\label{G17:quotdims}
 \begin{split}
 \dim A_D&=2n^2+r_\eta-13=u,\\
 \dim A_P&=\dim\mathcal P_{m,4}+r_\eta-12>u,
 \end{split}
 \qquad A_D=\mathbb C[X_{Y_D}]^{\mathrm{SL}_4},\quad
 A_P=\mathbb C[X_{Y_P}]^{\mathrm{SL}_4}.
\end{equation}
Indeed stable loci are nonempty and open, so one subtracts the
common orbit dimension $15$ from \eqref{G17:dimension}.
The original dimension gap is preserved for every flag type.

\paragraph{3. The two polynomial degrees and the effective bounds.}
The dimension of $W$ is $b$. Its group-action matrix coefficients
have degree at most $L=ns+3e$ in the entries of $g\in\mathrm{SL}_4$:
the first factor has degree $ns$, while $S_\eta(g^{-\mathsf T})$
has degree at most $3e$ because each entry of $g^{-\mathsf T}$ is
a cubic cofactor. The representation has finite kernel, since its
kernel is contained in the finite projective stabilizer of $w$.
Derksen's Theorem~1.1 and Proposition~1.2 therefore give generators
of $\mathbb C[W]^{\mathrm{SL}_4}$ of degrees at most
$\beta=6bL^{30}$, exactly as in G-16. Exactness of invariants makes
$A_D,A_P$ graded quotients with the same bound on their generators.

For the image-degree bound a smaller, different parameter degree
is valid, and the reason is important. Introduce an independent
matrix $B\in\operatorname{End}(V^*)$, without imposing any relation
between $B$ and the determinant-pencil parameters. The map
\[
 B\longmapsto S_\eta(B)v_\eta
\]
is a homogeneous polynomial map of degree $e$. One can see this
directly in the tensor construction of the Schur functor: it is
the restriction of $B^{\otimes e}$ to the relevant Schur component.
The invertible $B$ form a dense open set and give the full
$\mathrm{GL}(V^*)$ orbit, which is dense in $\mathcal C_\eta$
when $e>0$. Its scalar matrices supply every nonzero scalar
multiple, because the base field is $\mathbb C$.
Thus
\begin{equation}\label{G17:parameterization}
 (Z,B)\longmapsto
  \det(Z_0x_0+\cdots+Z_3x_3)^{\otimes s}
       \otimes S_\eta(B)v_\eta
\end{equation}
has dense image in $X_{Y_D}$ and is homogeneous of degree
$p=ns+e$ in its $4n^2+16$ independent parameters.
No inverse-matrix constraint or adjugate is needed for this
parametrization. If $e=0$, the second factor is constant and can
be omitted; all scalar multiples are already supplied by scaling
the first cone. Formula \eqref{G17:parameterization} then remains
valid with degree $p=ns$.

Apply the cumulative generator-filtration comparison from G-16,
now with quotient dimension $u$, generator bound $\beta$, and
parameter degree $p$. For clarity, its two inequalities become
\[
 \sum_{d=0}^{D}\dim(A_D)_d
       \leq(p\beta)^u\binom{D+u}{u},\qquad
 \sum_{d=0}^{\beta h}\dim(A_P)_d
       \geq\binom{h+u+1}{u+1}.
\]
The first uses the ordinary projective closure in the generator
embedding, of degree at most $(p\beta)^u$, and Chardin's theorem.
The second uses $u+1$ algebraically independent generator images.
The same exact binomial ratio is at least
$(h+u+1)/((u+1)\beta^u)$. Taking
$h=(u+1)((p\beta^2)^u-1)+1$ yields a strict cumulative gap through
$j_*=\beta h$. Since degree zero agrees, some positive stretch
$j\leq j_*$ has a strict invariant-dimension gap. Equation
\eqref{G17:extraction} proves \eqref{G17:gap} at that stretch.

For the individual progression, there are at most
$T=\binom{b+\beta}{b}-1$ homogeneous invariant generators.
The explicit grouping argument of G-11 makes the $E$th Veronese
rings standard graded for $E=4T(\beta!)$.
Their dimensions are $u$ and at least $u+1$, and composing their
degree-one invariant coordinates with \eqref{G17:parameterization}
gives parameter degree $pE$. Consequently their degree-$k$ Hilbert
functions satisfy
\[
 \dim(A_D)_{Ek}\leq Q\binom{k+u-1}{u-1},\qquad
 \dim(A_P)_{Ek}\geq\binom{k+u}{u},\qquad Q=(pE)^{u-1}.
\]
The ratio is strictly greater than one once
$k\geq u(Q-1)+1$. Applying \eqref{G17:extraction} gives the
specified successful multiples $j=Ek$.
On this progression the determinant-side multiplicity is positive as
well. The positive-dimensional graded domain $A_D$ has a nonzero
positive-degree generator $g$ of some degree $e_0\leq\beta$.
Since $e_0\mid E$, the elements $g^{Ek/e_0}$ are nonzero for every
$k\geq1$. Equation \eqref{G17:extraction} therefore gives
$m_{Ek\nu}(\det_n)\geq1$. Thus the displayed strict gaps compare
two positive multiplicities and are not occurrence gaps.
Finally the four-variable highest-weight restriction transfers
all the multiplicities back to the full orbit closures.
\end{proof}

\paragraph{Consistency with the earlier boundaries.}
If $\nu$ is rectangular, the strict inequality in
\eqref{G17:criterion} is $n/4>a$, equivalently $n<4m/3$,
recovering G-11/G-12 whenever the dimension hypothesis holds.
For the one-longer-row direction in G-16, it holds for all padding
because the extra first-row component compensates for $a$.
At equality $\nu_1=as$, every stretch remains excluded, rather
than merely failing this stability construction.
G-15 further excludes a finite initial segment when
$j(\nu_1-as)\leq3$; this is consistent with the eventual
progression in the strict interior.


\paragraph{Quantitative corollary: the first successful stretch has polynomial binary length.}
In the hypotheses and notation of G-17, put
$b_n=\binom{n+3}{3}$. The effective first-stretch bound satisfies
\begin{equation}\label{G17:bitlength}
 \log_2j_*
 \leq\log_2(u+1)+u\log_2p+(2u+1)\log_2\beta
 =O\!\left(n^5\log(s+1)+n^2\log n\right).
\end{equation}
For fixed $s$, the sharper estimate is $O_s(n^2\log n)$.
Thus a successful stretch, its original coordinate degree $sj$, and
its highest-weight entries admit polynomial binary length in $n$ and
$\log(s+1)$. This is a size bound for an existing successful stretch;
it does not provide an efficient search for that stretch or an efficient
certificate of the required multiplicity inequality.

\begin{proof}
The weak compositions of $s$ into $b_n$ parts are determined by their
first $b_n-1$ entries, each between zero and $s$. The Weyl dimension
formula for $\operatorname{GL}_4$ has six factors, each at most $ns+3$,
because $\eta_1\leq ns$. Consequently the ambient dimension in G-17 obeys
\[
 b=\binom{b_n+s-1}{s}\dim S_\eta\mathbb C^4
   \leq(s+1)^{b_n-1}(ns+3)^6.
\]
Moreover $|\eta|\leq ns$, so
\[
 L\leq4ns,\qquad p\leq2ns,\qquad u\leq2n^2,
 \qquad
 \beta\leq6(s+1)^{b_n-1}(ns+3)^6(4ns)^{30}.
\]
The exact bound
$j_*=\beta[(u+1)((p\beta^2)^u-1)+1]$
is at most $\beta(u+1)(p\beta^2)^u$. Taking logarithms gives
the first inequality in \eqref{G17:bitlength}, and $b_n=O(n^3)$
gives its stated order. If $s$ is fixed, instead use
$\binom{b_n+s-1}{s}=O_s(n^{3s})$; then $\log\beta=O_s(\log n)$.
Finally multiplication by $s$ or by a highest-weight entry bounded
by $ns$ adds only $O(\log(ns))$ bits.
\end{proof}


\begin{theorem}[G-18: Every large stretch in every strict four-row direction]\label{G-18}
Let $n>m\geq3$, $a=n-m$, and let $\nu\vdash ns$, $s\geq1$,
have at most four rows. Assume
\[
 \nu_4>0,\qquad \nu_1>as.
\]
Use $\eta,r_\eta,b,L,p,\beta,T,E$ from G-17, defined by the same
formulas even when its dimension inequality is not assumed. Put
\[
 c_m=\dim\mathcal P_{m,4},\qquad
 u=2n^2+r_\eta-13,\qquad v=c_m+r_\eta-12.
\]
The following conclusions hold in original coordinate degree $sj$.

First, both representations occur in every sufficiently large
integer stretch, with the uniform displayed bound
\begin{equation}\label{G18:occurrence}
 m_{j\nu}(\ell^a\operatorname{per}_m)\geq1,
 \qquad m_{j\nu}(\det_n)\geq1
 \qquad(j\geq\beta^2).
\end{equation}
This conclusion needs no dimension comparison between the two cones.

If $v>u$, equivalently $c_m+3>2n^2+2$, set
$Q_D=(pE)^{u-1}$. Then
\begin{equation}\label{G18:everylargepositive}
 m_{j\nu}(\ell^a\operatorname{per}_m)>m_{j\nu}(\det_n)
 \qquad\text{for every integer }j\geq E(4uQ_D+1).
\end{equation}
Thus every affirmative direction of G-17, including the rectangular
ones, gives a gap at every sufficiently large integer stretch.

There is also an effective converse range. Set
\[
 \bar v=4m^2-2m+r_\eta-10,\qquad Q_P=(pE)^{\bar v-1}.
\]
If
\begin{equation}\label{G18:reversecondition}
 2n^2>4m^2-2m+3,
\end{equation}
then
\begin{equation}\label{G18:everylargereverse}
 m_{j\nu}(\det_n)>m_{j\nu}(\ell^a\operatorname{per}_m)
 \qquad\text{for every integer }j\geq E(4\bar vQ_P+1).
\end{equation}
Condition \eqref{G18:reversecondition} holds, in particular,
whenever $n\geq\sqrt2\,m$.

More precisely, for every such fixed direction,
there is a constant $\kappa_\nu>0$ such that
\begin{equation}\label{G18:ratio}
 \frac{m_{j\nu}(\ell^a\operatorname{per}_m)}
      {m_{j\nu}(\det_n)}
  =\kappa_\nu j^{\,c_m-2n^2+1}(1+O(j^{-1}))
  \qquad(j\longrightarrow\infty)
\end{equation}
through all positive integers. When the exponent is zero, the
dimension comparison alone does not determine the sign of the gap.

For every fixed partition direction $\nu\vdash ns$ with at most
four rows, including those outside the two strict inequalities,
there is no occurrence obstruction at any sufficiently large
stretch. Here an occurrence obstruction means positive
padded-permanent multiplicity and zero determinant multiplicity.
\end{theorem}

\paragraph{Sources and the extra step.}
The enhanced coordinate rings and the initial stable tensor are
proved in G-17. We additionally use separation of distinct closed
orbits by invariants: Mukai (2003), Theorem~5.3, printed pp.~161--162,
and Corollary~5.4, p.~162. Openness of the stable locus is
Proposition~5.15, p.~166. The new step removes the finite projective
stabilizer by moving the padding line, then makes the resulting
degree-semigroup argument effective. The original coordinate degree
is $sj$; in a rectangular direction $ns=4\nu_4$ already incorporates
the necessary central congruence.

\begin{proof}
We retain the notation $X_{Y_D},X_{Y_P},W,A_D,A_P$ of G-17.
Their construction, common stable tensor, generator bound, and
dimensions $u,v$ do not require the inequality $v>u$.

\paragraph{1. A common stable tensor with no nontrivial scalar orbit symmetry.}
Write the balanced padded Fermat form of G-17 as $x_0^a f$, where
$f=\sum_{i=0}^3\gamma_i x_i^m$ and every $\gamma_i$ is nonzero.
The degree-$m$ hypersurface $f=0$ is smooth: its partial derivatives
have no common projective zero. It is irreducible, since two
positive-degree projective hypersurface components in $\mathbb P^3$
intersect, and their intersection would be singular in the product.
Its projective stabilizer $H\subseteq\mathrm{PGL}_4$ is finite.
Indeed, the equation $A\cdot f=cf$ forces all off-diagonal entries
of $A$ to vanish by the distinct monomials $x_i^{m-1}x_j$,
$i\ne j$, because $m\geq3$; its diagonal entries are then equal.
Thus the projective stabilizer has zero Lie algebra, so in
characteristic zero it is finite.

Keep $f$ and the auxiliary vector $v_\eta$ fixed, and vary a nonzero
padding linear form $\ell$. For every such $\ell$, the tensor
\[
 w_\ell=(\ell^a f)^{\otimes s}\otimes v_\eta
\]
belongs to both enhanced cones. This follows from the same binary
factorization and cyclic permanent/determinant representations of
$f$ used in G-17, followed by the independent padding block.
The inverse image of the stable locus of either cone under this
regular map is a nonempty open set, because it contains $\ell=x_0$.
Stability is unchanged by a nonzero scalar, so these sets descend
to nonempty open subsets of the projective space of linear forms.

For every nonidentity $h\in H$, its fixed locus in this projective
space is proper and closed: a linear transformation fixing every
line of the dual representation is scalar. The intersection of
the two stable opens with the complements of these finitely many
fixed loci is nonempty, by irreducibility of projective space.
Choose $[\ell]$ in this intersection. Because $a>0$ and $m\geq3$,
unique factorization implies that every projective automorphism of
$F=\ell^a f$ preserves the unique linear factor $\ell$ and the
unique irreducible factor $f$ of degree greater than one separately.
Its image therefore lies in $H$ and fixes $[\ell]$, so it is the
identity in $\mathrm{PGL}_4$.

Set $w=w_\ell$. A projective stabilizer of this decomposable tensor
preserves the line of $F^{\otimes s}$ and hence the line of $F$:
the Veronese map is injective on projective points. Thus its
projective stabilizer in $\mathrm{SL}_4$ is exactly the center
$\mu_4$. This center acts
on $W$ by the exponent
\[
 ns-|\eta|=4\nu_4,
\]
and hence acts trivially. It follows that
\begin{equation}\label{G18:noscaling}
 gw=cw,\quad g\in\mathrm{SL}_4,\ c\in\mathbb C^*
 \quad\Longrightarrow\quad c=1.
\end{equation}
The orbits of $w$ and of every $cw$ are closed: stability holds,
and scalar multiplication commutes with the group action.

\paragraph{2. Common nonzero invariant degrees have greatest common divisor one.}
Choose homogeneous generators $g_1,\ldots,g_t$ of
$\mathbb C[W]^{\mathrm{SL}_4}$, with positive degrees
$e_i\leq\beta$, as in G-17. Let
\[
 J=\{i:g_i(w)\ne0\}.
\]
This set is nonempty. Otherwise all positive-degree invariants
would vanish at $w$, so the distinct closed orbits of $w$ and zero
would fail to be separated.

Suppose the integer $q=\gcd\{e_i:i\in J\}$ exceeded one, and
choose a primitive $q$th root $c\ne1$. Homogeneity gives
$g_i(cw)=g_i(w)$ for every generator: for $i\in J$ use
$c^{e_i}=1$, and for $i\notin J$ both values are zero. The same
equality therefore holds for every invariant polynomial. Closed
orbit separation makes the two closed orbits equal, contrary to
\eqref{G18:noscaling}. Thus
\begin{equation}\label{G18:gcd}
 \gcd\{e_i:i\in J\}=1.
\end{equation}
Every product of these generators is nonzero at $w$, and consequently
has nonzero image in both $A_D$ and $A_P$.

Here is the required elementary numerical-semigroup bound. If
positive integers $e_i\leq\beta$ have greatest common divisor one,
their nonnegative semigroup contains every integer $N\geq\beta^2$.
Let $a_0$ be the least generator. In the directed graph on
$\mathbb Z/a_0\mathbb Z$ with edges adding the $e_i$, every residue
is reachable: a subsemigroup of a finite group is a subgroup, and
the greatest-common-divisor condition makes this subgroup the whole
group. A shortest path to each residue has no repeated vertex and
uses at most $a_0-1$ edges. Its integer weight $c_r$ therefore
satisfies $c_r\leq(a_0-1)\beta<\beta^2$. For $N\geq\beta^2$,
choose the representative of its residue; $N-c_r$ is a nonnegative
multiple of $a_0$. This proves the assertion, including $a_0=1$.

Applying this to \eqref{G18:gcd}, for every $N\geq\beta^2$ there
is a homogeneous ambient invariant $h_N$ of degree $N$ nonzero
at $w$. Its restrictions to both cones are nonzero. The exact
single-label identity of G-17 now proves \eqref{G18:occurrence}.

\paragraph{3. Transfer from the Veronese progression to every residue.}
We have $T\geq\beta$, because the positive-degree polynomial
space through degree $\beta$ contains the powers of one variable.
Since $\beta\geq2$ and $\beta!\geq\beta$, it follows that
$E=4T\beta!\geq\beta^2$.

For $j=Ek+r$, $0\leq r<E$, multiplication by the common nonzero
invariants $h_{E+r}$ and $h_{2E-r}$ is injective in each of the
domains $A_D,A_P$. In particular, for either $A$ and every $k\geq1$,
\begin{equation}\label{G18:sandwich}
 \dim A_{E(k-1)}\leq\dim A_{Ek+r}\leq\dim A_{E(k+2)}.
\end{equation}
Only existence of these invariants is used; no algorithm for
computing a minimal generating set is being asserted.

Suppose first $v>u$. G-17's standard-Veronese dimension and
image-degree bounds give
\[
 \dim(A_P)_{E(k-1)}\geq\binom{k+u-1}{u},\qquad
 \dim(A_D)_{E(k+2)}\leq Q_D\binom{k+u+1}{u-1}.
\]
The exact ratio of the binomials is
\[
 \frac{\binom{k+u-1}{u}}{\binom{k+u+1}{u-1}}
 =\frac{k(k+1)(k+2)}{u(k+u)(k+u+1)}
 \geq\frac{k}{4u}\qquad(k\geq u+1).
\]
For $k\geq4uQ_D+1$ this is strictly greater than $Q_D$.
Combining with \eqref{G18:sandwich} proves
\eqref{G18:everylargepositive} in every residue class.

For the converse bound, R-13's row--column torus gives
$c_m\leq4m^2-2m+2$, and hence $v\leq\bar v$.
Condition \eqref{G18:reversecondition} says exactly $u>\bar v$.
The enhanced permanent cone has a dense homogeneous
parametrization of degree $p=ns+|\eta|$:
\[
 (Z,L,B)\longmapsto
 \bigl(L^a\operatorname{per}_m(Z_0x_0+\cdots+Z_3x_3)\bigr)^{\otimes s}
       \otimes S_\eta(B)v_\eta.
\]
Here the pencil coefficients, coefficients of the linear form $L$,
and entries of $B\in\operatorname{End}(V^*)$ are independent.
The first factor has parameter degree $ns$, exactly as for the
determinant parametrization. Therefore the projective image of
the $E$th invariant Veronese has degree at most
$(pE)^{v-1}\leq Q_P$. The same Chardin upper bound, enlarged to
dimension $\bar v$, and Noether-normalization lower bound give
\[
 \dim(A_P)_{E(k+2)}\leq Q_P\binom{k+\bar v+1}{\bar v-1},
 \qquad
 \dim(A_D)_{E(k-1)}\geq\binom{k+\bar v-1}{\bar v}.
\]
Use \eqref{G18:sandwich} with the roles reversed and the preceding
binomial calculation with $\bar v$ in place of $u$.
This proves \eqref{G18:everylargereverse}. If
$n\geq\sqrt2\,m$, then
$2n^2-(4m^2-2m+3)\geq2m-3>0$, as required.

\paragraph{4. The all-integer asymptotic ratio.}
For a standard graded domain $B$ of dimension $q\geq2$, write
$e(B)=\deg\operatorname{Proj}B>0$. Its Hilbert function satisfies
\[
 \dim B_k=\frac{e(B)}{(q-1)!}k^{q-1}+O(k^{q-2}).
\]
This is the usual Hilbert-polynomial normalization used in G-10.
One can check the leading term directly after linear Noether
normalization: $B$ is a finite torsion-free module over a polynomial
ring in $q$ degree-one variables; a homogeneous basis over its
fraction field spans a free submodule of rank $e(B)$, and the
torsion quotient has growth $O(k^{q-2})$. An annihilating homogeneous
polynomial for that quotient and its finite generating set give
this error bound from the polynomial-ring Hilbert function.

Apply this formula to $B_D=A_D^{(E)}$ and $B_P=A_P^{(E)}$.
Their dimensions are $u,v$, both at least two here. Indeed
$c_m\geq20$: R-13 gives $c_3=20$, and a block diagonal pencil
with a $3\times3$ block and $m-3$ entries $x_0$ embeds
$x_0^{m-3}\mathcal P_{3,4}$ in $\mathcal P_{m,4}$.
Also $r_\eta\geq0$, including the rectangular case $\eta=0$.
Sandwich
\eqref{G18:sandwich}, uniformly over its finitely many residues,
gives for $A$ of dimension $q$
\[
 \dim A_j=\frac{e(A^{(E)})}{E^{q-1}(q-1)!}
                 j^{q-1}+O(j^{q-2}).
\]
Taking the quotient of these two positive leading terms proves
\eqref{G18:ratio}, with
\[
 \kappa_\nu=
 \frac{e(B_P)}{e(B_D)}\,
 \frac{(u-1)!}{(v-1)!}\,E^{u-v}>0,
 \qquad v-u=c_m-2n^2+1.
\]
All multiplicities are identified with the full orbit-closure
ones by the same four-variable highest-weight restriction as in
G-17. The final occurrence assertion follows from
\eqref{G18:occurrence} in the strict region. Outside that region,
G-05 and G-14 give the opposite weak multiplicity inequality at
every stretch, so an occurrence obstruction is already impossible.
\end{proof}

\paragraph{The precise remaining range.}
The reverse inequality concerns sufficiently large stretches of
each fixed direction. It does not exclude an individual
finite-degree gap, give a threshold uniform over changing
directions, or prove containment of orbit closures. Rectangular
directions obey the padding boundary of G-12 as well as the
all-stretch conclusion in the strict range. When
$c_m+3=2n^2+2$, the leading quotient degrees, and
possibly lower Hilbert coefficients, are the remaining comparison.


\subsection{A specified rectangular discriminant copy}

\begin{theorem}[G-19: an explicit rectangular discriminant module and its positive energy]
Let $m\geq3$, put $N=m^2$, and work over $\mathbb C$.
Here $V=\mathbb C^N$ is the argument space, forms lie in
$\operatorname{Sym}^m V^*$, and coefficient functions of degree
$D$ lie in $\operatorname{Sym}^D(\operatorname{Sym}^m V)$.
We use $g\cdot F=F\circ g^{-1}$, hence
$(g\cdot h)(F)=h(F\circ g)$.
Let $k\geq5$ satisfy
\[
 k\leq \operatorname{codim}_{\mathbb A^N}
                 \operatorname{Sing}(\operatorname{per}_m).
\]
In particular $k=5$ is allowed for every $m\geq3$, and $k=6$
is allowed for every $m\geq6$, by the established bounds in R-B17.
Define
\[
 D=k(m-1)^{k-1},\qquad \delta=m(m-1)^{k-1},\qquad
 \lambda=(\delta^k).
\]
There is an explicit highest-weight polynomial $h$ of coefficient
degree $D$ and weight $\lambda$, namely the discriminant of the
restriction to the first $k$ variables, with these properties.

Its generated irreducible module $W$ is contained in the ideal of
the determinant orbit closure, while its restriction to the
permanent orbit closure is nonzero and therefore injective on $W$.
With this explicit dual convention, $W$ has positive polynomial
type $S_\lambda V$. For normalized Haar measure, and any
$U(N)$-invariant Hermitian inner product on $W$, one has
\begin{equation}\label{G19:Haar}
 \int_{U(N)}|h(g\cdot F)|^2\,dg
   =\frac{\|h\|^2}{\dim W}\sum_{j=1}^{\dim W}|h_j(F)|^2,
 \qquad
 \dim W=\prod_{i=1}^k\prod_{j=k+1}^{N}
                     \frac{\delta+j-i}{j-i},
\end{equation}
where $(h_j)$ is an orthonormal basis of $W$.
This integral is zero at $F=\det_m$ and strictly positive at
$F=\operatorname{per}_m$.

There is also a positive integer Gaussian certificate. Normalize
the discriminant to have primitive integer coefficients, let
$A:\mathbb C^k\longrightarrow\mathbb C^N$ have independent
matrix entries, and write
\[
 H_m(A)=\operatorname{Disc}_{k,m}(\operatorname{per}_m\circ A)
        =\sum_{|\alpha|=mD}c_\alpha A^\alpha.
\]
Then $H_m$ is a nonzero integer polynomial, and for independent
standard complex Gaussian entries $Z$ normalized by
$\mathbb E|Z_{ij}|^2=1$,
\begin{equation}\label{G19:Gaussian}
 \mathbb E|H_m(Z)|^2=\sum_\alpha c_\alpha^2\alpha!\ \geq1.
\end{equation}
For the determinant in place of the permanent, the polynomial
and expectation are identically zero. There exists an integer
matrix $A$ with all entries in $\{0,1,\ldots,mD\}$ and
$H_m(A)\ne0$.

This separates a specified irreducible copy, not the total
multiplicities of its isomorphism type. It proves neither
$m_\lambda(\operatorname{per}_m)>m_\lambda(\det_m)$ nor
$m_\lambda(\det_m)=0$. Moreover the analogous degree-$n$
discriminant construction vanishes on every positively padded
form $\ell^{n-m}f_m$, so it supplies no such separator when
$n>m$.
\end{theorem}

\paragraph{Known inputs and the precise deduction.}
The geometric determinant input is Alper, Bogart, and Velasco
(2017), Theorem~1.7, author version printed p.~3, with proof
on p.~5: every linear determinant in at least five variables
defines a singular projective hypersurface. The permanent
singular-locus bound is the cited R-B17 Theorem~4.14, strengthened
by its Theorem~4.23 for $m\geq6$. The discriminant and its
invariance are Mukai (2003), Definition~5.20 and Corollary~5.21,
printed p.~170. Its degree and change-of-variable weight follow
from Benoist (2012), Theorem~1.3, equations (1.4)--(1.6),
printed pp.~1190--1191, specialized to one equation; the
coefficient-degree specialization is also Example~1.8, p.~1192.
The second displayed formula of that example omits the factor
$d_1$ (written $m$ here) in the variable-change degree; we use equation~(1.4),
Theorem~1.3(1.6), and scalar dilation for the correct weight.
These known theorems are not reproved. The deduction below
identifies a concrete rectangular module and its Haar/Gaussian
certificates, and states explicitly why this does not compare
the total multiplicities or survive padding. No historical
priority claim is made for the underlying discriminant separator.

\begin{proof}
Let $X=\{\operatorname{per}_m=0\}\subset\mathbb P^{N-1}$.
The codimension hypothesis implies that a general projective
$(k-1)$-plane misses $\operatorname{Sing}X$: its expected
intersection dimension is at most
$(k-1)+(N-1-k)-(N-1)=-1$.
Iterated Bertini on the smooth locus then makes its intersection
with $X$ smooth. This is the same linear-section argument used
in the cited proof of Alper--Bogart--Velasco Theorem~1.7, now
applied to the permanent's singular-locus bound. Consequently
\[
 \operatorname{Disc}_{k,m}(\operatorname{per}_m\circ A)\ne0
\]
for $A$ in a nonempty Zariski-open subset of the space of linear
maps. In contrast, the cited determinant theorem makes
$\operatorname{Disc}_{k,m}(\det_m\circ A)=0$ for every $A$.
If a restriction is the zero form, its discriminant is zero as
well, so no exceptional lower-rank map changes this conclusion.

The discriminant has coefficient degree $D=k(m-1)^{k-1}$.
Under substitution $x\mapsto Bx$ it transforms by
$\det(B)^\delta$, where $k\delta=mD$; thus
$\delta=m(m-1)^{k-1}$. Regard this determinant-character vector
as an element of the plethysm in $N$ variables using the first
$k$ coordinate vectors. It is killed by the raising operators
inside $\mathrm{GL}_k$, because it is a determinant character,
and by the raising operators from later coordinate indices,
because no such index occurs. It is therefore a highest-weight
vector of weight $(\delta^k,0^{N-k})$. Semisimplicity shows that
the module generated by this one vector is one irreducible copy
$W$ of that type, even if the ambient plethysm contains other
copies of the same type. The displayed Weyl dimension formula
is the usual formula with all equal-row factors canceled.

Every translate of $h$ vanishes on $\det_m$, and hence $W$
vanishes on its entire orbit closure. A full-rank map $A$ with
smooth permanent restriction can be completed to an invertible
change of the $N$ variables. Therefore some translate of $h$
is nonzero at $\operatorname{per}_m$. The restriction map from
$W$ to the permanent coordinate ring is equivariant and nonzero;
irreducibility makes it injective.

Apply the Schur-orthogonality projector identity used in R-02
and R-05 to this particular irreducible copy. It gives
\eqref{G19:Haar}. At least one basis vector evaluates nontrivially
at the permanent, so its energy is strictly positive. Every
basis vector evaluates to zero at the determinant. This proof
uses compact integration only after identifying the actual
copy and its two restriction maps.

Choose the primitive integer equation of the discriminant,
as permitted by the rational discriminant construction in
Benoist (2012), Theorem~1.7, p.~1191. The coefficients of
$\operatorname{per}_m\circ A$ are integer polynomials of degree
$m$ in the entries of $A$, so $H_m$ is integral and homogeneous
of degree $mD$. It is nonzero by the smooth-section argument.
The exact complex-Gaussian monomial orthogonality in R-06 gives
\eqref{G19:Gaussian}; some nonzero integer coefficient has
square at least one, and every $\alpha!$ is a positive integer.
Finally a nonzero polynomial of total degree at most $mD$
cannot vanish on the entire grid
$\{0,\ldots,mD\}^{kN}$, by successive univariate interpolation.
This proves the finite integer witness bound. It does not
give a polynomial-time search through that grid or an efficient
evaluation procedure for its permanent coefficients.

For the two limitations, the kernel on one specified copy $W$
does not determine the kernel on the whole $\lambda$-isotypic
space: other copies may survive on the determinant, while other
independent copies may vanish on the permanent. Thus no total
multiplicity inequality follows from this argument.
If $n>m$, every nonzero $k$-variable restriction of
$\ell^{n-m}f_m$ has a linear factor and a positive-degree
remaining factor. If the linear factor repeats, its hypersurface
is singular along that hyperplane. Otherwise the hyperplane
meets the remaining positive-degree hypersurface in projective
space, and all first derivatives of the product vanish at their
intersection. Here $k\geq5$ ensures that this intersection is
nonempty. A zero restriction again has zero discriminant.
Hence the degree-$n$ discriminant vanishes on the entire padded
orbit closure. The displayed module and energy cannot remove
this padding obstruction by themselves.
\end{proof}

\paragraph{An explicit smooth five-variable cubic slice.}
For $m=3$, the following matrix provides a concrete witness to
the permanent nonvanishing in G-19:
\[
 M(x)=\begin{pmatrix}
 x_1+x_2&x_3&x_4\\
 x_4&x_0&x_1\\
 x_2+x_3&x_1&x_0+x_2
 \end{pmatrix}.
\]
Writing $f=\operatorname{per}M$, direct expansion gives
\begin{align*}
 f={}&x_0^2x_1+x_0^2x_2+x_0x_1x_2+x_0x_2^2
        +x_0x_2x_4+2x_0x_3x_4\\
    &+x_1^3+x_1^2x_2+x_1x_2x_3+x_1x_3^2
        +x_1x_4^2+x_2x_3x_4.
\end{align*}
Here is a finite integer certificate of its smoothness, with the
complete matrix and two independently written reconstructions
included in the work archive. Put
$S=\mathbb Z[x_0,x_1,x_2,x_3,x_4]$, and form the coefficient
matrix of
\[
 \bigoplus_{j=0}^4 S_4\longrightarrow S_6,
 \qquad (g_j)_j\longmapsto\sum_{j=0}^4g_j\partial_{x_j}f.
\]
Order each monomial basis by decreasing lexicographic order on
its exponent tuple. The matrix has $210$ rows and $350$ columns;
column $70j+b$ corresponds to the $b$th degree-four monomial
times $\partial_{x_j}f$, with all indices starting at zero.
Keep all rows and select the following multiplier indices in
each derivative block:
\begin{center}
\begin{tabular}{r p{0.82\linewidth}}
\toprule
$j$ & Selected values of $b$ (inclusive intervals)\\
\midrule
0 & $0$--$69$\\
1 & $0$--$12$, $14$--$22$, $24$--$28$, $30$--$31$,
    $34$--$42$, $44$--$48$, $50$--$51$, $54$--$58$,
    $60$--$61$, $64$--$65$, $69$\\
2 & $0$--$12$, $15$--$22$, $25$--$28$, $35$--$42$,
    $45$--$48$, $55$--$58$\\
3 & $0$--$10$, $12$, $15$--$20$, $22$, $25$,
    $35$--$40$, $42$, $45$\\
4 & $0$--$8$, $10$, $12$, $15$--$18$, $35$\\
\bottomrule
\end{tabular}
\end{center}
The resulting $210\times210$ determinant, in the displayed
increasing row and column orders, is $94$ modulo $101$.
It is therefore a nonzero integer. After extension to
$\mathbb C$, every degree-six monomial belongs to the ideal
generated by the five partial derivatives. A common nonzero
projective zero of those derivatives would make all degree-six
monomials vanish, including the sixth power of a nonzero
coordinate, which is impossible. Thus $\{f=0\}\subset\mathbb P^4$
is smooth. This proof uses neither a numerical tolerance nor an
inference from a randomly sampled point.

The author reconstruction uses exact polynomial expansion; the
independent reconstruction multiplies all six permutations in
integer polynomial dictionaries, differentiates and convolves
the dictionaries, reproduces every one of the $73{,}500$ matrix
entries modulo $101$, and separately eliminates the selected
minor. The files are
\texttt{work/root/g19\_cubic/certificate.json},
\texttt{work/root/g19\_cubic/macaulay\_degree6.u8},
\texttt{work/root/check\_g19\_cubic\_slice.py}, and
\texttt{work/isolation/check\_G19\_cubic\_independent.py}.
The row-major byte matrix has SHA-256
\texttt{c739e7d83b6c43a3e48e6f20c92667fd77282af7e1240c3806d55989af0928e5}.

Consequently the specified-copy separator at $m=3$ has
coefficient degree $D=80$ and rectangular weight $(48^5)$.
Five is the minimum possible number of rows for such a
highest-weight separating equation: R-11 shows that the
four-variable determinant-cubic image is dense in the whole
cubic coefficient space. A highest-weight equation with at
most four rows lies in the corresponding coefficient
subalgebra, so vanishing on the determinant forces it to be
zero. G-19 and the displayed smooth slice supply a nonzero
five-row equation. This row-minimality assertion concerns
specified polynomial equations, and does not assert a total
multiplicity gap at weight $(48^5)$.


\begin{theorem}[G-20: The exact eventual four-row multiplicity phase]
\label{G-20}
Let $n>m\ge3$, put $a=n-m$, and define
\[
 C_m=\min\left\{\binom{m+3}{3},\,4m^2-2m+2\right\},
 \qquad N_m=\left\lfloor\sqrt{\left\lfloor C_m/2\right\rfloor}\right\rfloor.
\]
For a fixed partition direction $\nu\vdash ns$, $s\ge1$, with
at most four rows, assume the strict conditions
\[
 \nu_4>0,\qquad \nu_1>(n-m)s.
\]
Then exactly one of the following holds:
\begin{enumerate}
\item If $m<n\le N_m$, the padded-permanent multiplicity is
strictly larger than the determinant multiplicity in every
sufficiently large integer stretch $j\nu$.
\item If $n>N_m$, the determinant multiplicity is strictly
larger in every sufficiently large integer stretch $j\nu$.
\end{enumerate}
In both cases the original coordinate degree is $sj$, both
multiplicities are positive for every $j\ge\beta^2$ with
$\beta$ from G-17, and
\begin{equation}\label{G20:exact_ratio}
 \frac{m_{j\nu}(\ell^a\operatorname{per}_m)}
      {m_{j\nu}(\det_n)}
 =\kappa_\nu j^{C_m-2n^2+1}\bigl(1+O(j^{-1})\bigr),
 \qquad \kappa_\nu>0.
\end{equation}
There is no integer equality case in the exponent for the
stated values of $m,n$.

For $m\ge17$, the exact alternatives are
\begin{equation}\label{G20:large_boundary}
 \begin{aligned}
 n^2\le2m^2-m+1&\quad\Longrightarrow\quad
       \text{eventual permanent dominance},\\
 n^2\ge2m^2-m+2&\quad\Longrightarrow\quad
       \text{eventual determinant dominance}.
 \end{aligned}
\end{equation}
The small-size values are
\[
\begin{array}{c|rrrrrrr}
m&3&4&5&6&7&8&9\\ \hline
C_m&20&35&56&84&120&165&220\\
N_m&3&4&5&6&7&9&10
\end{array}
\]
\[
\begin{array}{c|rrrrrrr}
m&10&11&12&13&14&15&16\\ \hline
C_m&286&364&455&560&680&816&969\\
N_m&11&13&15&16&18&20&22.
\end{array}
\]

The corresponding exact criterion for eventual positive
four-row rectangular gaps is
\begin{equation}\label{G20:rectangular}
 m<n\le N_m,\qquad n<4m/3.
\end{equation}
Every admissible rectangular direction then has a positive
gap at every sufficiently large integer stretch. For every
$m\ge14$, condition \eqref{G20:rectangular} simplifies to
$m<n<4m/3$. The threshold $14$ is sharp for this uniform
eventual statement. No assertion is made that an individual
finite-degree gap is impossible in a direction that eventually
has determinant dominance.
\end{theorem}

\begin{proof}
R-13 and R-27 give $c_m=C_m$: sizes through sixteen fill the
four-variable form space, and every size at least seventeen
attains the row--column parameter bound. Substitute this exact
dimension in G-18. Its occurrence theorem gives positivity
of both multiplicities after $\beta^2$, and its ratio formula
becomes \eqref{G20:exact_ratio}.

The exponent is positive precisely when
$2n^2<C_m+1$, equivalently
$n^2\le\lfloor C_m/2\rfloor$. This is precisely $n\le N_m$.
For $m\ge17$, $C_m=4m^2-2m+2$ is even, so the exponent
$C_m-2n^2+1$ is odd and cannot vanish. The two adjacent integer
inequalities are exactly \eqref{G20:large_boundary}.
For $3\le m\le16$, the table is obtained by direct evaluation
of $\binom{m+3}{3}$. The only odd dimensions in that table are
$35,165,455,969$. Their half-successors are respectively
$18,83,228,485$, none of which is a square:
\[
 4^2<18<5^2,\quad9^2<83<10^2,\quad
 15^2<228<16^2,\quad22^2<485<23^2.
\]
Thus the exponent never vanishes in the small-size range
either. Its sign and the positive leading constant in
\eqref{G20:exact_ratio} prove the exact alternatives.

For completeness, the bounds are effective without waiting
for an unspecified asymptotic. With the notation of G-17 set
\[
 u=2n^2+r_\eta-13,\qquad v=C_m+r_\eta-12.
\]
If $v>u$, G-18 gives permanent dominance for every
\[
 j\ge E\bigl(4u(pE)^{u-1}+1\bigr).
\]
If $u>v$, its proof in paragraph~3 applies with the actual
dimension $v$ in place of the upper bound $\bar v$: the
permanent image degree is bounded by $(pE)^{v-1}$, and the
same standard-Veronese and residue-class sandwich gives
determinant dominance for every
\[
 j\ge E\bigl(4v(pE)^{v-1}+1\bigr).
\]
This is the same proved binomial comparison with the exact
dimension substituted, not a new claim about an uncomputed
leading coefficient.

For a rectangle $\nu=(q,q,q,q)$ with $4q=ns$, the first
strict condition is automatic and the second is
$ns/4>(n-m)s$, equivalently $n<4m/3$.
Within that range, the preceding exact sign criterion applies.
At or above its boundary, G-12 gives the opposite weak
multiplicity inequality at every stretch. This proves
\eqref{G20:rectangular}.

R-26 already proves that every $m\ge14$ and every
$m<n<4m/3$ lies in the positive dimension range; its four
small cases and uniform inequality are exact. To prove
sharpness for the eventual statement, take
\[
 m=13,\qquad n=17,\qquad s=4,\qquad\nu=(17,17,17,17).
\]
The direction is strict because $17>4(17-13)=16$.
Here $C_{13}=560$, and the ratio exponent is
\[
 C_{13}-2\cdot17^2+1=560-578+1=-17.
\]
Consequently the ratio is $\kappa_\nu j^{-17}(1+O(j^{-1}))$,
and the determinant multiplicity is larger at every sufficiently
large stretch. Any uniform threshold below fourteen would
include this pair, so it cannot have the asserted eventual
property. This argument does not decide any of its finitely
many earlier stretches.
\end{proof}

\paragraph{Scope.}
The theorem classifies the eventual sign along each fixed
four-row direction. It neither supplies a polynomial-time
positive extraction rule in growing ambient dimension nor
excludes changing-direction or isolated finite-degree gaps.
It does not establish containment of orbit closures, a general
permanent circuit lower bound, or a superlinear-padding GCT
separation. The hypotheses and explicit stretch bounds of
G-17/G-18 remain part of the conclusion.


\begin{theorem}[G-21: Maximal permanent-image dimension in every larger ambient space]
\label{G-21}
Let $m\ge17$ and $r\ge4$. Let $\mathcal P_{m,r}$ be the
Zariski closure of the forms obtained as permanents of
$m\times m$ pencils of linear forms in $r$ variables over
$\mathbb C$. Then
\begin{equation}\label{G21:dimension}
 \dim\mathcal P_{m,r}=rm^2-2m+2.
\end{equation}
Thus four variables already suffice to attain the full
parameter count modulo the row--column scaling directions;
each additional variable contributes exactly $m^2$ dimensions.
For a general coefficient specialization
$R=\mathbb C[X_{ij}]\longrightarrow S_r=\mathbb C[x_1,\ldots,x_r]$,
the corresponding graded first homology satisfies
\begin{equation}\label{G21:Tor}
 \operatorname{Tor}_1^R(R/J_m,S_r)_d=0
 \qquad(d\le m),
\end{equation}
where $J_m$ is the submaximal-permanent ideal used in R-23.
The assertion in \eqref{G21:Tor} concerns the specified graded
pieces and does not assert exactness of the entire restricted
resolution in higher degrees.
\end{theorem}

\begin{proof}
For a quaternary pencil $A$, write
$G_{ij}(A)=\partial\operatorname{per}(A)/\partial A_{ij}$.
R-23's cyclic-pencil proof shows that the $m^2$ degree-$(m-1)$
forms $G_{ij}(A)$ are linearly independent on a nonempty
Zariski open subset of the quaternary parameter space. R-27
shows that the differential image
\[
 \operatorname{span}\{x_kG_{ij}(A):1\le k\le4,\ 1\le i,j\le m\}
\]
has dimension $4m^2-2m+2$ on another nonempty open subset.
Their parameter space is irreducible, so choose one $A$
satisfying both conditions.

Regard this same pencil as an $r$-variable pencil with zero
coefficients on $x_5,\ldots,x_r$. The differential in every
new coefficient direction has the form $x_kG_{ij}(A)$,
$5\le k\le r$. For each fixed new variable these $m^2$ forms
are independent, because multiplication by $x_k$ is injective.
Their spans for different new variables, and the original
four-variable differential span, have disjoint multidegrees
in $x_5,\ldots,x_r$. Consequently the rank at this explicit
kind of parameter point is
\[
 (4m^2-2m+2)+(r-4)m^2=rm^2-2m+2.
\]
The generic image dimension is at least this rank. The
row--column torus has the same effective dimension $2m-2$
in every variable count, and gives the opposite upper bound
$rm^2-2m+2$, exactly as in R-13/R-23. This proves
\eqref{G21:dimension}. The use of a pencil with initially
unused variables is legitimate: a nonzero coefficient minor
at that point defines a nonempty open subset of the full
$r$-variable parameter space.

For the Tor statement, use the same first part of the graded
resolution and the same cited linear-syzygy theorem as in
R-23. In degree $m$, the relevant complex after specialization is
\[
 \mathbb C^{2m-2}\longrightarrow (S_r)_1^{m^2}
       \longrightarrow (S_r)_m.
\]
The old row--column relations remain independent on a nonempty
open subset: they are independent at the four-variable pencil
just used, regarded as an $r$-variable pencil. The last map
is the pencil differential and has the exact generic rank
\eqref{G21:dimension}. Its kernel therefore has dimension
$rm^2-(rm^2-2m+2)=2m-2$ and equals the image of the first
map. This proves the degree-$m$ assertion.
In degree $m-1$, the last map sends constants to the independent
$m^2$ cofactors. Below that degree its source is zero. These
give all lower-degree assertions in \eqref{G21:Tor}.
This tensor-complex argument remains meaningful when $r>m^2$;
in that case the coefficient map is an algebra homomorphism
into a larger polynomial ring, and is not called a quotient
by $m^2-r$ linear equations.
\end{proof}

\paragraph{Exact scope.}
This gives an image-dimension and low-degree derived-restriction
statement in all $r\ge4$. It does not compute individual
$r$-row multiplicities, identify all higher syzygies, or
turn the dimension statement into a polynomial-time positive
formula in growing $r$. Those are distinct remaining tasks.


\begin{proposition}[G-22: Fourier frames give stable common tensors in more rows]\label{G-22}
Let $n>m\ge r\ge2$, put $a=n-m$, and let
$\nu\vdash ns$, $s\ge1$, have $r$ positive parts. Suppose
\[
 \nu_r>0,\qquad \nu_1>as.
\]
Set $V=\mathbb C^r$, $\eta=\nu-\nu_r(1,\ldots,1)$,
$e=|\eta|$, and let $v_\eta\in S_\eta V^*$ be the unit
extremal flag vector of weight $-\eta$ used in G-17.
Then the two enhanced cones
\[
 X_Y=\{F^{\otimes s}\otimes v:F\in Y,
                  v\in\mathcal C_\eta\}^{\rm cl}
 \subseteq W=\operatorname{Sym}^s(\operatorname{Sym}^nV)
                         \otimes S_\eta V^*,
\]
where $Y=\mathcal D_{n,r}$ or
$Y=\{L^a q:q\in\mathcal P_{m,r}\}^{\rm cl}$,
have a common stable point for $\mathrm{SL}_r$.
Consequently their invariant coordinate rings have dimensions
\[
 \dim\mathbb C[X_Y]^{\mathrm{SL}_r}
 =\dim Y+r_\eta-(r^2-1),\qquad
 r_\eta=\#\{i<j:\eta_i>\eta_j\},
\]
and their degree-$j$ invariant dimensions equal the respective
multiplicities $m_{j\nu}(Y)$ in original coordinate degree $sj$.
No inequality comparing the dimensions of the two cones is needed.
\end{proposition}

\paragraph{Sources and the additional construction.}
The stability criteria are Losev (2006), Theorem~4, p.~1
(Kempf--Ness), and Mukai (2003), Theorem~7.4, printed p.~212,
with its linear-representation formulation (7.3), p.~213
and proof pp.~218--219 (Hilbert--Mumford). The exact quotient
locators are the same as in G-17. We prove the Fourier-frame
construction and preservation of every one-parameter minimum
weight under multiplication; the moment of a product of linear
forms is not assumed to equal the sum of their separate moments.

\begin{proof}
Write $x_0,\ldots,x_{r-1}$ for an orthonormal basis of $V$ and put
\[
 \mu_0=\frac{\nu_1}{s}-a,\qquad
 \mu_i=\frac{\nu_{i+1}}s\quad(1\le i<r).
\]
Every $\mu_i$ is positive and $\sum_i\mu_i=m$.
Choose a primitive $m$th root $\zeta$, and define the unit vectors
\begin{equation}\label{G22:frame}
 u_j=\sum_{i=0}^{r-1}\sqrt{\frac{\mu_i}{m}}\,
                    \zeta^{ij}x_i,\qquad 0\le j<m.
\end{equation}
The finite Fourier orthogonality identity gives
\begin{equation}\label{G22:framemoment}
 \|u_j\|^2=1,\qquad
 \sum_{j=0}^{m-1}u_ju_j^*=\operatorname{diag}(\mu_0,\ldots,\mu_{r-1}).
\end{equation}
Indeed $\sum_j\zeta^{(i-k)j}$ is $m$ for $i=k$ and zero
otherwise, since $0\le i,k<r\le m$.
Every $r$ of the vectors $u_j$ are independent: after dividing
the rows by their nonzero factors, their determinant is the
Vandermonde determinant on the distinct corresponding powers
of $\zeta$.

First work in the ordered tensor representation
\[
 \widetilde W=V^{\otimes ns}\otimes S_\eta V^*
\]
and consider the nonzero unit decomposable tensor
\[
 \xi=\bigl(x_0^{\otimes a}\otimes
             u_0\otimes\cdots\otimes u_{m-1}\bigr)^{\otimes s}
               \otimes v_\eta.
\]
The moment of a unit tensor product is the sum of the moments
of its factors: expand the infinitesimal action as a sum over
tensor positions and take the product Hermitian inner product.
By \eqref{G22:framemoment}, its diagonal contribution before
removing the trace is
\[
 sa e_0+s\mu-\eta=\nu-\eta=\nu_r(1,\ldots,1).
\]
The off-diagonal entries cancel in the Fourier sum. The
extremal weight vector $v_\eta$ has zero off-diagonal moment,
because distinct torus weights are orthogonal. Hence $\xi$
has zero $\mathrm{SL}_r$ moment map.

Its projective stabilizer is finite. A transformation preserving
the line of an ordered nonzero decomposable tensor preserves
the line of each tensor factor; this follows, for example, by
contracting all other factors with nonzero linear functionals.
If $m\ge r+1$, it consequently fixes $r+1$ full-spark lines
$[u_j]$. The first $r$ are a basis, so the transformation is
diagonal in that basis. Every coordinate of the next vector
in that basis is nonzero, since any $r$ of these $r+1$ vectors
are independent. Fixing its line makes all diagonal entries
equal. Thus the transformation is scalar.
If $m=r$, the $u_j$ are a basis and
\[
 \sum_{j=0}^{m-1}u_j=\sqrt{m\mu_0}\,x_0.
\]
All basis coefficients of $x_0$ are nonzero. Since $a>0$,
the ordered tensor also contains this padding vector and
the same argument again forces a scalar. The projective
stabilizer inside $\mathrm{SL}_r$ is therefore contained in
its finite center. Kempf--Ness now makes the affine orbit
of $\xi$ closed, and its finite stabilizer makes $\xi$ stable.

Let $q=\prod_{j=0}^{m-1}u_j\in\operatorname{Sym}^mV$,
$F=x_0^a q$, and
\[
 w=F^{\otimes s}\otimes v_\eta\in W.
\]
The diagonal $m\times m$ pencil with entries $u_j$ has both
permanent and determinant equal to $q$. Adding $a$ diagonal
entries $x_0$ proves membership of $F$ in both required
degree-$n$ cones. It remains to prove stability after the
symmetrizations; equivariance alone would not prove that.

For any nontrivial one-parameter subgroup
$\lambda:\mathbb G_m\to\mathrm{SL}_r$, diagonalize its action
and denote the smallest occurring integer weight of a nonzero
vector $z$ by $\omega_\lambda(z)$. Smallest weights add under
tensor product. They also add under multiplication of nonzero
polynomials in $\operatorname{Sym}(V)$: their smallest-weight
initial terms have nonzero product because this symmetric
algebra is an integral domain. The same statement applies
to the power $F^{\otimes s}$ in
$\operatorname{Sym}^s(\operatorname{Sym}^nV)$.
It follows that
\begin{equation}\label{G22:minweight}
 \omega_\lambda(w)
 =s\left(a\omega_\lambda(x_0)+
             \sum_j\omega_\lambda(u_j)\right)
                +\omega_\lambda(v_\eta)
 =\omega_\lambda(\xi).
\end{equation}
Both representations have scalar weight
$ns-e=r\nu_r>0$, so they are linear actions of ray type,
with irrelevant set the origin. Mukai's Theorem~7.4 says
that a nonzero vector is stable exactly when its limit
under every nontrivial $\lambda(t)$ as $t\to0$ fails to
exist. In the present weight notation, this means that
every $\omega_\lambda$ is negative. It holds for $\xi$,
and \eqref{G22:minweight} proves it for $w$. Thus $w$ is
the required common stable point.

For completeness the representation-theoretic extraction
uses the same argument as G-17 in rank $r$. The extremal
orbit cone has
$\mathbb C[\mathcal C_\eta]_j=S_{j\eta}V$ and projective
dimension $r_\eta$. The Segre--Veronese map consequently gives
\[
 \dim X_Y=\dim Y+r_\eta,\qquad
 \mathbb C[X_Y]_j=\mathbb C[Y]_{sj}\otimes S_{j\eta}V.
\]
These equalities follow from the span of translates of the
extremal tensor and from multiplication of paired coordinate
functions, as proved in G-17; they do not require normality.
The only degree-$sj$ type whose dual pairs to an invariant
with $S_{j\eta}V$ is
$\lambda=j\eta+t(1,\ldots,1)$.
The size equation $|\lambda|=nsj$ forces $t=j\nu_r$,
so that $\lambda=j\nu$. Exactness of invariants gives the
claimed multiplicity identity. The nonempty stable open
has orbit dimension $r^2-1$, which proves the quotient
dimension formula by the dimension-of-fibers theorem.
\end{proof}

\paragraph{Scope.}
The positive numbers $\mu_i$ need not be rational squares;
this is an existence witness over $\mathbb C$, not a claim
of polynomial bit complexity for a hitting set or for
multiplicity extraction. It works for every stated strict
direction and every $m\ge r$. Removing the scalar-orbit
period of the multiplied tensor requires a further step:
the trivial projective stabilizer of the ordered tensor
does not itself rule out permutations of the factors of $F$.



\begin{theorem}[G-23: Exact eventual multiplicity sign in every strict direction through $m$ rows]\label{G-23}
Let $n>m\ge17$ and $4\le r\le m$. Put $a=n-m$, and fix
$\nu\vdash ns$, $s\ge1$, with
\[
 \nu_r>0,\qquad \nu_1>as.
\]
Append no additional positive parts beyond row $r$.
Define
\[
 \begin{split}
 \eta&=\nu-\nu_r(1,\ldots,1),\qquad e=|\eta|,\qquad
 r_\eta=\#\{i<j:\eta_i>\eta_j\},\\
 u&=(r-2)n^2+r_\eta-r^2+3,\\
 v&=rm^2-2m+r+r_\eta-r^2+2,\\
 \Delta&=v-u=rm^2-(r-2)n^2-2m+r-1.
 \end{split}
\]
The integer $\Delta$ is never zero. Both multiplicities are
positive at every sufficiently large integer stretch, and
\begin{equation}\label{G23:ratio}
 \frac{m_{j\nu}(\ell^a\operatorname{per}_m)}
      {m_{j\nu}(\det_n)}
 =\kappa_\nu j^\Delta(1+O(j^{-1})),\qquad \kappa_\nu>0,
\end{equation}
through all integers $j\to\infty$, in original coordinate degree $sj$.
Thus the permanent multiplicity is eventually strictly larger
exactly when
\begin{equation}\label{G23:boundary}
 n\le N_{m,r}:=
 \left\lfloor\sqrt{\left\lfloor
       \frac{rm^2-2m+r-1}{r-2}\right\rfloor}\right\rfloor;
\end{equation}
for $n>N_{m,r}$ the determinant multiplicity is eventually strictly larger.

Here are explicit bounds. Let
\[
 \begin{aligned}
 B_n&=\binom{n+r-1}{r-1},&
 b&=\binom{B_n+s-1}{s}\dim S_\eta\mathbb C^r,\\
 L&=ns+(r-1)e,&p&=ns+e,\\
 \beta&=\left\lceil\frac{3r^2b}{8}\right\rceil
                          L^{2(r^2-1)},&
 T&=\binom{b+\beta}{b}-1,\qquad E=rT(\beta!).
 \end{aligned}
\]
The Schur dimension is the Weyl product
$\prod_{i<j}(\eta_i-\eta_j+j-i)/(j-i)$.
Both multiplicities are positive for every $j\ge\beta^2$.
If $\Delta>0$, the strict positive gap holds for every
\[
 j\ge E\bigl(4u(pE)^{u-1}+1\bigr).
\]
If $\Delta<0$, the strict reverse gap holds for every
\[
 j\ge E\bigl(4v(pE)^{v-1}+1\bigr).
\]
\end{theorem}

\paragraph{Proof dependencies and meaning.}
G-21 gives the permanent image dimension, R-28 gives the
determinant dimension and the projectively free common stable
tensor, and G-22 supplies stability for every stated direction.
The invariant-generator input is Derksen (2001), Theorem~1.1
and Proposition~1.2, printed p.~956. The Hilbert-function and
rational image-degree inputs are G-10 and its exact source
locators. The changes from G-18 are the growing row count,
the different group dimension, the new stable-point construction,
and the uniform arithmetic exclusion of equality. We give the
argument with all new parameters, rather than assuming that
the four-row constants remain valid.

\begin{proof}
Use the two enhanced cones of G-22, and let $A_D,A_P$ be
their invariant coordinate rings. They are graded domains.
Their degree-$j$ dimensions are the two specified multiplicities.
By R-28 and G-21,
\[
 \dim\mathcal D_{n,r}=(r-2)n^2+2,\qquad
 \dim\mathcal P_{m,r}=rm^2-2m+2.
\]
The padded permanent cone has dimension
$rm^2-2m+r+1$. Indeed its projective parametrization by a
padding line and a permanent form is generically injective:
on the nonempty open where the latter is irreducible of degree
$m>1$ and not divisible by the padding line, the unique
nonlinear factor and the linear factor recover both projective
inputs. The irreducible open is nonempty by the common cyclic
family in R-28. Thus padding adds $r-1$ affine-cone dimensions.
Subtracting $r^2-1$ from the enhanced dimensions in G-22 proves
\[
 \dim A_D=u,\qquad\dim A_P=v.
\]
Both are at least two: $r\le m<n$ gives
$u\ge2n^2-r^2+3>2$, and
$v\ge(r-1)m^2-2m+r+2>2$.

\paragraph{1. Invariant degrees and absence of a period.}
The ambient representation has dimension $b$. Its action
coefficients on $\mathrm{SL}_r$ have degree at most $L$:
the covariant factor contributes $ns$, and the entries of
the inverse transpose have degree $r-1$, contributing at
most $(r-1)e$ on the Schur factor. Its kernel is finite
because it is contained in the finite stabilizer of the
common stable tensor. In the embedding
$\mathrm{SL}_r\subset\mathbb A^{r^2}$, the defining equation
has degree $r$ and the group dimension is $r^2-1$.
Derksen's Proposition~1.2 therefore gives
\[
 \sigma(W)\le rL^{r^2-1},
\]
and Theorem~1.1 gives the displayed integer generator bound
$\beta$. Exactness of invariants makes $A_D,A_P$ graded
quotients of this same ambient invariant ring.

Choose homogeneous generators $g_i$ of degrees
$1\le e_i\le\beta$. R-28 supplies a common stable point
$w$ with
\[
 gw=cw,\quad g\in\mathrm{SL}_r,\ c\in\mathbb C^*
                  \quad\Longrightarrow\quad c=1.
\]
For precision, its projective stabilizer is the center
$\mu_r$, and this center acts with exponent
$ns-e=r\nu_r$, hence trivially.
The degrees of the generators nonzero at $w$ have greatest
common divisor one. If that divisor were $h>1$, a primitive
$h$th root $c$ would make every generator have equal values
at $w$ and $cw$, although their two closed orbits are distinct.
This contradicts closed-orbit separation by invariants
(Mukai, Theorem~5.3 and Corollary~5.4, pp.~161--162).
There is at least one such generator by the same separation
applied to the closed orbits of $w$ and zero.

The numerical-semigroup proof in G-18, part~2, now applies
with this actual $\beta$: shortest paths modulo the least
degree have at most $\beta-1$ edges, so each residue has a
representative of weight less than $\beta^2$. Hence every
integer $j\ge\beta^2$ is a sum of degrees nonzero at $w$.
The corresponding product gives an ambient homogeneous
invariant $h_j$ nonzero on both cones. This proves both
occurrence assertions at every such integer $j$.

\paragraph{2. Standard Veroneses and effective comparison.}
There are at most $T$ chosen positive-degree generators,
because the whole ambient coefficient-polynomial space
through degree $\beta$ has dimension $T$.
The generator-grouping argument of G-11 applies with
$E=rT\beta!$: remove pure generator blocks of degree
$\beta!$, leaving a residual monomial of degree less than
$T\beta!$, then complete it to a block of degree $E$ and
partition the remaining pure blocks. Thus $A_D^{(E)}$ and
$A_P^{(E)}$ are standard graded domains. Since $T\ge\beta$
and $\beta!\ge\beta$, we have $E\ge\beta^2$.

Both enhanced cones have dense homogeneous parametrizations
of degree $p$ in independent variables. They are
\[
 (Z,B)\longmapsto\det\left(\sum_{i=1}^r Z_i x_i\right)^{\otimes s}
                  \otimes S_\eta(B)v_\eta
\]
and
\[
 (Z,H,B)\longmapsto
 \left(H^a\operatorname{per}_m\left(\sum_{i=1}^rZ_i x_i\right)\right)^{\otimes s}
                  \otimes S_\eta(B)v_\eta.
\]
The parameter $H$ is an independent linear form and $B$ an
independent endomorphism of $V^*$. The first factors have
parameter degree $ns$, and the second has degree $e$.
If $e=0$ it is constant and may be omitted. These are parameter
degrees, distinct from the action degree $L$.
G-10's base-locus-aware degree argument bounds the two
projective invariant-Veronese image degrees by
\[
 (pE)^{u-1}\quad\hbox{and}\quad(pE)^{v-1},
\]
respectively.

Write $j=Ek+t$, $0\le t<E$. For either domain $A$, multiplication
by the common nonzero invariants $h_{E+t}$ and $h_{2E-t}$ gives
\begin{equation}\label{G23:sandwich}
 \dim A_{E(k-1)}\le\dim A_j\le\dim A_{E(k+2)}\qquad(k\ge1).
\end{equation}
If $v>u$, linear Noether normalization of $A_P^{(E)}$ and
the Chardin bound on $A_D^{(E)}$ yield
\[
 \dim(A_P)_{E(k-1)}\ge\binom{k+u-1}{u},\qquad
 \dim(A_D)_{E(k+2)}\le(pE)^{u-1}\binom{k+u+1}{u-1}.
\]
Their binomial ratio is
\[
 \frac{k(k+1)(k+2)}{u(k+u)(k+u+1)}
                 \ge\frac{k}{4u}\qquad(k\ge u+1).
\]
The asserted positive threshold follows from
\eqref{G23:sandwich}. If $u>v$, exchange the domains and
replace $u$ by the actual dimension $v$ in this argument.
This gives the asserted reverse threshold without an upper
bound substituted for the permanent dimension.

For a standard graded domain of dimension $q\ge2$, the
leading Hilbert term is
$\deg(\operatorname{Proj} A^{(E)})k^{q-1}/(q-1)!$,
with error $O(k^{q-2})$. Applying \eqref{G23:sandwich}
shows, through all integers, that
\[
 \dim A_j=c_A j^{q-1}+O(j^{q-2}),\qquad c_A>0.
\]
Taking the ratio for $A_P,A_D$ proves \eqref{G23:ratio}.

\paragraph{3. The integer dimensions can never be equal.}
If $r$ is even, $\Delta$ is odd. Suppose $r$ is odd and
$\Delta=0$. Reduction modulo two forces $m,n$ to have
the same parity. Whether both are even or both are odd,
reduction modulo four then gives
\[
 0=\Delta\equiv r-1\pmod4,
\]
so $r\equiv1\pmod4$. Reduction modulo $r-2$ gives
$2m^2-2m+1\equiv0$, and hence
\[
 (2m-1)^2\equiv-1\pmod{r-2}.
\]
But $r-2\equiv3\pmod4$, so one of its prime divisors
$p$ is congruent to three modulo four. A square root of
$-1$ modulo $p$ is impossible: its nonzero residue would
have $(p-1)$st power both $1$, by Fermat's theorem, and
$(-1)^{(p-1)/2}=-1$. This contradiction proves
$\Delta\ne0$. The precise integer boundary
\eqref{G23:boundary} follows immediately.
\end{proof}

\paragraph{Rectangular consequence and limits.}
For a rectangular direction, the strict condition is
$n<rm/(r-1)$. Throughout that range the exponent is positive:
\[
 \Delta>
 \frac{r m^2}{(r-1)^2}-2m+r-1
 =\frac r{(r-1)^2}\left(m-\frac{(r-1)^2}{r}\right)^2
                                      +\frac{r-1}{r}>0.
\]
Thus every such rectangle gives a gap at every sufficiently
large admissible coordinate degree. For a fixed $m,n,r$,
one may take $s=r/\gcd(n,r)$ and
$\nu=(ns/r)^r$, so all admissible degrees are $sj$.
The theorem concerns fixed strict directions with $r\le m$.
It does not decide the face $\nu_1=as$, the range $r>m$,
every earlier finite degree, or a uniform polynomial-time
positive rule when the parameters grow. Its bounds and
non-occurrence-obstruction conclusion retain these limits.



\begin{theorem}[G-24: The exact eventual rectangular range through $m$ rows]\label{G-24}
Let $n>m\ge17$ and $4\le r\le m$. At every sufficiently
large admissible coordinate degree $d$ (meaning $r\mid nd$),
the rectangular multiplicity gap
\[
 m_{(nd/r)^r}(\ell^{n-m}\operatorname{per}_m)
       >m_{(nd/r)^r}(\det_n)
\]
holds if and only if
\begin{equation}\label{G24:range}
 m<n<\frac{rm}{r-1}.
\end{equation}
Inside this range, G-23 supplies the explicit threshold.
Above the boundary, the permanent rectangular multiplicity
is zero in every positive admissible degree and the determinant
multiplicity is positive in every sufficiently large admissible
degree. At the equality boundary, the determinant multiplicity
is strictly larger in every sufficiently large degree.
No assertion excluding all earlier boundary gaps is made for
$r\ge5$; G-12 supplies that stronger weak inequality for $r=4$.

For an explicit equality-boundary threshold, write
$a=n-m$, so that $n=ra$, $m=(r-1)a$, and define
\[
 \begin{split}
 u&=(r-2)n^2-r^2+3,\\
 v&=\begin{cases}
       \binom{m+2}{2}-8,&r=4,\\
       (r-1)m^2-2m-r^2+2r+2,&r\ge5.
      \end{cases}
 \end{split}
\]
Let
\[
 b=\binom{n+r-1}{r-1},\qquad
 \beta=\left\lceil\frac{3r^2b}{8}\right\rceil n^{2(r^2-1)},
 \qquad T=\binom{b+\beta}{b}-1,\qquad E=rT(\beta!).
\]
Then $u>v\ge2$, and determinant dominance holds for every
\begin{equation}\label{G24:threshold}
 d\ge E\bigl(4v(nE)^{v-1}+1\bigr).
\end{equation}
\end{theorem}

\begin{proof}
The strict interior is the rectangular consequence of G-23.
For the other cases, set $a=n-m$ and use $r$ variables
$x_0,\ldots,x_{r-1}$. A nonzero padding form can be taken
to $x_0$ by a special-linear change of variables. The
one-parameter subgroup
\[
 \lambda(t)=\operatorname{diag}(t^{r-1},t^{-1},\ldots,t^{-1})
\]
gives weight $(r-1)a-m+rj=ra-n+rj$ to a monomial of
$x_0^a q$ coming from the $x_0^j$ part of $q$.
If $(r-1)n>rm$, all these weights are positive. The form
therefore tends to zero and every positive-degree invariant
vanishes. This proves the zero assertion above the boundary.
For eventual determinant occurrence, take the common free
point of R-28 with smaller size $n-1$ and padding one.
For the minimal admissible direction use
$s_0=r/\gcd(n,r)$ and $\nu=(ns_0/r)^r$.
Its strict condition is $ns_0/r>s_0$, since $n>r$.
R-28 consequently gives positive determinant multiplicity
at every sufficiently large integer stretch, that is, every
sufficiently large admissible coordinate degree.

\paragraph{1. Invariant functions inject into one fewer variables at equality.}
Now suppose $(r-1)n=rm$. Then $n=ra$ and $m=(r-1)a$.
The same limit is
\begin{equation}\label{G24:limit}
 \lim_{t\to0}\lambda(t)x_0^a q
                =x_0^a q(0,x_1,\ldots,x_{r-1}).
\end{equation}
Put $q_0=r-1$ for the number of variables after restriction,
and write
\[
 A_P=\mathbb C[\{L^a q:q\in\mathcal P_{m,r}\}^{\rm cl}]^{\mathrm{SL}_r},
 \qquad C=\mathbb C[\mathcal P_{m,q_0}]^{\mathrm{SL}_{q_0}}.
\]
Restriction along $q\mapsto x_0^a q$ gives a graded algebra
homomorphism
\begin{equation}\label{G24:injection}
 \rho:A_P\longrightarrow C.
\end{equation}
It is well defined because the indicated padded forms belong
to the original cone, and the subgroup
$\operatorname{diag}(1,\mathrm{SL}_{q_0})$ fixes $x_0$.
Its grading is unchanged: coefficients of $x_0^a q$ depend
linearly on coefficients of $q$.

To prove injectivity, let an invariant function $h$ restrict
to zero. For any original padded form, take its nonzero
padding line to $x_0$. The function $h$ is constant along
$\lambda(t)$ and hence, by regularity, at its limit
\eqref{G24:limit}. Restricting a permanent pencil to
$x_0=0$ puts the residual form in $\mathcal P_{m,q_0}$,
including its closure points. Thus the value at the limit
is zero. The dense padded image has zero value everywhere,
so $h=0$ on its closure. The same statement holds in degree
zero. This proves \eqref{G24:injection} without claiming
that every lower-rank invariant lifts.

Let $A_D=\mathbb C[\mathcal D_{n,r}]^{\mathrm{SL}_r}$.
For $r\ge5$ put $B=C$. For $r=4$ use instead
\[
 B=\mathbb C[\operatorname{Sym}^m\mathbb C^3]^{\mathrm{SL}_3}.
\]
In this latter case exactness of invariants gives a graded
surjection $B\to C$, not an asserted injection into the
full ternary space. Consequently in every case
\begin{equation}\label{G24:hilbertupper}
 \dim(A_P)_d\le\dim B_d.
\end{equation}

\paragraph{2. The two comparison rings have known dimensions and no period.}
The equality hypotheses and $r\le m$ imply $a\ge2$.
Apply R-28's common-free-point statement, which allows
row count three, first with degree $n$, size $n-1$, row
count $r$ and rectangular direction $(a^r)$ in degree one.
Apply it a second time with degree $m$, size $m-1$, row
count $q_0$ and direction $(a^{q_0})$.
Both applications satisfy the strict inequality $a>1$ and
the size condition. They supply a stable point with no
nontrivial scalar-orbit symmetry in $A_D$ and in the
permanent cone underlying $C$. The latter is also such a
point in the full ambient ternary representation when
$r=4$.

It follows that the stable quotients have dimensions
\[
 \dim A_D=(r-2)n^2+2-(r^2-1)=u,
 \qquad \dim B=v.
\]
For $r\ge5$, use G-21 in $q_0\ge4$ variables; for $r=4$
the whole ternary coefficient space has dimension
$\binom{m+2}{2}$ and the stable orbit dimension is eight.
The absence of scalar-orbit symmetry and invariant
separation give gcd one for the nonzero invariant degrees,
just as in R-28. We use this separately for the two rings;
no common invariant in their different ambient spaces is
being assumed.

The dimension inequality is strict. For $r\ge5$ it reads
\[
 u-v=\frac{r^2-3r+1}{(r-1)^2}m^2+2m-2r+1>0,
\]
since $m\ge r$ and $r^2-3r+1>0$.
For $r=4$ it is
\[
 u-v=\frac{55m^2-27m-108}{18}>0\qquad(m\ge17).
\]
The displayed dimensions are at least two in the stated range.

\paragraph{3. An explicit comparison at every integer degree.}
Derksen's general bound, specialized as in G-23 with $s=1$
and $\eta=0$, gives the displayed $\beta$ for the degree-$n$
ambient representation in $r$ variables. It also bounds the
generators for the degree-$m$ representation in $q_0<r$
variables: its vector-space dimension, action degree and
group-dimension exponent are all no larger. Exactness of
invariants passes these bounds to $A_D$ and to $B$.
The chosen $T$ bounds both generator counts. Hence the
$E$th Veroneses of both rings are standard graded by the
same generator-grouping argument as G-23.

The gcd-one argument and the bound $\beta$ give nonzero
elements in every degree at least $\beta^2$ in each ring
separately. For $d=Ek+t$, $0\le t<E$, multiplication by
elements of degrees $E+t$ and $2E-t$ therefore gives the
individual sandwiches
\[
 \dim A_{E(k-1)}\le\dim A_d\le\dim A_{E(k+2)},
 \qquad A=A_D\text{ or }B.
\]
The projective degree of $B^{(E)}$ is at most
$(nE)^{v-1}$. If $r\ge5$, compose its invariant-coordinate
map with the permanent-pencil parametrization of degree
$m<n$ and use G-10. If $r=4$, use the identity parametrization
of the full ternary coefficient space, of degree one, and
the same larger bound. No dominance of the ternary
permanent parametrization is needed here.

Noether normalization on $A_D^{(E)}$, whose dimension is
at least $v+1$, and the degree/Hilbert bound on $B^{(E)}$
now give
\[
 \dim(A_D)_{E(k-1)}\ge\binom{k+v-1}{v},\qquad
 \dim B_{E(k+2)}\le(nE)^{v-1}\binom{k+v+1}{v-1}.
\]
The same exact ratio as in G-23 is at least $k/(4v)$ for
$k\ge v+1$. For $k\ge4v(nE)^{v-1}+1$, this proves
$\dim(A_D)_d>\dim B_d\ge\dim(A_P)_d$ in every residue.
All degrees are admissible at this boundary because $r\mid n$.
The invariant multiplicities are the claimed rectangles,
by the same single-label extraction and inheritance as
G-11/G-22. This proves \eqref{G24:threshold} and the theorem.
\end{proof}

\paragraph{Exact scope.}
The equality-boundary argument is an injection of actual
invariant coordinate rings followed by a Hilbert-dimension
comparison. It neither identifies those rings nor replaces
orbit-closure multiplicities by ambient plethysm coefficients.
The theorem closes the eventual rectangular boundary through
$m$ rows in the stated size range. It leaves nonrectangular
boundary faces, earlier finite-degree gaps at $r\ge5$,
more than $m$ rows, and efficient positive extraction open.



\begin{theorem}[G-25: Exact strict-direction comparison through $2m-2$ rows]\label{G-25}
Every assertion and every explicit bound of G-23 remains valid when its
row range $4\le r\le m$ is replaced by
\[
             4\le r\le2m-2,
\]
with $n>m\ge17$ and the same strict conditions
$\nu_r>0$, $\nu_1>(n-m)s$, $\nu\vdash ns$.
In particular the exact eventual exponent is still
\[
 \Delta=rm^2-(r-2)n^2-2m+r-1\ne0,
\]
and both multiplicities are positive at every integer stretch
$j\ge\beta^2$, with the same definitions of $b,L,p,\beta,T,E,u,v$.
The effective thresholds for strict comparison, and the exact
floor boundary $N_{m,r}$, are those of G-23.

In the newly added range $m<r\le2m-2$, the exponent is positive
if and only if
\[
                   r=n=m+1.
\]
For these parameters it equals one. Thus, for every $m\ge17$,
the explicit nonrectangular direction
\[
 s=2,\qquad \nu=(3,2^{m-1},1),\qquad n=r=m+1
\]
has an eventually positive genuine-padding multiplicity gap in
every sufficiently large coordinate degree $2j$, and its
permanent-to-determinant multiplicity ratio is
$\kappa_\nu j(1+O(j^{-1}))$, with $\kappa_\nu>0$.
Every other pair of integer parameters in this newly added row
range has eventual determinant dominance in every fixed strict
direction.
\end{theorem}

\begin{proof}
R-29 gives both the common stable tensor and a common stable tensor
with no nontrivial scalar multiple in its orbit throughout the
enlarged range. Its enhanced-coordinate extraction agrees exactly
with that in G-22. The dimension formulas used by G-23 already
hold here: G-21 proves the permanent formula for every $r\ge4$
when $m\ge17$, and R-28 proves the determinant formula and the
padding dimension for every $r\ge4$ in the present size range.
Consequently the invariant-domain dimensions are precisely the
same $u,v$ displayed in G-23.

The only numerical estimate in that proof that used $r\le m$
beyond the stable-point input was $u,v\ge2$. We supply the
replacement explicitly. Since $r-2\ge r/2$, $r\le2m-2$,
$n\ge m+1$ and $r_\eta\ge0$,
\[
 \begin{aligned}
 u&=(r-2)n^2-r^2+r_\eta+3
       \ge\frac r2(m^2-2m+5)+3>2,\\
 v&=r(m^2-r+1)-2m+r_\eta+2
       \ge r(m^2-2m+3)-2m+2>2.
 \end{aligned}
\]
The last inequality follows already from $r\ge4$ and $m\ge17$.

All subsequent steps of G-23 use only these exact dimensions,
the ambient representation in rank $r$, and the two properties
of the common tensor supplied by R-29. Derksen's action-degree
and generator bounds are unchanged. The same nonzero generator
degrees have gcd one by closed-orbit separation, so the common
invariants $h_j$ exist for every $j\ge\beta^2$. The same
generator grouping gives the standard $E$th Veroneses, and the
same degree-$p$ parameter maps give their image-degree bounds.
Their every-residue Hilbert sandwiches, the Noether lower bound,
and the Chardin upper bound therefore prove the stated effective
thresholds and ratio with these actual values of $u,v$.
The parity and nonsquare argument excluding $\Delta=0$ in G-23
was valid for every integer $r\ge4$, with no upper row bound.
This proves the full extension, including all its constants.

For the additional sign classification, fix $m<r$.
The exponent is strictly decreasing as an integer function of
$n\ge m+1$, and direct substitution gives
\[
       \Delta\big|_{n=m+1}=2m(m-r+1)+1.
\]
If $r\ge m+2$, this is at most $1-2m<0$.
If $r=m+1$, its value at $n=m+1$ is one, while increasing $n$
to $m+2$ subtracts $(m-1)(2m+3)>1$. This proves the asserted
unique positive parameter case. Finally the displayed partition
has $m+1$ positive parts and total size $2m+2=ns$; its first
part is $3>2=(n-m)s$. It therefore meets the strict hypotheses,
and substituting $\Delta=1$ proves its claimed ratio and gap.
\end{proof}

\paragraph{Scope.}
The new positive example specifies a total multiplicity inequality,
not merely one nonzero equation copy. It is an eventual statement
with the explicit G-23 threshold and does not assert the comparison
at every earlier degree. No rectangle satisfies the strict
condition when $r>m$ and $n>m$ are integers: then
$rm/(r-1)\le m+1$. The rectangular boundary result G-24 keeps
its stated range $r\le m$. General nonrectangular equality
faces, more than $2m-2$ rows, changing directions, and
polynomial-resource coefficient extraction remain separate gates.


\begin{corollary}[G-25, endpoint extension from R-30]\label{G25:endpoints}
All assertions of G-25, including every explicit G-23 constant,
hold throughout $4\le r\le2m$. The unique positive case among
the added rows $r>m$ is still $r=n=m+1$.
\end{corollary}

\begin{proof}
R-30 supplies the common stable no-scaling points at
$r=2m-1$ and $r=2m$. The exact determinant and permanent
dimensions in R-28 and G-21 hold at both endpoints. For
completeness, the required lower bounds on their enhanced
invariant dimensions are, using $r\le2m$ and $n\ge m+1$,
\[
 \begin{aligned}
 u&\ge\frac r2(m-1)^2+3>2,\\
 v&\ge r(m-1)^2-2m+2>2,
 \end{aligned}
\]
because $r-2\ge r/2$ and $m\ge17$.
The exact proof of G-25 then applies with the unchanged
rank-$r$ representation, action degree, parameter degree,
generator bound, common degree semigroup and standard Veronese.
Its arithmetic sign classification used no upper row bound.
This proves the assertion at both new endpoints.
\end{proof}


\begin{theorem}[G-26: The complete first-row padding boundary through $m$ rows]\label{G-26}
Let $a=n-m>0$. For every coordinate degree $d\ge1$ and partition
$\lambda\vdash nd$ of length at most $r$, the padded permanent
multiplicity satisfies
\begin{equation}\label{G26:pieri_bound}
 \begin{cases}
 m_\lambda(\ell^a\operatorname{per}_m)=0,&\lambda_1<ad,\\[2pt]
 m_\lambda(\ell^a\operatorname{per}_m)
   \le m_{(\lambda_2,\ldots,\lambda_r)}(\operatorname{per}_m),
       &\lambda_1=ad.
 \end{cases}
\end{equation}
The multiplicity on the second line can be computed on the
permanent image in $r-1$ variables. This bound holds in every
finite degree over $\mathbb C$; no lifting or equality of these
multiplicity spaces is asserted.

Now suppose $n>m\ge17$ and $4\le r\le m$. Fix
$\nu\vdash ns$, $s\ge1$, with exactly $r$ positive parts.
If $\nu_1=(n-m)s$, the determinant multiplicity is strictly
larger than the padded permanent multiplicity at every
sufficiently large integer stretch $j$. The same conclusion
holds when $\nu_1<(n-m)s$, in which case the padded multiplicity
is zero at every positive stretch.

Consequently, in this row range every fixed full-row direction
has an exact eventual comparison: the padded permanent
multiplicity is eventually larger if and only if
\[
 \nu_1>(n-m)s
 \quad\hbox{and}\quad
 rm^2-(r-2)n^2-2m+r-1>0.
\]
In every other case the determinant multiplicity is eventually
strictly larger. In the strict-padding case the bounds are
those of G-23; an explicit equality-face bound follows below.
\end{theorem}

\paragraph{An explicit threshold on the equality face.}
Assume $\nu_1=as$, put $q=r-1$, and define
\[
 \begin{aligned}
 \mu&=(\nu_2,\ldots,\nu_r),&
 \eta&=\nu-\nu_r\mathbf1_r,&
 \theta&=\mu-\nu_r\mathbf1_q,\\
 e&=|\eta|,&e'&=|\theta|,&
 r_\eta&=\#\{i<j:\eta_i>\eta_j\},
 \end{aligned}
\]
with $r_\theta$ defined in the same way. Set
\[
 \begin{aligned}
 u&=(r-2)n^2-r^2+3+r_\eta,\\
 v&=\begin{cases}
       \binom{m+2}{2}-8+r_\theta,&r=4,\\
       (r-1)m^2-2m-r^2+2r+2+r_\theta,&r\ge5,
      \end{cases}\\
 b_D&=\binom{\binom{n+r-1}{r-1}+s-1}{s}\dim S_\eta\mathbb C^r,\\
 b_B&=\binom{\binom{m+q-1}{q-1}+s-1}{s}\dim S_\theta\mathbb C^q,\\
 b&=\max\{b_D,b_B\},&
 L&=ns+(r-1)e,&p&=ns+e,\\
 \beta&=\left\lceil\frac{3r^2b}{8}\right\rceil L^{2(r^2-1)},&
 T&=\binom{b+\beta}{b}-1,& E&=rT\beta!.
 \end{aligned}
\]
Then $u>v\ge2$, and strict determinant dominance holds for every
\begin{equation}\label{G26:threshold}
                  j\ge E\bigl(4v(pE)^{v-1}+1\bigr),
\end{equation}
in original coordinate degree $sj$. All Schur dimensions here
are given by the Weyl products used in G-23.

\paragraph{Background used.}
The Pieri formula is cited from Bump and Gamburd (2006),
\S2.1, equation with source label \texttt{e:pieri}, supplied
\texttt{acp20.tex}, lines 404--422. That source identifies it
with Macdonald's equation I~(5.16), p.~73, and Bump's
Theorem~42.4; we inspected the supplied Bump--Gamburd statement.
All other inputs are the exact invariant-extraction, degree,
and Hilbert bounds already used and located in G-22--G-24.

\begin{proof}
We first prove the finite-degree inequality and its change of
ambient dimension. We then compare two different enhanced
invariant algebras, with their degree semigroups treated separately.

\paragraph{1. Pieri removes exactly the first row.}
Let $V=\mathbb C^r$ and let $Y$ be the padded permanent image.
The parameter morphism
\[
 V\times\mathcal P_{m,r}\longrightarrow Y,
                    \qquad(\ell,F)\longmapsto\ell^a F
\]
has dense image by the definition of $Y$. Pullback of degree-$d$
coordinate functions is therefore an injective equivariant map
\begin{equation}\label{G26:parameter_injection}
 \mathbb C[Y]_d\lhook\joinrel\longrightarrow
       \operatorname{Sym}^{ad}V^*\otimes\mathbb C[\mathcal P_{m,r}]_d.
\end{equation}
The two parameter degrees are $(ad,d)$; they are not both $d$.
Use the contragredient convention for coordinate representations
fixed in the introduction. By complete reducibility over
$\mathbb C$ and Pieri, a copy of $S_\lambda V^*$ in the right
side comes from $S_\alpha V^*$ precisely when $\lambda/\alpha$
is a horizontal strip of size $ad$, with multiplicity one.

There are only $\lambda_1$ columns. If $\lambda_1<ad$, no such
strip exists. If $\lambda_1=ad$, the strip must contain exactly
one box in every column. Thus
\[
 \alpha'_k=\lambda'_k-1\quad(1\le k\le\lambda_1),
 \qquad
                 \alpha=(\lambda_2,\ldots,\lambda_r).
\]
This is the unique possibility, and its size is $md$.
Taking multiplicities in \eqref{G26:parameter_injection} gives
the asserted zero and upper bound.

For completeness, the multiplicity of this tail partition
equals that on $\mathcal P_{m,r-1}$. A highest-weight coordinate
polynomial with last weight zero uses only coefficient variables
of monomials not involving the last original variable; all these
coefficient weights are nonnegative in the polynomial convention.
Its vanishing on the $r$-variable permanent image is equivalent
to vanishing on the $(r-1)$-variable image: restriction of every
$r$-variable matrix pencil gives a pencil with one fewer variable,
and every such smaller pencil occurs by setting the unused
coefficient matrix to zero. The highest-weight coefficient
subalgebras and their ideal intersections consequently agree.
This is precisely the inheritance argument stated in the
introduction, now applied to the tail shape of length at most
$r-1$. It proves the last assertion of
\eqref{G26:pieri_bound} without replacing an actual coordinate
ring by its ambient plethysm ring.

\paragraph{2. The equality face gives two stable comparison algebras.}
Now assume the displayed equality-face hypotheses. Because
$\nu_1=as$ is the largest of the $r$ parts,
\[
 ns\le ras,\qquad
 n\ge\frac{rm}{r-1},\qquad
 a\ge\frac{m}{r-1}>1.
\]
Thus the integer $a$ is at least two. The tail $\mu$ has size
$ms$, exactly $q=r-1$ positive parts, and
\[
                    \mu_1\ge\frac{ms}{r-1}>s.
\]
Let $A$ be the enhanced invariant coordinate algebra for
the determinant cone $\mathcal D_{n,r}$, outer stretch $s$,
and flag weight $\eta$, as in G-22. Then
\[
                       \dim A_j=m_{j\nu}(\det_n).
\]
For $r\ge5$, let $B$ be the corresponding enhanced invariant
algebra for $\mathcal P_{m,q}$ and flag weight $\theta$.
For $r=4$, take instead the enhanced algebra for the full
ternary degree-$m$ coefficient space. In that case exactness
of invariants gives a degreewise surjection from $B$ to the
enhanced permanent algebra. In either case, part~1 gives
\begin{equation}\label{G26:degreewise_bound}
             m_{j\nu}(\ell^a\operatorname{per}_m)\le\dim B_j.
\end{equation}

Both $A$ and $B$ have a stable point with the no-scaling property.
For $A$, apply R-28 with degree $n$, matrix size $n-1$, row
count $r$, and the direction $\nu$: its strict condition is
$\nu_1=as>s$. The hypotheses $r\le n-1$ hold since
$r\le m<n$. For $B$, use degree $m$, matrix size $m-1$,
row count $q=r-1$, and the direction $\mu$. Its strict
condition is $\mu_1>s$, and $q\le m-1$.
The smaller padded permanent embeds in $\mathcal P_{m,q}$
by block diagonal padding with one linear entry. At $q=3$
the same point also lies in the full ternary coefficient space.
The common-point clause of R-28 applies in row count three
and does not require size at least seventeen for that clause.
The enhanced representations, and their scalar characters,
are exactly those just specified.

It follows from G-22's extraction and the stable quotient
dimension formula that
\[
                        \dim A=u,\qquad\dim B=v.
\]
Here R-28 supplies the determinant image dimension, G-21
supplies the permanent dimension for $q\ge4$, and the ternary
ambient dimension is $\binom{m+2}{2}$ for $q=3$.
No common invariant between two different groups is being
asserted. The no-scaling property gives gcd one for the
nonzero degrees in each ring separately, by the closed-orbit
separation argument already proved in R-28.

\paragraph{3. The determinant quotient has strictly larger dimension.}
Subtracting a constant from every part does not change flag
dimension. Deleting the first part removes exactly its strict
comparisons with the others, so
\[
 r_\eta-r_\theta
       =\#\{j>1:\nu_1>\nu_j\}\ge0.
\]
For $r\ge5$ this gives
\[
 \begin{aligned}
 u-v&=(r-2)n^2-(r-1)m^2+2m-2r+1+r_\eta-r_\theta\\
 &\ge\frac{r^2-3r+1}{(r-1)^2}m^2+2m-2r+1>0,
 \end{aligned}
\]
since $m\ge r$ and $r^2-3r+1>0$. For $r=4$,
\[
 u-v=2n^2-\binom{m+2}{2}-5+r_\eta-r_\theta
       \ge\frac{55m^2-27m-108}{18}>0.
\]
The explicit formulas give $v\ge2$; for instance at $r\ge5$
one has $v\ge3m^2-2m+2r+2>2$ by $r\le m$, and the ternary
case is immediate. Thus $u\ge v+1$.

\paragraph{4. Effective comparison in every original stretch.}
The two ambient representation dimensions are $b_D,b_B$.
Their action degrees are at most $L$: the smaller one has
degree $ms+(q-1)e'$, and
$e-e'=\nu_1-\nu_r\ge0$. Derksen's bounds used in G-23
therefore give the displayed common bound $\beta$ for the
two invariant generator degrees, even though the groups and
representations differ. Their numbers of generators are each
at most $T$. The identical generator grouping proves that
their $E$th Veroneses are standard graded. Each no-scaling
point gives nonzero elements in every degree at least
$\beta^2$, and $E\ge\beta^2$.

The determinant enhanced cone has its degree-$p$ dense
homogeneous parameter map from G-23. For $r\ge5$, the lower
permanent map has parameter degree $ms+e'\le p$. For $r=4$,
the full coefficient-space parameter has degree $s+e'\le p$.
The base-locus-aware image-degree bound thus gives
\[
                  \deg\operatorname{Proj} B^{(E)}\le(pE)^{v-1}.
\]
Write $j=Ek+t$, $0\le t<E$. In each ring separately, multiply
by its nonzero elements of degrees $E+t$ and $2E-t$ to obtain
the same two-sided Hilbert sandwich as G-23. Linear Noether
normalization and the Chardin upper bound yield
\[
 \dim A_j\ge\binom{k+v-1}{v},\qquad
 \dim B_j\le(pE)^{v-1}\binom{k+v+1}{v-1}.
\]
Their binomial ratio is at least $k/(4v)$ for $k\ge v+1$.
Condition \eqref{G26:threshold} implies
$k>4v(pE)^{v-1}$, so the first bound is strictly larger.
Together with \eqref{G26:degreewise_bound}, this proves the
equality-face conclusion at every stated integer stretch.

If instead $\nu_1<as$, part~1 makes every padded multiplicity
zero. The determinant multiplicity is eventually positive:
$\nu_1\ge ns/r>s$ since $r\le m<n$, so the first R-28
application in part~2 still applies. Its generator-semigroup
bound gives occurrence for every sufficiently large $j$.
Finally combine these two cases with G-23 in the strict
half-space and its proof that $\Delta$ never vanishes. This
proves the exact eventual classification in the theorem.
\end{proof}

\paragraph{Scope.}
The Pieri upper bound is an all-degree reduction and holds
beyond the row range of the eventual comparison. The strict
boundary comparison proved here uses $4\le r\le m$ and
$m\ge17$. It does not assert a positive-degree lifting from
the lower permanent algebra, eliminate earlier boundary gaps
when $r\ge5$, or settle boundary faces with more than $m$
rows. It is independent of the conditional analytic inputs.


\begin{theorem}[G-27: Nonrectangular boundary comparison through $2m-1$ rows]
\label{G-27}
Let $n>m\ge17$, $m<r\le2m-1$, and let
$\nu\vdash ns$, $s\ge1$, have exactly $r$ positive parts
and be nonrectangular. Put $a=n-m$. If $\nu_1\le as$,
the determinant multiplicity is strictly larger than the
padded permanent multiplicity at every sufficiently large
integer stretch. On the equality face the explicit threshold
below applies. Below the face the padded multiplicity is zero
at every positive stretch, and the determinant occurs at every
stretch beyond its explicit invariant-generator bound squared.

Together with G-26 and G-25, this classifies the eventual sign
in every fixed nonrectangular full-row direction with
$4\le r\le2m-1$. A positive permanent-minus-determinant gap
occurs eventually exactly when
\[
 \nu_1>(n-m)s,\qquad
 \Delta=rm^2-(r-2)n^2-2m+r-1>0.
\]
In every other case the determinant eventually dominates.
The strict-direction assertion already extends through $2m$
rows by G-25; the additional boundary assertion here stops
at $2m-1$ rows and excludes rectangular directions beyond
the range of G-26.
\end{theorem}

\paragraph{Explicit equality-face constants.}
Assume $\nu_1=as$. Put
\[
 q=r-1,\quad \mu=(\nu_2,\ldots,\nu_r),\quad
 \eta=\nu-\nu_r\mathbf1_r,\quad
 \theta=\mu-\nu_r\mathbf1_q,
 \quad e=|\eta|,\quad e'=|\theta|.
\]
As before, $r_\eta$ and $r_\theta$ count strict comparisons
between ordered parts. Define
\[
 \begin{aligned}
 u&=(r-2)n^2-r^2+3+r_\eta,\\
 v&=(r-1)m^2-2m-r^2+2r+2+r_\theta,\\
 b_D&=\binom{\binom{n+r-1}{r-1}+s-1}{s}
                    \dim S_\eta\mathbb C^r,\\
 b_B&=\binom{\binom{m+q-1}{q-1}+s-1}{s}
                    \dim S_\theta\mathbb C^q,\\
 b&=\max\{b_D,b_B\},\qquad
 L=ns+(r-1)e,\qquad p=ns+e,\\
 \beta&=\left\lceil\frac{3r^2b}{8}\right\rceil
                      L^{2(r^2-1)},\\
 T&=\binom{b+\beta}{b}-1,\qquad E=rT\beta!.
 \end{aligned}
\]
Then $u>v\ge2$, and strict determinant dominance holds for
every
\begin{equation}\label{G27:threshold}
                    j\ge E\bigl(4v(pE)^{v-1}+1\bigr)
\end{equation}
in coordinate degree $sj$. These are the same form of bounds
as G-26, but the lower invariant algebra is now allowed to
have a nontrivial degree period.

\begin{proof}
We compare the two enhanced invariant algebras and retain the
degree period only on the side used for an upper bound.
\paragraph{1. Separate stable algebras, with period one needed only above.}
On the equality face, the all-degree Pieri bound G-26 gives
\[
 m_{j\nu}(\ell^a\operatorname{per}_m)
       \le m_{j\mu}(\mathcal P_{m,q}).
\]
Let $A$ be the enhanced invariant algebra for
$\mathcal D_{n,r}$, outer degree $s$, and auxiliary weight
$\eta$, and let $B$ be that for $\mathcal P_{m,q}$ with
auxiliary weight $\theta$. Exact extraction yields
\[
 \dim A_j=m_{j\nu}(\det_n),\qquad
 \dim B_j=m_{j\mu}(\mathcal P_{m,q}).
\]
The two applications of R-32 are valid because
\[
 4\le r\le2m-1<2n-2,\qquad
 4\le q=r-1\le2m-2.
\]
Thus both enhanced cones have stable points with finite
projective stabilizers. Since $\nu$ is nonrectangular, the
first cone additionally has a stable point with no scalar
period. No such assertion is needed for $B$: the tail $\mu$
may be rectangular.

The exact image dimensions from R-28 and G-21 and the stable
quotient formula give $\dim A=u$ and $\dim B=v$.
The degree semigroup of $A$ has gcd one. The corresponding
gcd for $B$ is left unrestricted. This distinction is the
reason that no assertion about generic freeness of a
rectangular two-product tensor is required.

\paragraph{2. The dimension inequality persists beyond $m$ rows.}
Deleting the first part removes
\[
 r_\eta-r_\theta=\#\{j>1:\nu_1>\nu_j\}\ge0
\]
strict flag comparisons. Since $n\ge m+1$,
\[
 \begin{aligned}
 u-v
  &=(r-2)n^2-(r-1)m^2+2m-2r+1+r_\eta-r_\theta\\
  &\ge(2m-1)r-m^2-2m-1\\
  &\ge m^2-m-2>0.
 \end{aligned}
\]
For the last inequality use $r\ge m+1$. Moreover, writing
$q=r-1$, with $m\le q\le2m-2$, gives
\[
 v=q(m^2-q)-2m+3+r_\theta
        \ge m^3-2m^2+3>2.
\]
In particular $u\ge v+1$.

\paragraph{3. An upper Hilbert bound does not require degree period one.}
The degree and dimension estimates in G-23/G-26 give
homogeneous generator degrees at most $\beta$ in both rings,
at most $T$ generators, and standard $E$th Veroneses.
They also give $E\ge\beta^2$. The determinant no-scaling
point implies $A_j\ne0$ for every $j\ge\beta^2$.

For the lower algebra $B$, let $g$ be the gcd of the positive
degrees of a set of nonzero homogeneous generators. This
set is nonempty because $v\ge2$. Its degrees are at most
$\beta$. Divide them by $g$ and apply the elementary
residue-semigroup bound of G-18. It follows that every
multiple of $g$ at least $\beta^2$ belongs to their semigroup.
Since $B$ is a domain, a product of nonzero generators is
nonzero; hence $B_h\ne0$ at every such degree $h$.
Conversely $B_j\ne0$ implies $g\mid j$, by homogeneity of
the generators. This argument makes no saturation claim in
degrees not divisible by $g$.

Write $j=Ek+t$, $0\le t<E$. Since $g\le\beta$ and
$\beta!\mid E$, one has $g\mid E$. If $B_j\ne0$, then
$g\mid t$, and $2E-t\ge E\ge\beta^2$ is a multiple of
$g$. Multiplication by a nonzero element of $B_{2E-t}$
therefore embeds $B_j$ in $B_{E(k+2)}$. If $B_j=0$, the
same upper inequality holds without a multiplication map.
For $A$, multiplication by any nonzero element of
$A_{E+t}$ embeds $A_{E(k-1)}$ in $A_j$. Thus, for $k\ge1$,
\begin{equation}\label{G27:one_sided_sandwich}
 \dim A^{(E)}_{k-1}\le\dim A_j,\qquad
 \dim B_j\le\dim B^{(E)}_{k+2}.
\end{equation}
Only the lower bound in the first ring requires period one.

The lower permanent enhanced cone has a dense homogeneous
parameter map of degree $ms+e'\le p$. The image-degree
estimate used in G-23 therefore gives
\[
             \deg\operatorname{Proj}B^{(E)}\le(pE)^{v-1}.
\]
Linear Noether normalization for the standard domain
$A^{(E)}$ and the same Chardin bound for $B^{(E)}$ now give
\[
 \dim A_j\ge\binom{k+v-1}{v},\qquad
 \dim B_j\le(pE)^{v-1}\binom{k+v+1}{v-1}.
\]
As in G-26, their binomial ratio is at least $k/(4v)$ for
$k\ge v+1$, and \eqref{G27:threshold} makes the first
quantity strictly larger than the second. The Pieri upper
bound proves the desired strict comparison at every such
integer stretch, including all possible lower-ring residues.

\paragraph{4. The remaining two half-spaces.}
If $\nu_1<as$, G-26 gives zero padded multiplicity at every
positive stretch. R-32 still applies to the degree-$n$
determinant cone, because $r<2n-2$ and $\nu$ is nonrectangular.
It gives determinant occurrence at every
$j\ge\beta_D^2$, where
\[
 \beta_D=\left\lceil\frac{3r^2b_D}{8}\right\rceil
             \bigl(ns+(r-1)|\eta|\bigr)^{2(r^2-1)}.
\]
Here $b_D$ and $\eta$ have their displayed determinant
definitions; no lower tail degree is used in this case.
If $\nu_1>as$, G-25 supplies the comparison, with its
nonzero integer exponent $\Delta$. Combining these cases
with G-26 proves the stated eventual classification.
\end{proof}

\paragraph{Scope.}
The new semigroup argument allows a period on the algebra
used only for an upper bound; it does not remove that period.
Rectangular directions with more than $m$ rows, equality
faces at $2m$ rows, and strict directions beyond $2m$ rows
retain their separate stable-point gates. Earlier individual
degrees and varying-direction complexity questions are also
unsettled. The construction is independent of the stipulated
analytic inputs and makes no efficient-size claim for the
displayed thresholds.


\begin{theorem}[G-28: Boundary comparison from a generic stabilizer bound]
\label{G-28}
Let $n>m\ge17$ and
\[
                  m<r\le\min\{2n-2,m^2\}.
\]
Let $\nu\vdash ns$, $s\ge1$, have exactly $r$ positive
parts and be nonrectangular. Put $a=n-m$. If $\nu_1\le as$,
the determinant multiplicity is eventually strictly larger
than the padded permanent multiplicity at every integer
stretch. On the equality face, the explicit constants and
threshold of G-27 apply without change. Below the face,
the padded multiplicity is always zero and the determinant
occurs for every $j\ge\beta_D^2$, with $\beta_D$ as in G-27.

Consequently, for $n>m\ge17$ and $4\le r\le2m$, every
fixed nonrectangular full-row direction is classified
eventually: the padded permanent is larger exactly when
\[
 \nu_1>(n-m)s,\qquad
 rm^2-(r-2)n^2-2m+r-1>0,
\]
and the determinant is eventually larger otherwise.
The larger boundary range in the first paragraph does not
extend the strict-direction theorem past $2m$ rows.
\end{theorem}

\begin{proof}
We retain the Pieri tail reduction and replace the lower
stable quotient by the generic-orbit dimension upper bound
proved in R-33.

First suppose $\nu_1=as$. Define $q=r-1$, the tail $\mu$,
the weights $\eta,\theta$, and every constant $u,v,b_D,b_B,
b,L,p,\beta,T,E$ exactly as in G-27. In particular
\[
 \begin{aligned}
 u&=(r-2)n^2-r^2+3+r_\eta,\\
 v&=q(m^2-q)-2m+3+r_\theta.
 \end{aligned}
\]
Let $A$ and $B$ be the same determinant and lower permanent
enhanced invariant algebras. The Pieri bound still gives
\[
 m_{j\nu}(\ell^a\operatorname{per}_m)\le\dim B_j,
 \qquad \dim A_j=m_{j\nu}(\det_n).
\]

R-32 applies to $A$ because $r\le2n-2$, and the
nonrectangular label supplies a no-scaling stable point.
Thus $\dim A=u$ and every degree at least $\beta^2$ occurs
in $A$. For $B$, one has $4\le q<m^2$. R-33 gives only
\[
                         w:=\dim B\le v.
\]
No stable point, equality $w=v$, or period-one property is
assumed for this lower ring.

The same dimension difference computation as in G-27 uses
only $n\ge m+1$ and $r\ge m+1$, so
\[
 u-v\ge(2m-1)r-m^2-2m-1\ge m^2-m-2>0.
\]
Since $1\le q\le m^2-1$, the elementary identity
\[
 q(m^2-q)-(m^2-1)=(q-1)(m^2-q-1)\ge0
\]
also gives $v\ge m^2-2m+2>2$. Thus $u\ge v+1$,
regardless of whether the lower quotient attains its bound.

The same ambient representations and parameter maps give
the common generator bound $\beta$, the standard Veronese
degree $E$, and parameter degree at most $p$. If $w=0$,
the nonnegatively graded domain $B$ with $B_0=\mathbb C$
is just $\mathbb C$, so every positive padded multiplicity
is zero and the conclusion follows from occurrence in $A$.
Suppose henceforth $w\ge1$.

The arbitrary-period argument of G-27 applies to $B$ with
its actual nonzero generator degrees. For $j=Ek+t$,
$0\le t<E$, it gives
\[
 \dim A_j\ge\dim A^{(E)}_{k-1},\qquad
 \dim B_j\le\dim B^{(E)}_{k+2}.
\]
The projective image degree bound uses its actual dimension:
\[
 \deg\operatorname{Proj}B^{(E)}\le(pE)^{w-1}.
\]
For $w=1$ this is the degree-one reduced point case. The
standard-domain Hilbert upper bound, followed by $w\le v$,
gives
\[
 \begin{aligned}
 \dim B_j
   &\le(pE)^{w-1}\binom{k+w+1}{w-1}\\
   &\le(pE)^{v-1}\binom{k+v+1}{v-1}.
 \end{aligned}
\]
The lower Noether-normalization bound remains
$\dim A_j\ge\binom{k+v-1}{v}$. Therefore exactly the
binomial comparison in G-27 proves strict dominance for
\[
                    j\ge E\bigl(4v(pE)^{v-1}+1\bigr).
\]
No equality in a generic quotient-dimension formula has
entered this proof.

If $\nu_1<as$, G-26 makes the padded multiplicity zero.
The same R-32 application to $A$ yields determinant
occurrence at every $j\ge\beta_D^2$. This proves the
first assertion. Finally, $r\le2m$ implies both
$r\le2n-2$ and $r\le m^2$. G-28 therefore supplies
all nonrectangular on/below-face cases in that range;
G-26 handles the smaller rows and G-25 handles every
strict direction. Their nonzero integer exponent gives
the asserted complete nonrectangular classification.
\end{proof}

\paragraph{Scope.}
The additional input is a generic finite projective
stabilizer on the lower permanent image, proved from its
exact tangent kernel and Pl\"ucker coordinates in R-33.
It is used only for a quotient-dimension upper bound.
Rectangular directions beyond $m$ rows, strict directions
past $2m$ rows, individual earlier degrees and changing
directions retain their stated gates. This result does
not identify generic freeness with stability and does
not claim an efficient positive formula from the displayed
large thresholds.
